% ============================================================================
% Rozdział: Portcullis -- bramka na pull request (komponent C5).
% Źródło: migawka letsgolegacy.portcullis @ 6a018a2 (2026-09-26).
% ============================================================================
\section{Portcullis --- bramka na pull request}\label{sec:portcullis}

Portcullis to komponent C5 frameworku letsgolegacy: bramka uruchamiana na każdym pull
requeście (PR). Wykonuje analyzery Roslyn na kodzie C\#, raportuje findingi (pojedyncze
naruszenia reguł zgłoszone przez skan), ale blokuje tylko te, które leżą na liniach
zmienionych przez PR, i~zapisuje je jako SARIF 2.1.0 z~plikiem baseline znanych
findingów. Repozytorium \repopath{letsgolegacy.portcullis} jest osobnym narzędziem
z~własnym cyklem wydań; rdzeń Second Key czyta wyłącznie jego wynik.

Rozdział opisuje migawkę z~commitu \repopath{6a018a2} (2026-09-26). W~tym stanie
\textbf{nic nie zostało wydane}: nie istnieje żaden tag \repopath{v*}, w~nuget.org nie ma
żadnego pakietu, a~w~GHCR żadnego obrazu. Wersja \repopath{0.1.0}
w~\repopath{Directory.Build.props} to numer, który poniesie pierwsze wydanie, a~nie wersja,
którą można zainstalować. Działa budowanie ze źródeł (sekcja~\ref{sec:portcullis-tutorial},
krok~7).

\subsection{Po co istnieje i~miejsce w~łańcuchu}\label{sec:portcullis-cel}

\paragraph{Problem.} Agenci AI piszą dużą część commitowanego kodu. Ich charakterystyczny
błąd to według README nie usterka, tylko zmiana, która się kompiluje, przechodzi review
i~po cichu łamie granicę architektoniczną ustaloną przez zespół miesiące wcześniej.
Portcullis zamienia takie decyzje w~analyzery Roslyn, więc build zatrzymuje się na
\repopath{plik:linia}, zamiast dryf wyjść na jaw po roku.

README odróżnia Portcullis od narzędzi, które odpowiadają na inne pytania: SonarQube
ocenia jakość kodu, Snyk szuka podatności, CodeClimate mierzy złożoność. NetArchTest
i~ArchUnitNET sprawdzają zadeklarowaną architekturę, ale jako biblioteki testowe oparte na
refleksji, które zespół podłącza sam. Portcullis przyjmuje odwrotne założenia: tylko .NET,
analiza na modelu semantycznym Roslyn (rozwiązuje prawdziwe symbole i~wskazuje linię
zamiast dopasowywać tekst), dostarczanie jako diagnostyki kompilatora w~czasie builda plus
bramka merge.

\paragraph{Dwie właściwości, które odróżniają go od lintera.}
\begin{enumerate}[leftmargin=*]
  \item \textbf{Bramka zawężona do diffu PR.} Silnik reguł uruchomiony na prawdziwym
        repozytorium znajduje naruszenia sprzed PR. Bramka, która blokuje na wszystkich,
        blokuje każdy PR, także jednolinijkowy i~niezwiązany, i~według README zostaje
        wyłączona w~ciągu tygodnia. Portcullis raportuje wszystko, co znalazł, ale blokuje
        tylko na naruszeniach w~liniach zmienionych przez PR
        (sekcja~\ref{sec:portcullis-diff-gate}).
  \item \textbf{Pochodzenie AI jako wejście, nie produkt.} Linie przypisane agentowi AI
        -- ustalane z~prawdziwej historii git (trailery commitów i~autorstwo), nigdy ze
        stylu kodu -- dostają ostrzejszy próg poziomu niż to samo naruszenie napisane
        ręcznie. Sygnał nigdy nie jest odznaką ani wynikiem w~komentarzu; zmienia wyłącznie
        to, który próg obowiązuje.
\end{enumerate}

\paragraph{Miejsce w~łańcuchu Second Key.} W~łańcuchu opisanym w~README rdzenia
(sekcja~\ref{sec:architektura}) Portcullis jest krokiem 4 („Gate”): komponent C5, analyzery
Roslyn na zmienionych liniach, wynik \repopath{*.sarif}. Kontrakt między komponentami to
pliki:

\begin{itemize}[leftmargin=*]
  \item \textbf{wejście z~C4} -- pakiet standardów (\repopath{standards/}, repozytorium
        \path{letsgolegacy.secondkey-standards}) jest konsumowany przez agenta
        migracyjnego i~przez bramkę C5; pakiet cytuje dosłownie cztery identyfikatory
        reguł migracyjnych \repopath{PORTCULLIS\_MIG\_*} (sekcja~\ref{sec:standardy}),
        dlatego \repopath{RuleRegistryTests} je przypina;
  \item \textbf{wyjście do C10} -- plik SARIF 2.1 wchodzi do evidence pack (pakietu
        dowodowego) składanego przez C10; polecenie
        \repopath{sk gate} rdzenia streszcza SARIF per reguła i~kończy się kodem 1, gdy
        finding blokuje (sekcja~\ref{sec:cli}).
\end{itemize}

Bramka nie jest przenoszona do repozytorium rdzenia: ma własny cykl wydań, a~rdzeń czyta
tylko jej wynik. Z~Portcullis rdzeń przejął wzorzec bramki zawężonej do diffu (sekcja
bramki w~pakiecie dowodowym oddziela findingi blokujące od raportowanych), SARIF 2.1 jako
format wymiany, identyfikatory reguł pokazywane w~pakiecie dowodowym i~testy mutacyjne jako
nawyk.

\finding{Dla reszty łańcucha Portcullis to jeden artefakt: plik SARIF z~findingami, które
wprowadził dany pull request.}

Reguły zasad wdrażają zasady P2, P4, P9, P10, P11 i~P15 konstytucji architektury
\path{architecture-standards}. Użytkownik nie musi przyjmować tego standardu, ale tam
zapisano uzasadnienie każdej reguły. Cztery reguły migracyjne nie wynikają z~żadnej zasady
P1--P15; ich uzasadnienie jest w~\repopath{docs/rules/MIGRATION.md}
(sekcja~\ref{sec:portcullis-migracja}).

\zrodla{\path{letsgolegacy.portcullis/README.md}, \path{letsgolegacy.portcullis/docs/WORKPLAN.md},
\path{letsgolegacy.portcullis/CHANGELOG.md}, \path{letsgolegacy.portcullis/docs/rules/MIGRATION.md},
\path{letsgolegacy.secondkey/README.md}, \path{letsgolegacy.secondkey/docs/WORKPLAN.md},
\path{letsgolegacy.secondkey/docs/reference-notes.md}, \path{letsgolegacy.secondkey/docs/cli.md}}

\subsection{Budowa repozytorium i~przebieg skanu}\label{sec:portcullis-budowa}

\begin{xltabular}{\linewidth}{@{}>{\raggedright\arraybackslash}p{4.9cm} l >{\raggedright\arraybackslash}X@{}}
\toprule
\textbf{Projekt} & \textbf{Cel} & \textbf{Rola} \\
\midrule
\endfirsthead
\toprule
\textbf{Projekt} & \textbf{Cel} & \textbf{Rola} \\
\midrule
\endhead
\path{src/Portcullis.Rules} & \texttt{netstandard2.0} & Analyzery Roslyn; pakowane jako \texttt{Portcullis.Analyzers}. \\
\path{src/Portcullis.Engine} & \texttt{net10.0} & Harness skanu, model wyniku, pochodzenie z~git, baseline, SARIF. Nie jest publikowany jako pakiet. \\
\path{src/Portcullis.Cli} & \texttt{net10.0} & Narzędzie \texttt{portcullis} (pakiet \texttt{Portcullis.Cli}). \\
\path{src/Portcullis.CiComment} & \texttt{net10.0} & Narzędzie \texttt{portcullis-ci-comment} (pakiet \texttt{Portcullis.CiComment}). \\
\path{tests/Portcullis.Engine.Tests} & \texttt{net10.0} & Reguły i~ich mutanty, skaner, bramka, pochodzenie, baseline, SARIF, konfiguracja. \\
\path{tests/Portcullis.CiComment.Tests} & \texttt{net10.0} & Formatowanie komentarza, diff naruszeń, sticky comment, opcje. \\
\path{tests/Portcullis.Cli.Tests} & \texttt{net10.0} & Opcje CLI i~testy procesu end-to-end (SARIF, baseline). \\
\path{fixtures/analyzer-consumer} & \texttt{net8.0} & Konsument pakietu analyzerów na najstarszym wspieranym SDK (sekcja~\ref{sec:portcullis-roslyn}). \\
\path{fixtures/mock-engine-output} & -- & Ręcznie napisane wyniki skanu zgodne z~SPEC §3, używane w~testach komentarza. \\
\path{fixtures/consumer-violations.json} & -- & Lista „musi złapać” z~milestone'u M1D (sekcja~\ref{sec:portcullis-bootstrap}). \\
\bottomrule
\end{xltabular}

\paragraph{Przebieg skanu.} Rdzeniem jest metoda \repopath{Scanner.ScanAsync}; przyjmuje
ścieżkę skanu, opcjonalny raport pochodzenia, katalog, względem którego raport podaje
ścieżki, i~opcjonalny baseline:

\begin{enumerate}[leftmargin=*]
  \item Wylicza rekurencyjnie pliki \repopath{*.cs} pod ścieżką, pomija każdy plik, którego
        ścieżka względna zawiera segment \repopath{bin}, \repopath{obj} lub \repopath{.git}
        (bez względu na wielkość liter), i~sortuje listę porządkiem ordinalnym.
  \item Parsuje każdy plik przez \repopath{CSharpSyntaxTree.ParseText} i~buduje z~nich
        jedną kompilację \repopath{PortcullisScanTarget} z~referencjami
        \repopath{ScanReferences.All}.
  \item Wczytuje \repopath{portcullis.json} z~katalogu skanu i~zamienia go na
        \repopath{AnalyzerOptions}; katalog skanu trafia do klucza
        \repopath{portcullis\_path\_root} (sekcja~\ref{sec:portcullis-konfiguracja}).
  \item Uruchamia wszystkie analyzery \repopath{RuleRegistry.All} (11 analyzerów,
        15 identyfikatorów diagnostyk) przez \repopath{CompilationWithAnalyzers}.
  \item Mapuje diagnostyki na rekordy naruszeń, sortuje je po pliku, linii, kolumnie,
        identyfikatorze reguły i~komunikacie, nadaje fingerprinty
        (sekcja~\ref{sec:portcullis-baseline}) i~oznacza findingi akceptowane przez
        baseline.
  \item Liczy werdykt bramki: przy zdegradowanym raporcie pochodzenia -- cały skan
        z~podanym powodem; przy kompletnym -- eskalacja linii AI i~bramka zawężona do
        diffu; bez raportu -- cały skan (sekcja~\ref{sec:portcullis-diff-gate}).
\end{enumerate}

CLI wykonuje dookoła skanu resztę pracy: czyta baseline przed skanem, liczy pochodzenie
z~git na granicy CLI (silnik pozostaje niezależny od git), po skanie zapisuje baseline
(\repopath{--write-baseline}) i~SARIF (\repopath{--sarif}), wypisuje ostrzeżenia na stderr,
dokument JSON na stdout i~kończy się kodem bramki.

\paragraph{Co widzi kompilacja skanu.} Skan parsuje źródła bezpośrednio, bez MSBuild,
więc nie ma referencji projektu. \repopath{ScanReferences} dodaje dwie rzeczy:

\begin{itemize}[leftmargin=*]
  \item \textbf{zestawy frameworku działającego runtime'u} -- zaufane zestawy platformy
        z~katalogu, w~którym leży \repopath{System.Private.CoreLib}; zestawy aplikacji
        (Roslyn, sam Portcullis) są pominięte, żeby skanowane repozytorium nie wiązało się
        z~zależnościami Portcullis. Referencje powstają raz na proces (komentarz w~kodzie:
        około 170 zestawów);
  \item \textbf{deklaracje API .NET Framework} (\repopath{LegacyApiSurface}) -- same
        sygnatury, kompilowane w~pamięci do zestawu \repopath{Portcullis.LegacyApiSurface}:
        rdzeń System.Web (\repopath{HttpContext} ze statycznym \repopath{Current}, żądanie,
        odpowiedź, sesja, cache, aplikacja, interfejsy modułu i~handlera,
        \repopath{HttpRuntime}, \repopath{VirtualPathUtility}), przestrzenie
        \repopath{System.Web.SessionState}, \repopath{Caching}, \repopath{Hosting},
        \repopath{Security}, \repopath{Configuration} oraz rodzina
        \path{System.Configuration.ConfigurationManager}. Błąd kompilacji tej
        powierzchni rzuca wyjątek, zamiast po cichu wyłączyć reguły migracyjne.
\end{itemize}

Nie są deklarowane \repopath{System.Web.Mvc}, \repopath{System.Web.UI}
i~\repopath{System.Web.Http}. Skan nie widzi pakietów NuGet ani innych projektów, ani
globalnych dyrektyw \repopath{using} generowanych przez SDK dla
\path{<ImplicitUsings>enable</ImplicitUsings>} -- trafiają one do \repopath{obj/}, który
skan pomija. Pozostałe reguły to nie zmienia: działają tylko na typach zadeklarowanych
w~skanowanym źródle, a~typ frameworku nie ma lokalizacji w~źródle.

\zrodla{\path{letsgolegacy.portcullis/CONTRIBUTING.md}, \path{letsgolegacy.portcullis/portcullis.sln},
\path{letsgolegacy.portcullis/src/Portcullis.Engine/Scanner.cs},
\path{letsgolegacy.portcullis/src/Portcullis.Engine/Semantics/ScanReferences.cs},
\path{letsgolegacy.portcullis/src/Portcullis.Engine/Semantics/LegacyApiSurface.cs},
\path{letsgolegacy.portcullis/src/Portcullis.Cli/Program.cs},
\path{letsgolegacy.portcullis/src/Portcullis.Engine/Provenance/ProvenanceModels.cs},
\path{letsgolegacy.portcullis/docs/rules/MIGRATION.md}}

\subsection{Specyfikacja (SPEC)}\label{sec:portcullis-spec}

\repopath{docs/SPEC.md} to dokument milestone'u M0, napisany przed pierwszą regułą zgodnie
z~zasadą „najpierw specyfikacja, potem kod”. Obejmuje dokładnie pięć rzeczy: format
kontraktu reguły, schemat rekordu naruszenia, schemat JSON wyniku silnika, interfejs
sygnału pochodzenia i~werdykt dla każdej zasady P1--P15. Wszystko, co powstało później,
traktuje go jako zamrożony kontrakt: tor B (\repopath{track-b-ci}) budował i~testował akcję
GitHub na ręcznie napisanym pliku zgodnym z~§3, bez czekania na silnik toru A. Późniejsze
zmiany dopisano jako addenda (ticket R2 i~R3), a~pola dodano addytywnie.

\subsubsection{Format kontraktu reguły}

\repopath{architecture-standards} eksportuje reguły jako pluginy w~formacie Agent Plugins
1.0.0. Żadna z~zasad P1--P15 nie jest tam osobno czytelna maszynowo: to nagłówki H3
w~jednym dokumencie Markdown, cytowane gdzie indziej tylko jako goły token (\repopath{P7})
w~tablicy \repopath{skills[].principles} katalogu. SPEC mapuje więc pola reguły uczciwie,
rozróżniając odziedziczone od nowych:

\begin{itemize}[leftmargin=*]
  \item \textbf{nowe}: \repopath{ruleId} (w~kodzie \repopath{DiagnosticDescriptor.Id}),
        \repopath{defaultSeverity} (\repopath{error}, \repopath{warning} lub
        \repopath{info}; Agent Plugins nie zna poziomów) i~\repopath{roslynAnalyzer}
        (klasa \repopath{DiagnosticAnalyzer} zarejestrowana w~\repopath{RuleRegistry.All});
  \item \textbf{odziedziczone}: \repopath{principleIds} (goły token, np. \repopath{P2} --
        jedyny klucz łączący z~konstytucją; w~kodzie niesie go kategoria deskryptora),
        \repopath{sourcePlugin} (plugin, z~którego pochodzi tekst zasady, np.
        \repopath{architecture-core}) i~\repopath{sourceDoc} (dokument zasady);
  \item \textbf{luźno wzorowane}: \repopath{title} (na \repopath{description} w~\repopath{plugin.json})
        i~\repopath{description} (na \repopath{description} w~\repopath{SKILL.md}: „kiedy reguła
        strzela”).
\end{itemize}

\paragraph{Identyfikatory z~podkreśleniami.} Konwencji z~łącznikami nie da się zgłosić
przez potok Roslyn, którego używa skaner. \repopath{Diagnostic.Create} przyjmuje taki
identyfikator, ale \repopath{ReportDiagnostic} wewnątrz analyzera odrzuca go wyjątkiem
\repopath{ArgumentException} („which is not a valid identifier”), który Roslyn zgłasza jako
\repopath{AD0001}. Kropka jest odrzucana tak samo; podkreślenia i~same litery przechodzą.
Ponieważ SPEC jest zamrożony, problem zapisano w~\repopath{docs/MUTATIONS.md} §1, a~reguły
używają formy \repopath{PORTCULLIS\_P<n>\_<SLUG>}. \repopath{RuleRegistryTests} sprawdza, że
każdy identyfikator zaczyna się od \repopath{PORTCULLIS\_}, zawiera tylko litery, cyfry
i~podkreślenia, a~kategoria to \repopath{P1}--\repopath{P15}, \repopath{Meta} albo
\repopath{Migration}. Addendum R2: kategorie \repopath{Meta} (para reguł pokrycia
konwencji) i~\repopath{Migration} nie mają identyfikatora zasady; reguły migracyjne mają
postać \repopath{PORTCULLIS\_MIG\_<SLUG>}.

\subsubsection{Rekord naruszenia}

Rekord \repopath{Portcullis.Engine.Model.Violation} ma pięć pól ustalonych na starcie
i~dwa dodane w~tickecie R3:

\begin{xltabular}{\linewidth}{@{}>{\raggedright\arraybackslash}p{2.4cm} >{\raggedright\arraybackslash}p{2.6cm} >{\raggedright\arraybackslash}X@{}}
\toprule
\textbf{Pole JSON} & \textbf{Typ} & \textbf{Znaczenie} \\
\midrule
\endfirsthead
\toprule
\textbf{Pole JSON} & \textbf{Typ} & \textbf{Znaczenie} \\
\midrule
\endhead
\texttt{ruleId} & string & Identyfikator reguły, która strzeliła. \\
\texttt{filePath} & string & Ścieżka w~stylu POSIX, względna wobec skanowanego katalogu; pusta dla findingu dotyczącego całego repozytorium (wtedy linia to 1). \\
\texttt{line} & integer & Numer linii liczony od 1. \\
\texttt{message} & string & Komunikat konkretnego naruszenia. \\
\texttt{severity} & \texttt{error} | \texttt{warning} | \texttt{info} & Mapowane 1:1 z~\texttt{DiagnosticSeverity}; \texttt{Hidden} przechodzi w~\texttt{info}. \\
\texttt{fingerprint} & string lub \texttt{null} & Tożsamość findingu między commitami (SHA-256; sekcja~\ref{sec:portcullis-baseline}); \texttt{null} tylko w~JSON sprzed R3. \\
\texttt{baselined} & boolean & \texttt{true}, gdy baseline przekazany do skanu akceptuje finding: jest raportowany, ale nie blokuje. \\
\bottomrule
\end{xltabular}

\begin{lstlisting}
{ "ruleId": "PORTCULLIS_P9_ORPHAN_ENTITY", "filePath": "Domain/QrCode.cs",
  "line": 11, "message": "...", "severity": "warning" }
\end{lstlisting}

\subsubsection{Wynik skanu}

\repopath{portcullis scan <ścieżka>} drukuje na stdout jeden dokument JSON: rekord
\path{Portcullis.Engine.Model.ScanResult}, nazwy pól w~camelCase, wartości \repopath{null}
są wypisywane.

\begin{xltabular}{\linewidth}{@{}>{\raggedright\arraybackslash}p{3.6cm} >{\raggedright\arraybackslash}X@{}}
\toprule
\textbf{Pole} & \textbf{Znaczenie} \\
\midrule
\endfirsthead
\toprule
\textbf{Pole} & \textbf{Znaczenie} \\
\midrule
\endhead
\texttt{schemaVersion} & Wersja schematu, dziś \texttt{1.0.0}; podbijana przy zmianie łamiącej. \\
\texttt{engineVersion} & Wersja informacyjna zestawu silnika, z~commitem źródła (np. \texttt{1.0.0+<sha>}). \\
\texttt{scannedPath} & Bezwzględna ścieżka skanu. \\
\texttt{scanStartedUtc} & Według SPEC: czas rozpoczęcia skanu (ISO 8601, UTC). \\
\texttt{scanDurationMs} & Czas trwania skanu w~milisekundach. \\
\texttt{filesScanned} & Liczba sparsowanych plików \texttt{*.cs} (bez \texttt{bin/}, \texttt{obj/}, \texttt{.git/}). \\
\texttt{rulesEvaluated} & Liczba analyzerów w~\texttt{RuleRegistry.All}, które się wykonały (w~migawce 11). \\
\texttt{violations} & Tablica rekordów naruszeń, sortowana po \texttt{filePath}, \texttt{line}, kolumnie, \texttt{ruleId}, \texttt{message} (addendum R3; wcześniej tylko plik i~linia). \\
\texttt{summary} & \texttt{errorCount}, \texttt{warningCount}, \texttt{infoCount} -- liczby naruszeń według poziomu. \\
\texttt{gate} & Dodane w~M5A: \texttt{blocked}, \texttt{scope} (\texttt{all} lub \texttt{diff}), \texttt{blockingErrorCount}, \texttt{degradedReason} (\texttt{null} poza fallbackiem), \texttt{acceptedByBaselineCount} (R3). \texttt{null} tylko w~JSON sprzed M5A. \\
\texttt{configuration} & \texttt{path} wczytanego \texttt{portcullis.json} (\texttt{null}, gdy go nie ma) i~\texttt{error}, gdy plik istnieje, ale nie dał się sparsować. \\
\texttt{baseline} & Dodane w~R3: \texttt{path}, \texttt{entryCount}, \texttt{acceptedCount}; \texttt{null} bez \texttt{--baseline}. \\
\bottomrule
\end{xltabular}

Pola \repopath{gate}, \repopath{configuration}, \repopath{baseline}, \repopath{fingerprint}
i~\repopath{baselined} dodano addytywnie, bez podbijania \repopath{schemaVersion}: są
opcjonalne i~mają wartości domyślne, więc każda wcześniejsza konstrukcja, każdy fixture
JSON sprzed zmiany i~każdy konsument, który ich nie zna, działa bez zmian.

SPEC zawiera przykład wyniku uruchomiony naprawdę 2026-08-17 na 156 plikach repozytorium
referencyjnego, zaraz po szkielecie M2: zero reguł, zero naruszeń, kod wyjścia 0
(sekcja~\ref{sec:portcullis-bootstrap}). Przykład z~dzisiejszymi polami jest
w~sekcji~\ref{sec:portcullis-tutorial}, krok~5.

\paragraph{Kody wyjścia CLI.} \repopath{0} -- bramka nie blokuje, \repopath{1} -- bramka
blokuje, \repopath{2} -- błędne użycie. Tylko \repopath{2} oznacza awarię samego
Portcullis; \repopath{1} to werdykt. Kod 2 zwracają też: ścieżka skanu, która nie jest
istniejącym katalogiem (skan, który się nie odbył, nie może dać zielonej bramki), baseline,
którego nie da się użyć, oraz pliki SARIF lub baseline, których nie da się zapisać.

\subsubsection{Interfejs sygnału pochodzenia}

Typy interfejsu są zdefiniowane w~pliku
\path{src/Portcullis.Engine/Provenance/ProvenanceModels.cs}. Raport pochodzenia ma pola
\repopath{schemaVersion}, \repopath{commit} (commit lub zakres, dla którego go policzono)
i~listę zakresów:

\begin{xltabular}{\linewidth}{@{}>{\raggedright\arraybackslash}p{4.4cm} >{\raggedright\arraybackslash}X@{}}
\toprule
\textbf{Pole} & \textbf{Znaczenie} \\
\midrule
\endfirsthead
\toprule
\textbf{Pole} & \textbf{Znaczenie} \\
\midrule
\endhead
\texttt{ranges[].filePath} & Ścieżka POSIX względna wobec katalogu głównego repozytorium git. \\
\texttt{ranges[].startLine}, \texttt{.endLine} & Zakres linii liczonych od 1, obustronnie domknięty -- nigdy cały plik ani cała sesja. \\
\texttt{ranges[].source} & \texttt{human}, \texttt{ai}, \texttt{mixed} lub \texttt{unknown}. \\
\texttt{ranges[].confidence} & Liczba 0,0--1,0. \\
\path{ranges[].attribution.tool} & Identyfikator narzędzia (w~kodzie m.in. \texttt{github-copilot}, \texttt{cursor}) albo \texttt{null}, gdy nie da się go ustalić. \\
\path{ranges[].attribution.method} & Sposób ustalenia: \texttt{commit-trailer}, \texttt{git-blame-authorship}, \texttt{no-attribution-signal}. \\
\path{ranges[].attribution.commitSha} & Commit albo \texttt{null}. \\
\bottomrule
\end{xltabular}

Kontrakt konsumpcji: naruszenie w~zakresie o~źródle \repopath{ai} dostaje ostrzejszy próg
poziomu. Pochodzenie nigdy nie jest osobnym raportem ani poleceniem CLI -- „tę linię napisało
AI” bez żadnego skutku to według SPEC ciekawostka, nie produkt. Kształt jest odpowiedzią na
nieudany sygnał „edit survival” narzędzia \repopath{copilot-scope}: średnią na poziomie sesji,
bez pliku, linii i~hunku, emitowaną przez jeden edytor. Dlatego każdy zakres musi mieć plik
i~linie, a~pola \repopath{attribution} jawnie mówią, skąd sygnał pochodzi.

Później do kontraktu doszło pole \repopath{Status} (\repopath{Complete} albo
\repopath{Degraded}) i~\repopath{DegradedReason} -- zob. sekcja~\ref{sec:portcullis-diff-gate}.

\subsubsection{Werdykt dla zasad P1--P15}

„Expressible” oznacza zasadę, którą pierwsza reguła mogła wdrożyć od razu; „Deferred” --
częściowo mechaniczną, ale wymagającą osądu międzyplikowego, czasowego lub behawioralnego;
„Out of scope” -- artefakt, który nie jest kodem C\#, więc żaden analyzer Roslyn go nie
sięgnie.

\begin{xltabular}{\linewidth}{@{}l >{\raggedright\arraybackslash}p{4.4cm} >{\raggedright\arraybackslash}p{2.0cm} >{\raggedright\arraybackslash}X@{}}
\toprule
\textbf{ID} & \textbf{Zasada} & \textbf{Werdykt} & \textbf{Uzasadnienie} \\
\midrule
\endfirsthead
\toprule
\textbf{ID} & \textbf{Zasada} & \textbf{Werdykt} & \textbf{Uzasadnienie} \\
\midrule
\endhead
P1 & AppHost jest composition rootem & Deferred & Kształt deklaracji zasobów jest sprawdzalny; reszta to osąd całego repozytorium. \\
P2 & Wspólne jądro, nie wspólna domena & Expressible & Sufit ok. 800 linii i~kierunek zależności jądra. \\
P3 & Usługa i~baza per bounded context & Deferred & Graf zależności; współdzielenie logiczne czy fizyczne -- osąd. \\
P4 & Persystencja przenośna i~migrowana & Expressible & \texttt{EnsureCreated()}, \texttt{MigrateAsync}, \texttt{HasData}. \\
P5 & Konfiguracja przez środowisko, sekrety przez platformę & Deferred & Literały sekretów -- tak; jeden właściciel klucza -- niezmiennik między repozytoriami. \\
P6 & Jeden kontener na usługę, wieloetapowy Dockerfile & Out of scope & Dockerfile, nie C\#. \\
P7 & Topologia Fly.io & Out of scope & \texttt{fly.toml}, nie C\#. \\
P8 & Opcjonalne zależności degradują się & Deferred & Warunkowa rejestracja DI -- tak; czy fallback działa -- nie. \\
P9 & \texttt{Program.cs} jest manifestem; ścisłe warstwy & Expressible & Kontroler bez bezpośredniego \texttt{DbContext}. \\
P10 & Rozszerzalność przez interfejs i~rejestrację & Expressible & Wykrywanie własnych klas bazowych. \\
P11 & Antykorupcja na krawędzi & Expressible & Typy SDK dostawcy tylko w~adapterze. \\
P12 & CI/CD sterowane tagami & Out of scope & YAML CI, nie C\#. \\
P13 & Testy na warstwie, która trzyma logikę & Deferred & Istnienie testów -- tak; jakość asercji i~kolejność w~czasie -- nie. \\
P14 & Dokumentacja zapisuje uzasadnienie & Out of scope & Proza. \\
P15 & Obserwowalność decydowana w~czasie builda & Expressible & Obecność \texttt{AddServiceDefaults()}. \\
\bottomrule
\end{xltabular}

Bilans: 6~Expressible (P2, P4, P9, P10, P11, P15), 5~Deferred (P1, P3, P5, P8, P13),
3~Out of scope jako inny typ artefaktu (P6, P7, P12) i~1~Out of scope jako proza (P14).
P1 i~P5 są rozdzielone (częściowo Expressible) i~tabela SPEC mówi to wprost, zamiast
wymuszać jeden werdykt.

\subsubsection{Co jest findingiem, a~co blokuje}

\begin{itemize}[leftmargin=*]
  \item \textbf{Finding} to każda diagnostyka zgłoszona przez analyzer z~\repopath{RuleRegistry.All}
        na kompilacji skanu, zmapowana na rekord naruszenia. Każdy finding jest
        raportowany w~\repopath{violations} i~w~\repopath{summary}, niezależnie od zakresu
        bramki i~baseline.
  \item \textbf{Blokuje} wyłącznie finding, który (a)~ma poziom \repopath{error} po
        ewentualnej eskalacji pochodzenia, (b)~mieści się w~zakresie bramki -- cały skan
        (\repopath{scope} \repopath{all}) albo zmienione linie zakresu commitów
        (\repopath{scope} \repopath{diff}) -- i~(c)~nie jest akceptowany przez baseline.
        Liczba takich findingów to \repopath{gate.blockingErrorCount}; bramka blokuje, gdy
        jest większa od zera.
  \item \repopath{warning} i~\repopath{info} same nie blokują. Ostrzeżenie może zablokować
        tylko po eskalacji do \repopath{error} na linii przypisanej AI. Podniesienie poziomu
        w~\repopath{.editorconfig} działa tylko w~kanale pakietu analyzerów (w~buildzie), nie
        w~bramce CLI (sekcja~\ref{sec:portcullis-konfiguracja}).
  \item SARIF zawiera findingi wszystkich poziomów, filtrowane tym samym zakresem i~tym
        samym baseline, którego używa bramka (sekcja~\ref{sec:portcullis-sarif}).
\end{itemize}

\zrodla{\path{letsgolegacy.portcullis/docs/SPEC.md}, \path{letsgolegacy.portcullis/docs/MUTATIONS.md},
\path{letsgolegacy.portcullis/docs/TUTORIAL.md},
\path{letsgolegacy.portcullis/src/Portcullis.Engine/Model/ScanResult.cs},
\path{letsgolegacy.portcullis/src/Portcullis.Engine/Model/Violation.cs},
\path{letsgolegacy.portcullis/src/Portcullis.Engine/Scanner.cs},
\path{letsgolegacy.portcullis/src/Portcullis.Engine/Provenance/ProvenanceModels.cs},
\path{letsgolegacy.portcullis/src/Portcullis.Rules/RuleRegistry.cs},
\path{letsgolegacy.portcullis/src/Portcullis.Cli/Program.cs},
\path{letsgolegacy.portcullis/tests/Portcullis.Engine.Tests/Rules/RuleRegistryTests.cs}}

\subsection{Reguły zasad architektury i~reguły pokrycia konwencji}\label{sec:portcullis-reguly}

\repopath{RuleRegistry.All} rejestruje 11 analyzerów: sześć reguł zasad (P2, P9, P10 --
tor A; P4, P11, P15 -- przebieg uzupełniający), cztery reguły migracyjne i~jeden analyzer
pokrycia konwencji. Wszystkie wyłączają analizę kodu generowanego
(\repopath{GeneratedCodeAnalysisFlags.None}) i~działają współbieżnie. Reguły zasad
rozpoznają role (jądro, domena, adapter, punkt wejścia) po nazwach folderów i~plików, bo
w~zwykłym C\# nic tych ról nie deklaruje; nazwy są konfigurowalne
(sekcja~\ref{sec:portcullis-konfiguracja}).

\begingroup\small
\begin{xltabular}{\linewidth}{@{}>{\raggedright\arraybackslash}p{4.3cm} >{\raggedright\arraybackslash}p{1.25cm} >{\raggedright\arraybackslash}X >{\raggedright\arraybackslash}X@{}}
\toprule
\textbf{Identyfikator} & \textbf{Poziom} & \textbf{Co wykrywa} & \textbf{Mechanizm} \\
\midrule
\endfirsthead
\toprule
\textbf{Identyfikator} & \textbf{Poziom} & \textbf{Co wykrywa} & \textbf{Mechanizm} \\
\midrule
\endhead
\path{PORTCULLIS_P2_KERNEL_LOC_CEILING} & \texttt{error} & Suma linii plików \texttt{*.cs} we wszystkich folderach jądra przekracza sufit (domyślnie 800). & Konwencja \texttt{kernelFolders}, \texttt{kernelLineCeiling}; jedno zgłoszenie w~pierwszej linii największego pliku, komunikat podaje folder, sumę, sufit i~największy plik. \\
\path{PORTCULLIS_P2_KERNEL_ENTITY_REFERENCE} & \texttt{error} & Plik jądra odwołuje się do typu zadeklarowanego w~folderze encji. & Semantyczny: nazwa wiązana z~typem (albo składowa, lokalna, parametr tego typu), którego lokalizacja w~źródle leży w~\texttt{entityFolders}; raz na typ w~pliku. \\
\path{PORTCULLIS_P4_ENSURE_CREATED_OUTSIDE_TEST} & \texttt{error} & \texttt{X.Database.EnsureCreated()} lub \texttt{EnsureCreatedAsync()} bez nadrzędnego \texttt{if}, którego warunek wywołuje \texttt{IsInMemory()}. & Składniowy (typy EF Core nie rozwiązują się w~skanie). Strażnik w~stylu wczesnego \texttt{return} nie jest rozpoznawany. \\
\path{PORTCULLIS_P4_SEED_DATA_IN_MODEL} & \texttt{warning} & Wywołanie \texttt{HasData(...)} w~metodzie \texttt{OnModelCreating}. & Składniowy; komunikat podaje klasę. \\
\path{PORTCULLIS_P9_CONTROLLER_NO_DBCONTEXT} & \texttt{error} & Kontroler używa typu, którego nazwa kończy się na \texttt{DbContext}. & Składniowy. Kontroler: nazwa klasy kończy się na \texttt{Controller}, atrybut zawiera \texttt{ApiController} albo baza to \texttt{ControllerBase}/\texttt{Controller}. Raz na nazwę typu w~klasie. \\
\path{PORTCULLIS_P9_ORPHAN_ENTITY} & \texttt{warning} & Encja z~właściwości \texttt{DbSet<T>} klasy dziedziczącej po \texttt{DbContext}, zadeklarowana w~źródle, bez żadnego odwołania poza własnym plikiem. & Semantyczny; sama deklaracja \texttt{DbSet<T>} nie liczy się jako użycie. Zgłoszenie w~deklaracji encji. \\
\path{PORTCULLIS_P10_CUSTOM_BASE_CLASS} & \texttt{warning} & Klasa dziedziczy po klasie zadeklarowanej w~skanowanym kodzie. & Semantyczny: baza o~\texttt{TypeKind.Class}, różna od \texttt{object}, z~lokalizacją w~źródle. Bazy nierozwiązane (np. \texttt{ControllerBase} w~skanie CLI) i~bazy z~metadanych są pomijane. \\
\path{PORTCULLIS_P11_VENDOR_SDK_OUTSIDE_ADAPTER} & \texttt{warning} & Dyrektywa \texttt{using} importuje przestrzeń nazw dostawcy poza adapterem. & Składniowy; dopasowanie na granicy segmentu kropki, bez wielkości liter. Pomijane pliki w~\texttt{adapterFolders} i~pliki punktu wejścia (composition root). \\
\path{PORTCULLIS_P15_MISSING_SERVICE_DEFAULTS} & \texttt{warning} & Plik punktu wejścia buduje hosta, a~nigdzie nie wywołuje \texttt{AddServiceDefaults()}. & Składniowy: plik o~nazwie z~\texttt{entryPointFileNames}, wywołanie \texttt{CreateBuilder} lub \texttt{CreateApplicationBuilder}; zgłoszenie w~miejscu fabryki. \\
\path{PORTCULLIS_CONVENTION_UNMATCHED} & \texttt{warning} & Konwencja skonfigurowana jawnie nie pasuje do żadnej ścieżki skanu. & Kategoria \texttt{Meta}; bez lokalizacji, raz na konwencję. \\
\path{PORTCULLIS_NO_CONVENTION_MATCHED} & \texttt{warning} & Nic nie skonfigurowano i~żadna z~czterech konwencji nie pasuje. & Kategoria \texttt{Meta}; bez lokalizacji, najwyżej raz. \\
\bottomrule
\end{xltabular}
\endgroup

\paragraph{Szczegóły reguł zasad.}
\begin{itemize}[leftmargin=*]
  \item \textbf{P2.} Sufit liczy się dla wszystkich drzew składniowych, których ścieżka ma
        segment pasujący do \repopath{kernelFolders}, łącznie. Gdy żaden plik nie jest
        w~folderze jądra, reguła kończy pracę bez zgłoszeń -- tę ciszę raportuje analyzer
        pokrycia.
  \item \textbf{P4.} Rozpoznawany strażnik to dokładnie kształt z~repozytorium
        referencyjnego (przykład z~komentarza w~kodzie reguły):
\begin{lstlisting}
if (context.Database.IsInMemory()) { await context.Database.EnsureCreatedAsync(...); }
\end{lstlisting}
        Komunikat: \repopath{EnsureCreated} nie zapisuje migracji, więc schemat
        zamarza w~stanie z~pierwszego uruchomienia, gdy migracje przybywają w~kodzie.
  \item \textbf{P9.} \repopath{CONTROLLER\_NO\_DBCONTEXT} to reguła, którą SPEC wprost
        wskazuje dla P9. \repopath{ORPHAN\_ENTITY} jest własnym rozszerzeniem tej samej
        dyscypliny warstw na jedyną kategorię fixture'a sprawdzalną Roslynem (encje
        EF Core bez kontrolera, \repopath{CONSUMER-003}--\repopath{CONSUMER-011}); żaden
        wpis fixture'a nie cytuje P9 dosłownie i~źródła to zaznaczają. Reguła jest
        grubsza niż fixture: liczy „zero odwołań poza plikiem deklaracji”, a~nie „zero
        odwołań z~kontrolera”, więc właściwość nawigacyjna w~sąsiedniej klasie domeny
        oznacza „używana”.
  \item \textbf{P10.} Ograniczenie do baz zadeklarowanych w~źródle nie jest przypadkowe:
        typy ASP.NET Core i~EF Core nie rozwiązują się w~skanie, więc \repopath{: ControllerBase}
        czy \repopath{: DbContext} -- prawowite punkty rozszerzeń frameworku -- nie są
        zgłaszane.
  \item \textbf{P11.} Lista dostawców jest celowo krótka; reguła „każdy pakiet zewnętrzny”
        utonęłaby w~fałszywych alarmach na EF Core i~ASP.NET Core. Plik punktu wejścia
        (\repopath{Program.cs}) jest drugim uznanym miejscem adaptera od czasu, gdy
        pierwszy przebieg na repozytorium referencyjnym zgłosił w~nim idiomatyczną
        rejestrację klienta SDK w~DI (sekcja~\ref{sec:portcullis-mutacje}). Wcześniejsze
        dopasowanie tylko pierwszego segmentu sprawiało, że wpis wielosegmentowy
        \repopath{Google.Cloud} nigdy niczego nie łapał; zastąpiono je dopasowaniem na
        granicy segmentu (\repopath{Azure} łapie \repopath{Azure.Storage.Blobs}, nie łapie
        \repopath{AzureFoo}).
  \item \textbf{P15.} Plik bez wywołania fabryki hosta (np. zwykłe narzędzie konsolowe)
        nie jest punktem wejścia usługi i~nie jest sprawdzany.
\end{itemize}

\paragraph{Pokrycie konwencji (\texttt{Meta}).} Każda reguła konwencyjna reaguje na
konwencję, która nic nie dopasowała, wczesnym powrotem. Dla reguły to poprawne, dla bramki --
nie: zespół, którego jądro nazywa się \repopath{Common}, a~adaptery leżą w~\repopath{External/},
dostawał zieloną, pustą bramkę. \repopath{ConventionCoverageAnalyzer} sprawdza cztery
konwencje: foldery jądra, foldery encji, foldery adapterów i~nazwy plików punktu wejścia,
na ścieżkach względnych wobec katalogu skanu. \repopath{PORTCULLIS\_CONVENTION\_UNMATCHED}
zgłasza każdą konwencję, którą zespół skonfigurował sam i~która nic nie dopasowała (prawie
zawsze literówka albo folder przemianowany po napisaniu konfiguracji).
\repopath{PORTCULLIS\_NO\_CONVENTION\_MATCHED} strzela najwyżej raz, gdy żadnej z~tych
czterech nie skonfigurowano i~żadna nie pasuje -- wtedy reguły konwencyjne nie oceniły
niczego. Decyzje:

\begin{itemize}[leftmargin=*]
  \item raport łączny „wszystko albo nic”, bo wiele zdrowych usług nie ma np. wspólnego
        jądra; wcześniejsza wersja per konwencja dawała cztery ostrzeżenia na zwykłym kodzie
        i~uczyła wyłączania reguły;
  \item poziom \repopath{warning}, bo \repopath{error} zablokowałby build każdemu przy
        pierwszej instalacji;
  \item przestrzenie nazw dostawców nie są sprawdzane (brak importów SDK to nic
        podejrzanego); pusta kompilacja nie daje raportu; jawnie pusta lista konwencji
        liczy się jako spełniona;
  \item o~tym, czy zespół „coś skonfigurował”, decydują tylko klucze tych czterech
        konwencji -- samo ustawienie \repopath{kernelLineCeiling} nie wycisza raportu;
  \item diagnostyki nie mają lokalizacji; skaner mapuje je na pustą ścieżkę, a~komentarz na
        PR pokazuje je pod nagłówkiem „Repository-wide”.
\end{itemize}

\zrodla{\path{letsgolegacy.portcullis/README.md}, \path{letsgolegacy.portcullis/src/Portcullis.Rules/PACKAGE-README.md},
\path{letsgolegacy.portcullis/src/Portcullis.Rules/RuleRegistry.cs},
\path{letsgolegacy.portcullis/src/Portcullis.Rules/KernelBoundaryAnalyzer.cs},
\path{letsgolegacy.portcullis/src/Portcullis.Rules/PersistencePortabilityAnalyzer.cs},
\path{letsgolegacy.portcullis/src/Portcullis.Rules/ProgramManifestLayeringAnalyzer.cs},
\path{letsgolegacy.portcullis/src/Portcullis.Rules/ExtensibilityInheritanceAnalyzer.cs},
\path{letsgolegacy.portcullis/src/Portcullis.Rules/AntiCorruptionEdgeAnalyzer.cs},
\path{letsgolegacy.portcullis/src/Portcullis.Rules/ObservabilityBuildTimeAnalyzer.cs},
\path{letsgolegacy.portcullis/src/Portcullis.Rules/ConventionCoverageAnalyzer.cs},
\path{letsgolegacy.portcullis/docs/CONFIGURATION.md}, \path{letsgolegacy.portcullis/docs/MUTATIONS.md}}

\subsection{Reguły migracyjne}\label{sec:portcullis-migracja}

Cztery analyzery Roslyn dla kodu przeniesionego z~.NET Framework na nowoczesny .NET
(.NET~10). Każdy szuka idiomu, który migracja mechaniczna -- człowiek w~pośpiechu albo
agent migracyjny -- zwykle przenosi bez zmian: kod nadal się kompiluje (przez shim
zgodności albo pakiet NuGet) i~przechodzi testy, ale niesie założenie, które stary runtime
spełniał, a~nowy nie. Identyfikatory są kontraktem: pakiet standardów Second Key cytuje je
dosłownie, a~\repopath{RuleRegistryTests} je przypina. Wszystkie mają kategorię diagnostyki
\repopath{Migration}, bo żadna zasada \repopath{architecture-standards} nie nazywa tych
idiomów.

\begin{xltabular}{\linewidth}{@{}>{\raggedright\arraybackslash}p{2.9cm} >{\raggedright\arraybackslash}X@{}}
\toprule
\multicolumn{2}{@{}l@{}}{\textbf{Reguła, poziom domyślny i~opis}} \\
\midrule
\endfirsthead
\toprule
\multicolumn{2}{@{}l@{}}{\textbf{Reguła, poziom domyślny i~opis (cd.)}} \\
\midrule
\endhead
\multicolumn{2}{@{}l@{}}{\path{PORTCULLIS_MIG_SYSTEM_WEB} \quad poziom \texttt{warning}} \\*
Co wykrywa & Każde użycie przestrzeni nazw lub typu z~\texttt{System.Web} i~przestrzeni w~niej zagnieżdżonych: \texttt{using}, nazwa kwalifikowana, referencja typu. \\*
Ryzyko migracji & System.Web nie istnieje w~nowoczesnym .NET; kod buduje się tylko przez adaptery System.Web lub pozostałą referencję, a~potok żądań, moduły, sesja i~cache są emulowane, nie natywne. \\*
Przykład & \texttt{using System.Web;} i~parametr \texttt{HttpRequest request} (typ \path{System.Web.HttpRequest}). \\
\midrule
\multicolumn{2}{@{}l@{}}{\path{PORTCULLIS_MIG_HTTPCONTEXT_CURRENT} \quad poziom \texttt{error}} \\*
Co wykrywa & Odczyt lub zapis statycznego \path{System.Web.HttpContext.Current}, pod dowolną nazwą (alias, \texttt{using static}, pełna kwalifikacja). \\*
Ryzyko migracji & Działa tylko na przepływie żądania, chowa zależność od warstwy web w~kodzie, który wygląda na niezależny; na ASP.NET Core istnieje tylko przez adaptery System.Web. \\*
Przykład & \path{HttpContext.Current.User.Identity.Name} \\
\midrule
\multicolumn{2}{@{}l@{}}{\path{PORTCULLIS_MIG_SYNC_OVER_ASYNC} \quad poziom \texttt{warning}} \\*
Co wykrywa & Blokujące czekanie na task: \texttt{.Result}, \texttt{.Wait()}, \texttt{Task.WaitAll}, \texttt{Task.WaitAny}, \path{.GetAwaiter().GetResult()}. \\*
Ryzyko migracji & Pod kontekstem synchronizacji klasycznego ASP.NET zakleszcza; na ASP.NET Core trzyma wątek puli przez cały czas czekania, jeden na każde współbieżne żądanie -- usługa pada przy pierwszym realnym obciążeniu. \\*
Przykład & \path{_client.GetStringAsync(url).Result} \\
\midrule
\multicolumn{2}{@{}l@{}}{\path{PORTCULLIS_MIG_CONFIGURATION_MANAGER} \quad poziom \texttt{error}} \\*
Co wykrywa & Odczyt ustawień przez statyczne składowe \texttt{ConfigurationManager}, \texttt{WebConfigurationManager} lub \texttt{ConfigurationSettings} zamiast \texttt{IOptions<T>} albo \texttt{IConfiguration}. \\*
Ryzyko migracji & Pakiet zgodności kompiluje kod, ale czyta \texttt{<app>.dll.config}, nie \texttt{web.config} ani \texttt{appsettings.json}; ustawienie wraca w~runtime jako \texttt{null}, bez błędu builda i~testu. \\*
Przykład & \path{ConfigurationManager.AppSettings} z~kluczem ustawienia. \\
\bottomrule
\end{xltabular}

\subsubsection{Jak reguły decydują}

\textbf{Na symbolach, nigdy na tekście.} Typ liczy się jako System.Web, bo przestrzeń
nazw, w~której jest zadeklarowany, to System.Web; \repopath{.Result} się liczy, bo
właściwość deklaruje \repopath{Task<TResult>}; odczyt konfiguracji się liczy, bo składową
deklaruje \path{System.Configuration.ConfigurationManager}. Dzięki temu reguły łapią
idiom w~każdej postaci -- przez alias (\repopath{using Ctx = System.Web.HttpContext;}),
\repopath{using static}, pełną kwalifikację lub grupę metod -- i~zostawiają sobowtóry:
\path{Microsoft.AspNetCore.Http.HttpContext}, własne klasy zespołu \repopath{HttpContext}
czy \repopath{ConfigurationManager} oraz nowoczesny
\path{Microsoft.Extensions.Configuration.ConfigurationManager}, który zwraca
\repopath{WebApplicationBuilder.Configuration} (wyszukiwanie tekstowe oznaczyłoby go w~każdym
\repopath{Program.cs}).

Typy są identyfikowane po w~pełni kwalifikowanej nazwie metadanych
(metoda \path{MigrationSymbols.MetadataName}, np. \path{System.Threading.Tasks.Task`1}),
a~nie po zestawie, który je deklaruje, więc prawdziwy \repopath{System.Web.dll}, pakiet
adapterów i~ręcznie napisany shim są rozpoznawane tak samo. Kod celowo nie używa
\repopath{Compilation.GetTypeByMetadataName}: zwraca \repopath{null}, gdy dwa zestawy
deklarują tę samą nazwę -- czyli dokładnie w~pół-zmigrowanym rozwiązaniu z~System.Web
i~shimem. Kod reguł nie używa też zmiennych wzorca w~warunkach
(\repopath{x is T t \&\& ...}); powód opisuje sekcja~\ref{sec:portcullis-mutacje}.

\textbf{Oba kanały widzą te same symbole.} W~pakiecie analyzerów kompilator wiąże nazwy
z~prawdziwymi referencjami projektu. CLI, akcja i~kontener kompilują źródła same, bez
builda projektu, więc dodają zestawy frameworku i~deklaracje API .NET Framework
(sekcja~\ref{sec:portcullis-budowa}). Skan nie widzi natomiast globalnych
\repopath{using} z~\repopath{obj/}: w~pliku, który na nich polega, nazwa typu bez
\repopath{using} (\repopath{Task}, \repopath{HttpClient}) się nie wiąże, a~składowa
wyrażenia, którego typ pochodzi z~wywołania, nadal tak (\repopath{client.GetStringAsync(url).Result}
jest widoczne, parametr zadeklarowany jako \repopath{Task<int>} w~takim pliku -- nie).

\paragraph{Konfiguracja.}
\begin{xltabular}{\linewidth}{@{}>{\raggedright\arraybackslash}p{4.7cm} >{\raggedright\arraybackslash}p{3.0cm} >{\raggedright\arraybackslash}X@{}}
\toprule
\textbf{\texttt{portcullis.json} / \texttt{.globalconfig}} & \textbf{Domyślnie} & \textbf{Efekt} \\
\midrule
\endfirsthead
\toprule
\textbf{\texttt{portcullis.json} / \texttt{.globalconfig}} & \textbf{Domyślnie} & \textbf{Efekt} \\
\midrule
\endhead
\path{migrationExemptFolders} \newline \path{portcullis_migration_exempt_folders} & brak & Foldery, w~których żadna z~czterech reguł nic nie zgłasza: warstwa zgodności, która świadomie rozmawia z~System.Web w~trakcie migracji przyrostowej, albo projekt jeszcze niezmigrowany. \\
\path{systemWebAllowedTypes} \newline \path{portcullis_system_web_allowed_types} & \path{System.Web.HttpUtility}, \path{System.Web.IHtmlString} & W~pełni kwalifikowane typy System.Web, których \texttt{PORTCULLIS\_MIG\_SYSTEM\_WEB} nie zgłasza. \\
\bottomrule
\end{xltabular}

Foldery pasują tak jak w~każdej innej konwencji folderów: segment ścieżki równy nazwie albo
kończący się na \repopath{.Nazwa}, bez wielkości liter, względem katalogu skanu. Ustawienie
listy zastępuje wartość domyślną, a~jawnie pusta lista jest honorowana. Pusta domyślna
lista wyjątków ma uzasadnienie w~kodzie: wyjątek, o~który nikt nie prosił, to dziura
w~bramce. Poziom zmienia się zwykłą drogą w~pakiecie analyzerów:

\begin{lstlisting}
[*.cs]
dotnet_diagnostic.PORTCULLIS_MIG_SYNC_OVER_ASYNC.severity = error
\end{lstlisting}

\subsubsection{\texttt{PORTCULLIS\_MIG\_SYSTEM\_WEB} (warning)}

\textbf{Co wykrywa.} Każde użycie przestrzeni nazw lub typu z~\repopath{System.Web} albo
przestrzeni w~niej zagnieżdżonej: dyrektywę \repopath{using} (także \repopath{using static}
i~\repopath{using X = ...}), nazwę kwalifikowaną, referencję typu -- typy parametrów i~pól,
\repopath{new}, \repopath{typeof}, \repopath{catch}, argumenty generyczne, atrybuty.

\textbf{Dlaczego to ryzyko.} System.Web to ASP.NET z~.NET Framework i~nie istnieje
w~nowoczesnym .NET. Kod, który go nazywa, buduje się tylko przez adaptery System.Web albo
przez referencję pozostałą po starym projekcie; w~obu przypadkach potok żądań, moduły,
handlery, sesja i~cache są emulowane, a~nie natywne. Każde trafienie to fragment migracji
przeniesiony, a~nie wykonany.

\begin{lstlisting}[escapeinside={(*}{*)}]
// Reported: 'System.Web.HttpRequest' is System.Web (*\ldots*)
using System.Web;

public class AuditService
{
    public void Record(HttpRequest request) => Write(request.UserHostAddress);
}
\end{lstlisting}

Ta sama klasa z~dyrektywą \repopath{using Microsoft.AspNetCore.Http;}, czyli z~typem
\repopath{HttpRequest} z~ASP.NET Core, nie jest zgłaszana.

Jedna referencja to jedno zgłoszenie: \repopath{System.Web.HttpContext.Current} jest
zgłaszane raz, na \repopath{System.Web.HttpContext} (najbardziej zewnętrzna nazwa, która
nadal jest przestrzenią lub typem System.Web). Nazwa, której przestrzeń System.Web nie
rozwiązuje się w~skanie (\repopath{System.Web.Mvc.Controller}), jest zgłaszana na prefiksie
\repopath{System.Web}, a~nie gubiona.

\textbf{Celowo niezgłaszane.}
\begin{itemize}[leftmargin=*]
  \item Samo \repopath{using System.Web;}: leżą tam \repopath{System.Web.HttpUtility}
        i~\repopath{System.Web.IHtmlString}, które nowoczesny .NET dostarcza
        (\repopath{System.Web.HttpUtility.dll}). Typy legacy używane przez tę dyrektywę są
        zgłaszane w~miejscu użycia. Po usunięciu obu z~\repopath{systemWebAllowedTypes}
        zgłaszana jest też sama dyrektywa.
  \item Nazwy w~deklaracji \repopath{namespace} (deklarowanie typów w~System.Web to jego
        implementacja, jak w~shimie) i~nazwy w~komentarzach dokumentacji.
  \item Tylko w~skanie CLI: nazwa niekwalifikowana w~pozycji wyrażenia wewnątrz klasy,
        której klasy bazowej skan nie rozwiązuje. Chodzi o~zmigrowany kontroler ASP.NET Core:
        \repopath{HttpContext.Request} to tam \repopath{ControllerBase.HttpContext}, którego
        skan nie widzi; zgłoszenie byłoby fałszywym alarmem na dokładnie tym kodzie, który
        produkuje migracja. Nazwy kwalifikowane i~nazwy w~pozycji typu są w~takiej klasie
        nadal zgłaszane.
\end{itemize}

\textbf{Znane luki.} W~skanie CLI typ z~\repopath{System.Web.Mvc}, \repopath{System.Web.UI},
\repopath{System.Web.Http} i~innych niezadeklarowanych przestrzeni jest łapany w~dyrektywie
\repopath{using} i~na kwalifikowanym prefiksie \repopath{System.Web}, ale nie w~miejscu
użycia po krótkiej nazwie. Użycie, które nigdy nie zapisuje nazwy System.Web
(\repopath{var context = GetContext();}, a~potem \repopath{context.Request}), jest
zgłaszane tylko tam, gdzie \repopath{GetContext} nazywa typ. Wszystko, co nie jest C\#, zostaje poza zasięgiem:
\repopath{.cshtml}, \repopath{.aspx}, \repopath{.ascx}, \repopath{Global.asax},
\repopath{web.config}, a~także nazwy w~łańcuchach znaków (np. nazwa typu System.Web
przekazana jako tekst do \repopath{Type.GetType}).

\subsubsection{\texttt{PORTCULLIS\_MIG\_HTTPCONTEXT\_CURRENT} (error)}

\textbf{Co wykrywa.} Odczyt lub zapis statycznego \repopath{System.Web.HttpContext.Current},
jakkolwiek nazwanego: przez \repopath{using System.Web;}, w~pełni kwalifikowanego, przez
alias lub \repopath{using static}. Implementacja działa na drzewie operacji: każda
referencja właściwości jest sprawdzana, czy to statyczne \repopath{Current} typu o~nazwie
metadanych \repopath{System.Web.HttpContext}.

\textbf{Dlaczego to ryzyko.} \repopath{HttpContext.Current} to sposób, w~jaki kod
.NET Framework sięga po bieżące żądanie z~dowolnego miejsca -- z~repozytorium, helpera,
loggera. Działa tylko na przepływie żądania, chowa zależność od warstwy web w~kodzie, który
wygląda na niezależny od frameworku, a~na ASP.NET Core istnieje tylko przez adaptery
System.Web. To idiom, który przeżywa resztę System.Web: zespół może przenieść każdy typ na
odpowiednik ASP.NET Core i~nadal czytać stan żądania przez statyk. Dlatego to osobna reguła,
a~linia z~\repopath{HttpContext.Current} dostaje obie diagnostyki (tę i~\repopath{MIG\_SYSTEM\_WEB}):
pierwszą zespół może spłacić przejściem na adaptery, drugą nadal będzie winien.

\begin{lstlisting}[escapeinside={(*}{*)}]
// Reported: 'HttpContext.Current' reaches the current request through the static (*\ldots*)
public class OrderService
{
    public void Place(Order order) => order.PlacedBy = HttpContext.Current.User.Identity.Name;
}
\end{lstlisting}

Wersja zgodna z~regułą ma jawną zależność: usługa dostaje \repopath{IHttpContextAccessor}
przez konstruktor i~da się ją testować bez żądania. Jeszcze lepiej przekazać użytkownika
z~endpointu -- wtedy usługa nie potrzebuje żadnego typu web.

\textbf{Celowo niezgłaszane.} Własne statyczne \repopath{Current} zespołu na własnym
\repopath{HttpContext}; składowe instancyjne \repopath{HttpContext}, \repopath{HttpContextBase}
i~\repopath{HttpContextWrapper} (te należą do \repopath{MIG\_SYSTEM\_WEB});
\repopath{nameof(HttpContext.Current)}, które niczego nie czyta.

\textbf{Znane luki.} \repopath{IHttpContextAccessor} użyty jako service locator głęboko
w~warstwie domeny ma ten sam kształt problemu, ale to nowoczesne API i~nie jest zgłaszany.
Własny statyczny wrapper zespołu wokół \repopath{HttpContext.Current} jest zgłaszany raz,
wewnątrz wrappera; jego wywołania nie. Refleksja nie jest widoczna.

\subsubsection{\texttt{PORTCULLIS\_MIG\_SYNC\_OVER\_ASYNC} (warning)}

\textbf{Co wykrywa.} Blokowanie wątku na tasku:
\begin{itemize}[leftmargin=*]
  \item \repopath{.Result} na \repopath{Task<T>} i~\repopath{ValueTask<T>};
  \item \repopath{.Wait()} w~każdym przeciążeniu na \repopath{Task}/\repopath{Task<T>} oraz
        statyczne \repopath{Task.WaitAll} i~\repopath{Task.WaitAny};
  \item \path{.GetAwaiter().GetResult()} na czterech typach tasków (\repopath{Task},
        \repopath{Task<T>}, \repopath{ValueTask}, \repopath{ValueTask<T>}), także po
        \repopath{ConfigureAwait(...)};
  \item \repopath{GetResult()} na każdym z~ośmiu typów awaiterów tych tasków trzymanym
        w~zmiennej.
\end{itemize}

\textbf{Dlaczego to ryzyko.} To idiom, który migracja mechaniczna produkuje najczęściej:
synchroniczne miejsce wywołania spotyka API, które w~nowoczesnym .NET jest tylko
asynchroniczne (\repopath{HttpClient}, \repopath{Stream}, EF Core), i~dostaje doklejone
\repopath{.Result}. Pod kontekstem synchronizacji klasycznego ASP.NET to zakleszczenie; na
ASP.NET Core zakleszczenia nie ma, ale wątek puli jest zajęty przez cały czas czekania --
jeden na każde współbieżne żądanie. Tak usługa, która przeszła wszystkie testy, pada przy
pierwszym realnym obciążeniu.

\begin{lstlisting}[escapeinside={(*}{*)}]
// Reported: '_client.GetStringAsync(url).Result' blocks the calling thread (*\ldots*)
public string Rates(string url) => _client.GetStringAsync(url).Result;
\end{lstlisting}

Wersja zgodna z~regułą to metoda \repopath{async} zwracająca \repopath{Task<string>},
która robi \repopath{await} na wywołaniu klienta.

\textbf{Celowo niezgłaszane} -- dwa kształty, w~których reguła może udowodnić z~kodu przed
sobą, że task już się zakończył:
\begin{itemize}[leftmargin=*]
  \item antecedent wewnątrz kontynuacji:
\begin{lstlisting}
task.ContinueWith(previous => previous.Result)
\end{lstlisting}
  \item synchroniczna szybka ścieżka: ta sama zmienna lokalna, parametr, pole lub właściwość
        na gałęzi \repopath{true} warunku, który jako koniunkcja wymaga
        \repopath{IsCompleted} lub \repopath{IsCompletedSuccessfully}:
\begin{lstlisting}
if (task.IsCompletedSuccessfully) return task.Result;
\end{lstlisting}
        Zaprzeczone
        sprawdzenie (\repopath{!task.IsCompleted}), sprawdzenie w~alternatywie, gałąź
        \repopath{else} albo sprawdzenie innego taska niczego nie dowodzą i~są zgłaszane.
\end{itemize}
Nie są też zgłaszane rzeczy, które tylko wyglądają na taski: \repopath{SemaphoreSlim.Wait()},
\repopath{ManualResetEventSlim.Wait()}, \repopath{Monitor.Wait},
\path{Task.Yield().GetAwaiter().GetResult()}, własna właściwość \repopath{Result}
zespołu, a~także \repopath{nameof(Task<int>.Result)}.

\textbf{Znane fałszywe alarmy} (task zakończony, ale dowód wymaga przepływu sterowania,
którego reguła nie śledzi): \repopath{.Result} po \repopath{await Task.WhenAll(...)}
(zalecenie: zrobić \repopath{await} na zakończonym tasku, to nic nie kosztuje), po wczesnym
\repopath{return} na \repopath{!IsCompleted}, po sprawdzeniu
\path{Status == TaskStatus.RanToCompletion} oraz w~kontynuacji
\path{Task.Factory.ContinueWhenAll} lub \path{ContinueWhenAny} (wyjątek obejmuje tylko
\repopath{ContinueWith}). Blokujące czekanie w~miejscu, które nie może być asynchroniczne
(konstruktor, \repopath{Dispose}, getter), nadal jest blokującym czekaniem i~jest zgłaszane.

\textbf{Znane luki.} W~skanie CLI task zwracany przez API pakietu NuGet
(\repopath{query.ToListAsync().Result} z~EF Core) nie jest widoczny, bo pakiet nie jest
referencjonowany; API frameworku i~własne metody asynchroniczne kodu są widoczne, a~pakiet
analyzerów widzi wszystko. Również tylko w~CLI plik polegający na niejawnych globalnych
\repopath{using} traci kształty wymagające samej nazwy \repopath{Task}:
\repopath{Task.WaitAll(...)} i~\repopath{.Result} na parametrze lub polu typu
\repopath{Task<T>}. Widać to w~skanie samego Portcullis: w~\repopath{GitCommandRunner.cs}
dwa odczyty \repopath{.Result} na taskach z~\repopath{StreamReader.ReadToEndAsync()} są
zgłaszane (to dwa wpisy baseline repozytorium), a~poprzedzające je \repopath{Task.WaitAll}
nie. Metody i~lambdy \repopath{async void} to inne zagrożenie i~nie należą do tej reguły.

\subsubsection{\texttt{PORTCULLIS\_MIG\_CONFIGURATION\_MANAGER} (error)}

\textbf{Co wykrywa.} Odczyt ustawień przez API konfiguracji .NET Framework, czyli każdą
statyczną składową jednego z~trzech typów -- jako odczyt właściwości, wywołanie, grupę
metod albo przez \repopath{using static}:
\begin{itemize}[leftmargin=*]
  \item \path{System.Configuration.ConfigurationManager} (\repopath{AppSettings},
        \repopath{ConnectionStrings}, \repopath{GetSection}, \repopath{OpenExeConfiguration}
        i~inne);
  \item jej bliźniak z~ASP.NET: \path{System.Web.Configuration.WebConfigurationManager};
  \item dawno przestarzały \path{System.Configuration.ConfigurationSettings}.
\end{itemize}
Komunikat podaje wyrażenie i~typ deklarujący.

\textbf{Dlaczego to ryzyko.} To defekt, który przeżywa każdy build. Pakiet
\path{System.Configuration.ConfigurationManager} sprawia, że kod kompiluje się na
nowoczesnym .NET, ale czyta tam \repopath{<app>.dll.config} -- nie \repopath{web.config},
nigdy \repopath{appsettings.json} ani zmienne środowiskowe -- więc ustawienie, które kiedyś
przychodziło, wraca w~runtime jako \repopath{null}, bez błędu builda i~bez nieudanego testu
jednostkowego.

\begin{lstlisting}[escapeinside={(*}{*)}]
// Reported: 'ConfigurationManager.AppSettings' reads settings through System.Configuration.ConfigurationManager (*\ldots*)
public static string Currency() => ConfigurationManager.AppSettings["Currency"];
\end{lstlisting}

Wersja zgodna z~regułą wiąże klasę opcji jeden raz, przy starcie aplikacji i~z~walidacją.
Rejestracja to wywołanie \repopath{AddOptions<ShopOptions>()}, po nim \repopath{Bind} sekcji
konfiguracji i~\repopath{ValidateOnStart}; odbiorcy dostają opcje przez wstrzyknięcie jako
\repopath{IOptions<ShopOptions>}.

\textbf{Celowo niezgłaszane.} \path{Microsoft.Extensions.Configuration.ConfigurationManager}
(nowoczesna klasa o~tej samej krótkiej nazwie), \repopath{IConfiguration}, własny
\repopath{ConfigurationManager} zespołu oraz \repopath{typeof(ConfigurationManager)}
i~\repopath{nameof(...)}, które niczego nie czytają.

\textbf{Znane luki.} Składowe obiektu \repopath{Configuration} zwracanego przez
\repopath{OpenExeConfiguration} nie są zgłaszane osobno (samo wywołanie jest). Klasy
ustawień z~designera dziedziczące po \repopath{ApplicationSettingsBase}
(\repopath{Properties.Settings.Default.X}) też czytają \repopath{app.config} i~nie są
zgłaszane. Deklaracja własnej \repopath{ConfigurationSection} nie jest odczytem.

\subsubsection{Poziom a~bramka i~dowody}

Dwie reguły, których trafienie to defekt runtime'u na nowoczesnym .NET (stan żądania przez
statyk i~ustawienia wracające jako \repopath{null}), mają domyślnie \repopath{error}. Dwie,
których trafienie to dług, który nadal może działać (System.Web przez shim i~blokujące
czekanie), mają \repopath{warning}. Na bramce PR pochodzenie zmienia to dla kodu AI:
ostrzeżenie na linii przypisanej agentowi AI jest eskalowane do błędu, więc na migracyjnym
PR napisanym przez agenta blokują wszystkie cztery reguły -- i~tylko na liniach, które ten PR
zmienił.

\finding{Na pull requeście napisanym przez agenta migracyjnego wszystkie cztery reguły
migracyjne blokują, ale wyłącznie na zmienionych liniach.}

Każda reguła ma testy w~obu kierunkach w~\repopath{tests/Portcullis.Engine.Tests/Rules/},
skan end-to-end przez ścieżkę CLI (\repopath{ScannerMigrationTests}), ręcznie napisane
mutanty i~przebieg Stryker.NET (sekcja~\ref{sec:portcullis-mutacje}). Oba kanały
uruchomiono naprawdę 2026-09-26. Pakiet analyzerów, spakowany i~użyty z~jednorazowego
projektu web \repopath{net10.0} z~prawdziwymi pakietami zgodności
(\path{Microsoft.AspNetCore.SystemWebAdapters} 2.3.0,
\path{System.Configuration.ConfigurationManager} 10.0.12), zgłosił:

\begin{lstlisting}[escapeinside={(*}{*)}]
LegacyOrders.cs(9,34):  error   PORTCULLIS_MIG_CONFIGURATION_MANAGER: 'LegacyConfig.AppSettings' reads settings through System.Configuration.ConfigurationManager. (*\ldots*)
LegacyOrders.cs(11,30): error   PORTCULLIS_MIG_HTTPCONTEXT_CURRENT: 'System.Web.HttpContext.Current' reaches the current request through the static (*\ldots*)
LegacyOrders.cs(11,30): warning PORTCULLIS_MIG_SYSTEM_WEB: 'System.Web.HttpContext' is System.Web, which modern .NET does not have: (*\ldots*)
LegacyOrders.cs(13,30): warning PORTCULLIS_MIG_SYNC_OVER_ASYNC: '_client.GetStringAsync("https://example.test/rates").Result' blocks the calling thread (*\ldots*)
\end{lstlisting}

\path{System.Web.HttpUtility.UrlEncode(...)} i~\repopath{builder.Configuration[...]}
(nowoczesny \repopath{ConfigurationManager}) w~tym samym projekcie nie zostały zgłoszone.
Przebieg pokazał też, dlaczego reguły są semantyczne: z~niejawnymi \repopath{using} Web SDK
gołe \repopath{ConfigurationManager} nawet się nie kompiluje (\repopath{CS0104},
niejednoznaczność między dwiema klasami), więc odczyt legacy trzeba było zapisać przez
alias -- a~reguła i~tak go złapała. \repopath{portcullis scan} na katalogu podobnego kodu
legacy, bez żadnego builda, zgłosił wszystkie cztery reguły; ten skan jest testem
\repopath{ScannerMigrationTests} (z~oczekiwanymi liniami i~poziomami, a~w~drugim teście
z~\repopath{migrationExemptFolders} wyciszającym tylko wskazany folder). To jeszcze nie są
dowody z~prawdziwego zmigrowanego kodu -- to ticket P8 w~benchu nopCommerce
(sekcja~\ref{sec:bench}).

\zrodla{\path{letsgolegacy.portcullis/docs/rules/MIGRATION.md},
\path{letsgolegacy.portcullis/src/Portcullis.Rules/MigrationSymbols.cs},
\path{letsgolegacy.portcullis/src/Portcullis.Rules/SystemWebUsageAnalyzer.cs},
\path{letsgolegacy.portcullis/src/Portcullis.Rules/HttpContextCurrentAnalyzer.cs},
\path{letsgolegacy.portcullis/src/Portcullis.Rules/SyncOverAsyncAnalyzer.cs},
\path{letsgolegacy.portcullis/src/Portcullis.Rules/ConfigurationManagerAnalyzer.cs},
\path{letsgolegacy.portcullis/src/Portcullis.Rules/PortcullisConventions.cs},
\path{letsgolegacy.portcullis/src/Portcullis.Engine/Semantics/ScanReferences.cs},
\path{letsgolegacy.portcullis/tests/Portcullis.Engine.Tests/ScannerMigrationTests.cs},
\path{letsgolegacy.portcullis/tests/Portcullis.Engine.Tests/Rules/RuleRegistryTests.cs},
\path{letsgolegacy.portcullis/portcullis-baseline.json}, \path{letsgolegacy.portcullis/docs/WORKPLAN.md}}

\subsection{Diff gate: bramka zawężona do zmienionych linii}\label{sec:portcullis-diff-gate}

Do milestone'u M5 CLI i~narzędzie komentarza liczyły kod wyjścia jako
\repopath{errorCount > 0} dla całego skanowanego drzewa. Repozytorium referencyjne miało
4~istniejące błędy \path{PORTCULLIS_P9_CONTROLLER_NO_DBCONTEXT}, więc bramka na nim
blokowałaby każdy PR, także taki, który niczego w~pobliżu nie dotyka. M4 nazwał tę lukę,
M5 odtworzył ją deterministycznie w~dwóch commitach, M5A (\repopath{docs/DIFF-GATE.md})
ją zamknął.

\finding{Wszystko, co skan znajdzie, jest raportowane; blokować może tylko błąd na linii
zmienionej przez PR, którego nie akceptuje baseline.}

\paragraph{Werdykt.} \repopath{ScanResult} dostał addytywne pole \repopath{gate}
(\repopath{GateResult}): \repopath{blocked}, \repopath{scope}, \repopath{blockingErrorCount},
później \repopath{degradedReason} i~\repopath{acceptedByBaselineCount}. \repopath{scope}
\repopath{all} to pierwotne zachowanie (każdy błąd w~skanie blokuje) i~obowiązuje, gdy nie
podano zakresu commitów; \repopath{scope} \repopath{diff} obowiązuje przy
\repopath{--provenance-range} i~liczy tylko błędy leżące na zmienionych liniach. Werdykt
powstaje w~jednym miejscu, w~\repopath{Scanner}; CLI i~\repopath{Portcullis.CiComment}
zwracają kod wyjścia z~\repopath{Gate?.Blocked ?? Summary.ErrorCount > 0} (fallback dotyczy
tylko historycznego JSON bez pola \repopath{gate}).

\paragraph{Skąd biorą się zmienione linie.} \repopath{--provenance-range <commitOrRange>}
wskazuje zakres, \repopath{--provenance-repo <ścieżka>} -- repozytorium git (domyślnie
skanowana ścieżka; trzeba go podać, gdy skanuje się podkatalog, np. \repopath{<repo>/src}).
\repopath{GitProvenanceProvider} wykonuje:

\begin{enumerate}[leftmargin=*]
  \item \textbf{Rozwiązanie zakresu.} \repopath{A..B} i~\repopath{A...B} są porównywane
        wprost jako dwa refy (bez semantyki merge-base dla trzech kropek -- świadome
        uproszczenie). Pojedynczy commit to zakres od pierwszego rodzica (dla merge'a --
        linia główna, jak widok PR w~GitHub) albo od pustego drzewa dla commita bez rodzica.
  \item \textbf{Diff.} \repopath{git diff --no-color -U0 <baza> <czubek>}. Przy
        \repopath{-U0} każdy hunk to dokładnie dodane lub zmienione linie; z~nagłówka
        \repopath{@@ -a,b +start,count @@} powstaje zakres \repopath{start..start+count-1}.
        Hunki czystego usunięcia (\repopath{count} 0) i~sekcje binarne nie dają zakresów.
        Nagłówek pliku to \repopath{+++ b/<ścieżka>}, ta sama ścieżka w~cudzysłowie (git
        cytuje ścieżki z~cudzysłowem, ukośnikiem wstecznym lub znakiem sterującym; ósemkowe
        sekwencje są dekodowane jako bajty UTF-8) albo \repopath{+++ /dev/null} dla
        usuniętego pliku.
  \item \textbf{Autorstwo.} Dla każdego zakresu \repopath{git blame --porcelain} z~opcją
        \repopath{-L start,koniec} na czubku zakresu; kolejne linie z~tego samego commita
        łączą się w~jeden zakres.
        Dzięki blame zakres wielocommitowy przypisuje linię commitowi, którego zmiana
        przetrwała. Porażka blame daje zakres ze źródłem \repopath{unknown}.
  \item \textbf{Klasyfikacja commita} (raz na commit) z~\repopath{git log -1} (autor,
        e-mail, trailery) przez \repopath{AiToolClassifier}.
\end{enumerate}

Każde wywołanie git dostaje \repopath{-c core.quotePath=false}
i~\repopath{-c diff.noprefix=false}, niezależnie od konfiguracji repozytorium -- inaczej ścieżki nie-ASCII i~brak prefiksu
\repopath{b/} po cichu psułyby parsowanie. Proces startuje z~listą argumentów
(\repopath{ProcessStartInfo.ArgumentList}), bez powłoki, więc komunikaty commitów, refy
i~ścieżki trafiają do git jako dosłowne argumenty; stdout i~stderr są czytane równolegle.

\paragraph{Klasyfikacja commitów.} Autorstwo pochodzi wyłącznie ze strukturalnych sygnałów
git, nigdy ze stylu kodu. Klasyfikator zna trzy narzędzia (identyfikatory w~kodzie to m.in.
\repopath{github-copilot} i~\repopath{cursor}) i~rozpoznaje je po domenie lub adresie
e-mail, nazwie autora oraz trailerach charakterystycznych dla narzędzia
(np. \repopath{Made-with: Cursor}).

\begin{xltabular}{\linewidth}{@{}>{\raggedright\arraybackslash}X l l >{\raggedright\arraybackslash}p{4.2cm}@{}}
\toprule
\textbf{Sygnał} & \textbf{Źródło} & \textbf{Pewność} & \textbf{Metoda} \\
\midrule
\endfirsthead
\toprule
\textbf{Sygnał} & \textbf{Źródło} & \textbf{Pewność} & \textbf{Metoda} \\
\midrule
\endhead
Trailer \texttt{Co-authored-by} z~tożsamością narzędzia AI albo nazwany trailer narzędzia, brak współautora-człowieka & \texttt{ai} & 0,85 & \texttt{commit-trailer} \\
Sygnał narzędzia AI i~co najmniej jeden współautor-człowiek w~\texttt{Co-authored-by} & \texttt{mixed} & 0,7 & \texttt{commit-trailer} \\
Autor commita to tożsamość narzędzia AI (konto bota, domena e-mail narzędzia) & \texttt{ai} & 0,95 & \texttt{git-blame-authorship} \\
Nieznane konto o~nazwie kończącej się na \texttt{[bot]} & \texttt{unknown} & 0,5 & \texttt{git-blame-authorship} \\
Brak sygnału AI & \texttt{human} & 0,7 & \texttt{git-blame-authorship} \\
Porażka blame albo odczytu commita & \texttt{unknown} & 0,0 & \texttt{no-attribution-signal} \\
\bottomrule
\end{xltabular}

\paragraph{Eskalacja poziomu.} Tylko zakresy \repopath{ai} eskalują, o~jeden stopień:
\repopath{info} \textrightarrow{} \repopath{warning}, \repopath{warning} \textrightarrow{}
\repopath{error}; \repopath{error} zostaje. \repopath{mixed} celowo nie eskaluje:
współautor-człowiek jest traktowany jak przegląd, a~nie jak nienadzorowany wynik AI (decyzja
zapisana jako polityka, do rewizji). Eskalacja jest wykonywana raz, w~\repopath{Scanner},
jako przebieg po zebranych naruszeniach, a~nie w~każdym analyzerze (sekcja~\ref{sec:portcullis-m4}).

\paragraph{Jedna flaga, dwa skutki.} \repopath{--provenance-range} steruje eskalacją
(tylko zakresy \repopath{ai}) i~zakresem bramki (wszystkie zakresy, bez względu na
źródło: „czy linia była częścią diffu” nie zależy od autora). Eskalacja może tylko
zaostrzyć poziom, więc wspólna flaga nigdy nie daje gorszego wyniku niż dwie osobne.
Bramka liczy się po eskalacji, więc eskalowane ostrzeżenie blokuje. O~tym, czy naruszenie
leży na zmienionej linii, decyduje wyłącznie \repopath{ChangedLines}: ścieżka naruszenia
(względem katalogu skanu) i~ścieżka zakresu (względem katalogu repozytorium) są
zamieniane na bezwzględne i~porównywane bez wielkości liter, a~linia musi leżeć w~zakresie
obustronnie domkniętym. Finding bez pliku (całe repozytorium) nigdy nie leży na zmienionej
linii. Tej samej klasy używają bramka, eskalacja i~filtr SARIF, więc SARIF nie może się
z~bramką nie zgodzić.

\paragraph{Fail-safe zamiast fail-open.} \repopath{GitProvenanceProvider} łapał błąd git
i~zwracał \emph{pusty} raport. Przy bramce zawężonej do diffu każda awaria git (płytki
klon, uszkodzony checkout, base przepisany force-pushem, katalog, który nie jest
repozytorium, brak binarki git) dawała zieloną bramkę. Odtworzono to przed poprawką na
katalogu z~jednym błędem i~bez repozytorium git:

\begin{lstlisting}
$ portcullis scan "$SB" --provenance-range HEAD~1..HEAD --provenance-repo "$SB"
CLI exit code = 0
gate: {'blocked': False, 'scope': 'diff', 'blockingErrorCount': 0}
summary: {'errorCount': 1, 'warningCount': 0, 'infoCount': 0}
stderr:
\end{lstlisting}

Poprawka wprowadza \repopath{ProvenanceStatus}: \repopath{Complete} albo
\repopath{Degraded}. Raport jest zdegradowany, gdy git nie rozwiązał zakresu albo gdy
diff zawierał hunk bez parsowalnego nagłówka \repopath{+++} (git kończy się wtedy
sukcesem, a~jedynym objawem byłaby krótsza lista zakresów, czyli łagodniejsza bramka).
\repopath{Scanner} traktuje zdegradowany raport jak brak zakresu: bramka wraca do
\repopath{scope} \repopath{all} z~powodem w~\repopath{degradedReason}, CLI pisze dwa
ostrzeżenia na stderr, SARIF dostaje notyfikację, a~komentarz na PR -- adnotację.
\repopath{ChangedLines} odmawia zbudowania się ze zdegradowanego raportu (wyjątek), więc
tego fallbacku nie da się ominąć przez pomyłkę. Fallback nigdy nie jest łagodniejszy od
zakresu, który zastępuje: zakres pusty blokuje na 0~naruszeniach, cały skan na N~$\geq$~0.

\begin{lstlisting}
$ portcullis scan "$SB" --provenance-range HEAD~1..HEAD --provenance-repo "$SB"
CLI exit code = 1
gate: {'blocked': True, 'scope': 'all', 'blockingErrorCount': 1,
       'degradedReason': "git could not resolve 'HEAD~1..HEAD' in '...': ... Could not access 'HEAD~1'"}
stderr:
portcullis: warning: could not scope the gate to 'HEAD~1..HEAD' — git could not resolve ...
portcullis: warning: falling back to gating on the whole scan (Gate.Scope "all").
\end{lstlisting}

Degradacja dotyczy zbioru linii, nie atrybucji. Porażka \repopath{git blame} lub
klasyfikacji commita nie degraduje raportu: zbiór zmienionych linii jest wtedy znany
i~poprawny, traci się tylko informację, kto je napisał, a~to wpływa wyłącznie na eskalację,
która może tylko zaostrzać. Degradacja w~tym miejscu przełączałaby cały skan na bramkę
bezwzględną z~powodu jednego nieczytelnego commita i~przywracała blokowanie na starym
długu. Przy okazji zamknięto dwa sąsiednie błędy: wyjątek z~\repopath{git log}
w~klasyfikacji zabijał proces (teraz daje klasyfikację bez atrybucji), a~parser po cichu
gubił hunki z~nieczytelnym nagłówkiem (teraz ustawia \repopath{SawUnattributedHunk}).
Naiwny test „jest patch, nie ma zakresów” byłby błędny: czyste usunięcie i~commit tylko
binarny dają dokładnie to, poprawnie. Powód degradacji jest spłaszczany do jednej linii,
bo trafia dosłownie do komentarza Markdown.

\paragraph{Kiedy PR jest blokowany.} Bramka blokuje, gdy istnieje co najmniej jeden
finding, który jednocześnie:
\begin{enumerate}[leftmargin=*]
  \item ma poziom \repopath{error} po eskalacji (domyślny \repopath{error} albo ostrzeżenie
        na linii przypisanej AI);
  \item leży na linii zmienionej w~zakresie commitów -- albo gdziekolwiek w~skanie, gdy
        zakresu nie podano lub raport pochodzenia jest zdegradowany;
  \item nie jest akceptowany przez baseline (sekcja~\ref{sec:portcullis-baseline}).
\end{enumerate}
W~GitHub Actions krok \repopath{actions/checkout} musi mieć \repopath{fetch-depth: 0}:
płytki klon nie zawiera commita bazowego PR, więc git zawodzi i~bramka (poprawnie, ale
surowiej) przechodzi na cały skan.

\paragraph{Weryfikacja.} \repopath{ScannerGateTests} sprawdza na prawdziwych regułach
m.in.: zgodność bez pochodzenia z~bramką bezwzględną; błąd poza zakresem nie blokuje, choć
jest raportowany; błąd w~zakresie blokuje; ostrzeżenie eskalowane przez zakres AI
blokuje; błąd w~zakresie autorstwa człowieka blokuje (bramka nie zależy od źródła);
kompletny, pusty zakres nie blokuje; zdegradowany raport wraca do bramki bezwzględnej.
Scenariusz 4~z~M5 (czysty commit E po commicie D z~błędem) z~\repopath{--provenance-range
D..E} daje \repopath{blocked: false}, \repopath{scope: diff} i~kod 0 (wcześniej 1).
Najmocniejszy dowód to prawdziwy, już scalony PR repozytorium referencyjnego (dodanie
klienta web), przeskanowany względem pierwszego rodzica: \repopath{errorCount: 9},
\repopath{blockingErrorCount: 5}. Pięć ostrzeżeń \path{PORTCULLIS_P10_CUSTOM_BASE_CLASS}
na nowych klasach klienta, napisanych przez agenta AI, eskalowało do błędów i~blokuje;
cztery stare błędy \repopath{CONTROLLER\_NO\_DBCONTEXT} poza diffem nie. Cały łańcuch
z~\repopath{Portcullis.CiComment} wyrenderował linię „Blocking this PR — 5 errors within
this PR's own changed lines.” i~zakończył się kodem 1.

\zrodla{\path{letsgolegacy.portcullis/docs/DIFF-GATE.md}, \path{letsgolegacy.portcullis/docs/SPEC.md},
\path{letsgolegacy.portcullis/docs/M4-INTEGRATION.md}, \path{letsgolegacy.portcullis/docs/FINDINGS.md},
\path{letsgolegacy.portcullis/src/Portcullis.Engine/Scanner.cs},
\path{letsgolegacy.portcullis/src/Portcullis.Engine/Model/ScanResult.cs},
\path{letsgolegacy.portcullis/src/Portcullis.Engine/Provenance/ChangedLines.cs},
\path{letsgolegacy.portcullis/src/Portcullis.Engine/Provenance/GitProvenanceProvider.cs},
\path{letsgolegacy.portcullis/src/Portcullis.Engine/Provenance/GitOutputParsers.cs},
\path{letsgolegacy.portcullis/src/Portcullis.Engine/Provenance/GitCommandRunner.cs},
\path{letsgolegacy.portcullis/src/Portcullis.Engine/Provenance/AiToolClassifier.cs},
\path{letsgolegacy.portcullis/src/Portcullis.Engine/Provenance/ProvenanceModels.cs},
\path{letsgolegacy.portcullis/src/Portcullis.Cli/Program.cs},
\path{letsgolegacy.portcullis/tests/Portcullis.Engine.Tests/ScannerGateTests.cs}}

\subsection{Baseline znanych findingów i~wyciszenia}\label{sec:portcullis-baseline}

Zakres diffu już zapobiega blokowaniu na długu, którego PR nie dotyka. Baseline (ticket R3)
pokrywa to, czego zakres diffu nie obejmuje:
\begin{itemize}[leftmargin=*]
  \item \textbf{stare findingi na liniach, które „zmieniły się” bez zmiany treści} --
        przeformatowanie bloku, przeniesienie metody w~obrębie pliku, zmiana końców linii
        wciągają stare findingi do diffu; bez baseline blokują, choć PR niczego nie
        wprowadził;
  \item \textbf{bramki na całym drzewie} -- zaplanowany audyt, CI bez użytecznego zakresu
        commitów i~fallback po degradacji; bez baseline blokują na całym długu naraz;
  \item \textbf{SARIF dla code scanning} -- ogranicza się do nowych findingów, zamiast przy
        każdym wysłaniu raportować zaakceptowany dług.
\end{itemize}
Finding zaakceptowany przez baseline jest nadal \textbf{raportowany} -- w~JSON
z~\repopath{baselined: true}, w~komentarzu z~dopiskiem \emph{(accepted by the baseline)} --
ale nie \textbf{blokuje} w~żadnym zakresie (\repopath{all}, \repopath{diff}, fallback po
degradacji). \repopath{gate.blockingErrorCount} plus \repopath{gate.acceptedByBaselineCount}
to wszystkie błędy w~zakresie.

\paragraph{Fingerprint.} Wpis baseline to fingerprint findingu w~schemacie
\repopath{portcullisFingerprint/v1}: szesnastkowy (małe litery) SHA-256 z~pięciu wierszy --
wersji schematu, identyfikatora reguły, ścieżki pliku (względem skanowanego katalogu),
tekstu linii findingu po normalizacji białych znaków (przycięcie i~zwinięcie ciągów do
jednej spacji) oraz numeru wystąpienia. Dwa findingi tej samej reguły na identycznym
tekście w~jednym pliku rozróżnia wystąpienie (0, 1, \ldots{} w~kolejności pliku). Finding
bez lokalizacji zamiast tekstu linii używa komunikatu. Test przypina wartość policzoną
niezależnie od .NET.

\begin{xltabular}{\linewidth}{@{}>{\raggedright\arraybackslash}X >{\raggedright\arraybackslash}p{6.4cm}@{}}
\toprule
\textbf{Zmiana} & \textbf{Fingerprint} \\
\midrule
\endfirsthead
\toprule
\textbf{Zmiana} & \textbf{Fingerprint} \\
\midrule
\endhead
Linie dodane lub usunięte w~innym miejscu pliku & bez zmian \\
Linia przeindentowana, zmienione białe znaki & bez zmian \\
Przeredagowany komunikat reguły w~nowej wersji & bez zmian (poza findingiem całego repozytorium, którego komunikat zastępuje linię) \\
Zmieniony poziom (eskalacja, konfiguracja) & bez zmian \\
Linia przepisana & \textbf{zmienia się} -- finding jest nowy i~wymaga ponownego przeglądu \\
Plik przemianowany lub przeniesiony & \textbf{zmienia się} \\
Skan uruchomiony na innym katalogu & \textbf{zmienia się}: baseline zapisuje się i~czyta z~tą samą ścieżką skanu \\
\bottomrule
\end{xltabular}

Numery linii celowo nie są wejściem: baseline oparty na liniach starzeje się przy pierwszej
edycji powyżej findingu -- każdy zaakceptowany finding poniżej się przesuwa, przestaje
pasować i~wraca jako nowy, więc baseline blokowałby dokładnie te PR, które miał przepuścić.

\paragraph{Plik.} Własny baseline repozytorium (\repopath{portcullis-baseline.json}) zapisano
poleceniem \repopath{--write-baseline}; zawiera dwa findingi \path{PORTCULLIS_MIG_SYNC_OVER_ASYNC}
w~\path{Portcullis.Engine/Provenance/GitCommandRunner.cs} (linie 74 i~77). Pierwszy wpis:

\begin{lstlisting}[escapeinside={(*}{*)}]
{
  "schemaVersion": "1.0.0",
  "fingerprintVersion": "portcullisFingerprint/v1",
  "about": "Findings accepted as pre-existing. Portcullis still reports them, but they do not block and are left out of the filtered SARIF. (*\ldots*)",
  "findings": [
    {
      "fingerprint": "3ded74b18cca77948520679cf45b3fb644383312223aadfe5141c50c9c84cd1c",
      "ruleId": "PORTCULLIS_MIG_SYNC_OVER_ASYNC",
      "filePath": "Portcullis.Engine/Provenance/GitCommandRunner.cs",
      "line": 74,
      "message": "'stderrTask.Result' blocks the calling thread until the task completes (Task<TResult>.Result). (*\ldots*)"
    },
    (*\ldots*)
  ]
}
\end{lstlisting}

Dopasowywany jest wyłącznie \repopath{fingerprint}; reguła, plik, linia i~komunikat są dla
recenzenta PR, który powiększa baseline (widzi, co jest akceptowane, a~nie kolumnę
skrótów). \repopath{line} to miejsce findingu w~chwili zapisu i~nie jest później
aktualizowane. Wpisy są sortowane po pliku, linii i~regule, z~końcami linii \repopath{\textbackslash{}n}
na każdej platformie, więc przepisany baseline różni się tylko tam, gdzie zmieniły się
findingi.

\paragraph{Zachowanie CLI.}
\begin{itemize}[leftmargin=*]
  \item \repopath{--write-baseline} zapisuje wszystkie findingi skanu do pliku
        \repopath{--baseline}, zastępując go, i~ocenia ten przebieg względem nowego pliku --
        więc przebieg przechodzi, a~następne blokują tylko na nowych. Bez
        \repopath{--baseline} to błąd użycia, nie flaga, która nic nie robi.
  \item Pliku, którego nie da się użyć, CLI nie traktuje jak pustego baseline: brak pliku,
        błąd odczytu, niepoprawny JSON, JSON \repopath{null}, \repopath{schemaVersion}
        inna niż 1.x, inny \repopath{fingerprintVersion}, brak tablicy
        \repopath{findings} albo wpis bez fingerprintu kończą przebieg kodem 2
        z~komunikatem nazywającym plik i~poprawkę. Odczyt jako „nic nie akceptuje” byłby
        bezpieczny dla bramki, ale zamieniłby literówkę w~workflow w~bramkę blokującą na
        całym długu, bez słowa wyjaśnienia.
  \item Wpisy, które niczego nie dopasowały (dług spłacony albo linia przepisana), są
        liczone: \repopath{baseline.entryCount} minus \repopath{baseline.acceptedCount};
        CLI wypisuje notę na stderr, SARIF ma notyfikację \repopath{note}. Nigdy niczego nie
        blokują; przepisanie baseline je usuwa.
\end{itemize}

Ograniczenia: baseline jest związany ze skanowanym katalogiem (skan \repopath{src} w~jednym
workflow i~katalogu głównego w~drugim wymaga dwóch plików); identyczne linie są
rozróżniane tylko kolejnością, więc przy dodaniu kopii zaakceptowanej linii powyżej
oryginału nowy finding zostanie zgłoszony poprawnie co do liczby, ale na drugiej
z~identycznych linii.

\paragraph{Wyciszenia.} Źródła Portcullis nie opisują własnego mechanizmu wyciszania
pojedynczych findingów w~kodzie, a~SARIF Portcullis nie zawiera obiektów
\repopath{suppressions}. Rolę wyciszeń pełnią:
\begin{itemize}[leftmargin=*]
  \item baseline -- akceptacja konkretnych findingów, widoczna w~PR, który go zmienia;
  \item \repopath{.editorconfig}: \repopath{dotnet\_diagnostic.<ID>.severity = none} albo
        \repopath{suggestion} -- działa w~kanale pakietu analyzerów, per reguła i~per
        ścieżka; skan CLI (akcja, kontener) tych nadpisań nie stosuje;
  \item \repopath{migrationExemptFolders} i~\repopath{systemWebAllowedTypes} dla reguł
        migracyjnych; jawnie pusta lista konwencji dla reguł zasad („ten projekt nie ma
        wspólnego jądra”);
  \item \repopath{fail-on-block: false} w~akcji GitHub -- raport bez egzekwowania.
\end{itemize}
Konwencję wyciszania z~uzasadnieniem definiuje pakiet standardów C4
(sekcja~\ref{sec:standardy}):
\begin{lstlisting}
#pragma warning disable <DIAGNOSTIC> // SECONDKEY-FLAG <rule-id>: <why>
\end{lstlisting}
Źródła Portcullis nie opisują, jak skan CLI traktuje takie dyrektywy.

\zrodla{\path{letsgolegacy.portcullis/docs/SARIF.md}, \path{letsgolegacy.portcullis/docs/TUTORIAL.md},
\path{letsgolegacy.portcullis/docs/CONFIGURATION.md}, \path{letsgolegacy.portcullis/portcullis-baseline.json},
\path{letsgolegacy.portcullis/src/Portcullis.Engine/Findings/FindingFingerprint.cs},
\path{letsgolegacy.portcullis/src/Portcullis.Engine/Findings/Baseline.cs},
\path{letsgolegacy.portcullis/src/Portcullis.Engine/Findings/BaselineFile.cs},
\path{letsgolegacy.portcullis/src/Portcullis.Engine/Sarif/SarifReport.cs},
\path{letsgolegacy.portcullis/src/Portcullis.Cli/Program.cs},
\path{letsgolegacy.portcullis/src/Portcullis.Cli/ScanCommandOptions.cs},
\path{letsgolegacy.portcullis/action.yml},
\path{letsgolegacy.secondkey-standards/catalog/skill.template.md}}

\subsection{Konfiguracja}\label{sec:portcullis-konfiguracja}

Każda reguła zasad rozpoznaje jądro, warstwę domeny, granicę adaptera i~punkt wejścia po
nazwach folderów i~plików. Dopóki nie było to konfigurowalne, zespół z~jądrem
\repopath{Common/} albo adapterami w~\repopath{External/} dostawał zielony skan, który
niczego nie sprawdził, a~jedyne pokrętło -- poziom w~\repopath{.editorconfig} -- potrafi
regułę wyłączyć, ale nie powie jej, gdzie patrzeć. Reguły migracyjne nie potrzebują
konwencji; ich dwa ustawienia to wyjątki.

\paragraph{Dwa kanały, jeden model.} Analyzer jest ładowany przez kompilator, może działać
poza procesem w~IDE i~przy \repopath{EnforceExtendedAnalyzerRules} nie może używać
\repopath{System.IO} (błąd \repopath{RS1035}), więc konfigurację dostaje przez Roslyn.
Model konwencji (\repopath{PortcullisConventions}) i~nazwy kluczy
(\repopath{PortcullisConventionKeys}) są jedne:

\begin{xltabular}{\linewidth}{@{}>{\raggedright\arraybackslash}p{4.4cm} >{\raggedright\arraybackslash}p{4.4cm} >{\raggedright\arraybackslash}X@{}}
\toprule
\textbf{Kanał} & \textbf{Zespół pisze} & \textbf{Droga do reguł} \\
\midrule
\endfirsthead
\toprule
\textbf{Kanał} & \textbf{Zespół pisze} & \textbf{Droga do reguł} \\
\midrule
\endhead
\texttt{Portcullis.Analyzers} (NuGet, w~buildzie) & \texttt{.globalconfig} albo \texttt{.editorconfig} z~\texttt{is\_global = true} & \texttt{AnalyzerConfigOptions.GlobalOptions} Roslyn \\
CLI \texttt{portcullis}, akcja, kontener & \texttt{portcullis.json} w~katalogu skanu & \texttt{PortcullisConfigFile} spłaszcza JSON na te same klucze \\
\bottomrule
\end{xltabular}

Wcześniej CLI przekazywał \repopath{options: null} do \repopath{CompilationWithAnalyzers},
co odrzucało całą konfigurację analyzerów: pakiet dało się skonfigurować, CLI nie -- te
same reguły, dwa zachowania. To poprawiono.

\paragraph{\texttt{portcullis.json}.} Plik leży w~katalogu, który się skanuje (ta sama
ścieżka co w~\repopath{portcullis scan}); nie jest szukany w~górę drzewa, bo niejawna
konfiguracja z~katalogu nadrzędnego utrudnia odpowiedź na pytanie „czemu ta reguła nie
strzela”. Każda właściwość jest opcjonalna; wartość domyślna obowiązuje per klucz. Nazwy
właściwości są porównywane bez wielkości liter, komentarze i~przecinki końcowe są
dozwolone, a~nieznane właściwości ignorowane (konfiguracja pisana dla nowszej wersji działa
ze starszą).

\begin{lstlisting}
{
  "kernelFolders": ["Common", "Shared"],
  "entityFolders": ["Model"],
  "adapterFolders": ["External", "Gateways"],
  "vendorNamespaces": ["Stripe", "Google.Cloud", "Acme.Payments"],
  "entryPointFileNames": ["Program.cs", "Startup.cs"],
  "kernelLineCeiling": 500,
  "migrationExemptFolders": ["Legacy"],
  "systemWebAllowedTypes": ["System.Web.HttpUtility", "System.Web.IHtmlString"]
}
\end{lstlisting}

\paragraph{\texttt{.globalconfig}.} Te same wartości jako płaskie klucze (pełna lista
w~tabeli poniżej; przykład z~dokumentacji ustawia wszystkie osiem). To opcje
\emph{globalne}: wartość w~zwykłej sekcji \repopath{[*.cs]} jest opcją per plik i~nie
zostanie odczytana. Listy rozdziela przecinek, białe znaki wokół są przycinane.

\begin{lstlisting}
is_global = true

portcullis_kernel_folders = Common, Shared
portcullis_kernel_line_ceiling = 500
portcullis_migration_exempt_folders = Legacy
\end{lstlisting}

\begingroup\small
\begin{xltabular}{\linewidth}{@{}>{\raggedright\arraybackslash}p{5.0cm} >{\raggedright\arraybackslash}X >{\raggedright\arraybackslash}p{4.2cm}@{}}
\toprule
\textbf{Klucz / właściwość JSON} & \textbf{Domyślnie} & \textbf{Używa} \\
\midrule
\endfirsthead
\toprule
\textbf{Klucz / właściwość JSON} & \textbf{Domyślnie} & \textbf{Używa} \\
\midrule
\endhead
\path{portcullis_kernel_folders} \newline \path{kernelFolders} & \texttt{ServiceDefaults}, \texttt{Kernel}, \texttt{SharedKernel} & obie reguły P2 \\
\path{portcullis_entity_folders} \newline \path{entityFolders} & \texttt{Domain}, \texttt{Entities} & \path{PORTCULLIS_P2_KERNEL_ENTITY_REFERENCE} \\
\path{portcullis_adapter_folders} \newline \path{adapterFolders} & \texttt{Infrastructure}, \texttt{Adapters}, \texttt{Integrations} & wyjątek adaptera w~P11 \\
\path{portcullis_vendor_namespaces} \newline \path{vendorNamespaces} & \texttt{Anthropic}, \texttt{OpenAI}, \texttt{Stripe}, \texttt{Twilio}, \texttt{SendGrid}, \texttt{PayPal}, \texttt{Amazon}, \texttt{Azure}, \texttt{Google.Cloud}, \texttt{Firebase}, \texttt{MailKit} & P11 \\
\path{portcullis_entry_point_file_names} \newline \path{entryPointFileNames} & \texttt{Program.cs} & P15 i~wyjątek composition root w~P11 \\
\path{portcullis_kernel_line_ceiling} \newline \path{kernelLineCeiling} & \texttt{800} & \path{PORTCULLIS_P2_KERNEL_LOC_CEILING} \\
\path{portcullis_migration_exempt_folders} \newline \path{migrationExemptFolders} & brak & cztery reguły \texttt{PORTCULLIS\_MIG\_*} \\
\path{portcullis_system_web_allowed_types} \newline \path{systemWebAllowedTypes} & \texttt{System.Web.HttpUtility}, \texttt{System.Web.IHtmlString} & \path{PORTCULLIS_MIG_SYSTEM_WEB} \\
\path{portcullis_path_root} & brak (CLI ustawia katalog skanu) & wszystkie konwencje; nie liczy się jako „skonfigurowane” \\
\bottomrule
\end{xltabular}
\endgroup

\paragraph{Reguły dopasowania.}
\begin{itemize}[leftmargin=*]
  \item \textbf{Foldery}: segment ścieżki równy nazwie (\repopath{Common}) albo zakończony
        na \repopath{.Nazwa} (\repopath{Acme.Common}), bez wielkości liter.
  \item \textbf{Przestrzenie nazw dostawców}: dopasowanie na granicy segmentu kropki --
        \repopath{Azure} łapie \repopath{using Azure;} i~\repopath{using Azure.Storage.Blobs;},
        nie łapie \repopath{using AzureFoo;}; wpisy wielosegmentowe
        (\repopath{Google.Cloud}) działają.
  \item \textbf{Ścieżki względem katalogu skanu}: checkout w~\repopath{\textasciitilde{}/Domain/}
        nie czyni wszystkich plików „encjami”, a~ścieżka workspace CI z~\repopath{Integrations}
        nie zwalnia całego drzewa z~reguły dostawców. Pakiet analyzerów nie zna katalogu
        skanu (kompilator podaje ścieżki bezwzględne), więc dopasowuje pełną ścieżkę, chyba
        że ustawi się \repopath{portcullis\_path\_root}.
  \item \textbf{Konfiguracja zastępuje, nie rozszerza}: lista z~jednym elementem
        \repopath{Common} oznacza tylko \repopath{Common}; model addytywny uniemożliwiłby
        zawężenie konwencji. Jawnie pusta lista (w~JSON \repopath{kernelFolders} równe
        \repopath{[]}) to deklaracja „tego projektu to nie dotyczy”: wyłącza reguły bez
        tłumienia ich diagnostyk, a~pominięcie klucza przywraca wartości domyślne.
  \item \textbf{Sufit linii}: wartość niebędąca dodatnią liczbą całkowitą przechodzi na
        domyślną (analyzer rzucający wyjątek dawałby \repopath{AD0001} i~wyłączał pozostałe
        reguły) i~nie liczy się jako „skonfigurowana”.
\end{itemize}

\paragraph{Poziom reguł.} W~kanale pakietu analyzerów poziom zmienia się zwykłą drogą
Roslyn (\repopath{dotnet\_diagnostic.<ID>.severity}), także dla diagnostyk pokrycia.
\textbf{W~kanale CLI, akcji i~kontenera to nie działa}: \repopath{portcullis scan} buduje
własną kompilację i~raportuje każdą diagnostykę na jej \repopath{defaultSeverity}; nie da
się z~CLI ani podnieść reguły do blokującej, ani jej wyciszyć.

\paragraph{Co zostało użyte.} Wynik skanu ma obiekt \repopath{configuration}
(\repopath{path} -- \repopath{null}, gdy pliku nie ma; \repopath{error} -- gdy plik jest, ale
nie dał się sparsować). Zepsuty \repopath{portcullis.json} nie wyłącza bramki: skan idzie
dalej na wartościach domyślnych, ale CLI wypisuje ostrzeżenie, a~SARIF ma notyfikację --
zespół pracujący miesiącami na domyślnych przy przekonaniu, że jego konfiguracja działa,
to ta sama pusta zieleń w~innym przebraniu.

\paragraph{Czego nie da się skonfigurować.} Kształtu reguły (można wskazać adaptery, nie
można przedefiniować, czym jest wyciek typu dostawcy), zestawu reguł
(\repopath{RuleRegistry.All} jest stały) ani poziomu w~CLI.

\zrodla{\path{letsgolegacy.portcullis/docs/CONFIGURATION.md}, \path{letsgolegacy.portcullis/docs/rules/MIGRATION.md},
\path{letsgolegacy.portcullis/src/Portcullis.Rules/PortcullisConventions.cs},
\path{letsgolegacy.portcullis/src/Portcullis.Rules/PathConventions.cs},
\path{letsgolegacy.portcullis/src/Portcullis.Engine/Configuration/PortcullisConfigFile.cs},
\path{letsgolegacy.portcullis/src/Portcullis.Rules/PACKAGE-README.md}}

\subsection{SARIF 2.1.0}\label{sec:portcullis-sarif}

Ticket R3: „SARIF 2.1 z~filtrem zmienionych linii i~plikiem baseline”, zamknięty, gdy „pull
request produkuje przefiltrowany SARIF”. \repopath{--sarif <plik>} zapisuje log SARIF
2.1.0 (standard OASIS) obok zwykłego JSON na stdout; stdout i~kod wyjścia się nie zmieniają.

\begin{lstlisting}
portcullis scan ./src \
  --provenance-range "$BASE_SHA..$HEAD_SHA" \
  --provenance-repo . \
  --baseline portcullis-baseline.json \
  --sarif portcullis.sarif
\end{lstlisting}

\begingroup\small
\begin{xltabular}{\linewidth}{@{}>{\raggedright\arraybackslash}p{3.6cm} >{\raggedright\arraybackslash}X@{}}
\toprule
\textbf{Część logu} & \textbf{Co zapisuje Portcullis} \\
\midrule
\endfirsthead
\toprule
\textbf{Część logu} & \textbf{Co zapisuje Portcullis} \\
\midrule
\endhead
\texttt{\$schema}, \texttt{version} & URI schematu OASIS 2.1.0 (errata 01) i~\texttt{2.1.0}; jeden run. \\
\texttt{tool.driver} & \texttt{name} \texttt{Portcullis}, \texttt{version} (wersja informacyjna silnika z~commitem), \texttt{informationUri}, \texttt{rules}: każda zarejestrowana reguła (także bez wyników), posortowana po identyfikatorze; tytuł jako opis krótki i~pełny, \texttt{help} i~\texttt{helpUri} do dokumentacji rodziny reguł (\texttt{Migration} → \texttt{docs/rules/MIGRATION.md}, \texttt{Meta} → \texttt{docs/CONFIGURATION.md}, reguły zasad → README), \texttt{defaultConfiguration.level}, kategoria jako \texttt{properties.category} i~tag. Wynik reguły spoza rejestru dostaje deskryptor z~kategorią \texttt{Other}. \\
\texttt{invocations[0]} & \texttt{executionSuccessful: true} i~\texttt{toolExecutionNotifications}: \texttt{warning}, gdy nie dało się zastosować zakresu diffu, \texttt{warning}, gdy nie dało się odczytać pliku konfiguracji, \texttt{note}, gdy wpisy baseline niczego nie dopasowały. \\
\texttt{originalUriBaseIds} & \texttt{\%SRCROOT\%}: katalog główny repozytorium jako bezwzględny URI \texttt{file:}. \\
\texttt{results} & Jeden wynik na finding, który liczy się przeciw zmianie. \\
\texttt{properties} & \texttt{scope} (\texttt{diff} albo \texttt{all}), \texttt{findingsInScan}, \texttt{excludedOutsideChangedLines}, \texttt{excludedByBaseline} -- to, co odrzuciły filtry, jest policzone w~samym logu. \\
\bottomrule
\end{xltabular}
\endgroup

Każdy wynik ma pola:
\begin{itemize}[leftmargin=*]
  \item \repopath{ruleId} i~\repopath{ruleIndex} -- reguła i~jej indeks w~tablicy
        \repopath{tool.driver.rules};
  \item \repopath{level} -- poziom po eskalacji; \repopath{error} i~\repopath{warning}
        przechodzą bez zmian, \repopath{info} staje się \repopath{note};
  \item \repopath{message.text} -- komunikat naruszenia;
  \item \path{locations[0].physicalLocation} -- URI pliku względem katalogu głównego
        repozytorium (bazą jest \repopath{\%SRCROOT\%}, każdy segment ścieżki jest
        zakodowany procentowo) i~linia początkowa w~\repopath{region.startLine};
  \item \repopath{partialFingerprints} -- fingerprint pod kluczem
        \repopath{portcullisFingerprint/v1};
  \item \repopath{baselineState} o~wartości \repopath{new} -- tylko wtedy, gdy skan dostał
        baseline (wszystko, co baseline zaakceptował, zostało pominięte).
\end{itemize}

\paragraph{Lokalizacje względem katalogu repozytorium.} Konsument SARIF rozwiązuje ścieżki
względem swojego checkoutu: GitHub code scanning szuka \repopath{src/Shop/Orders.cs}, nie
\repopath{Shop/Orders.cs}. Katalogiem głównym jest \repopath{--provenance-repo}, a~bez niego
skanowana ścieżka -- przy skanowaniu podkatalogu trzeba więc podać
\repopath{--provenance-repo}. Finding całego repozytorium (diagnostyki pokrycia) nie ma
pliku i~jest zapisywany bez \repopath{locations}, co SARIF dopuszcza; nigdy nie leży na
zmienionej linii, więc trafia tylko do logu bez zakresu diffu, a~GitHub code scanning
takiego wyniku nie wyświetla.

\paragraph{Które findingi trafiają do logu.} Te, które liczą się przeciw zmianie: definicja
zakresu bramki zastosowana do wszystkich poziomów, nie tylko błędów. Przy
\repopath{--provenance-range} zostają tylko findingi na zmienionych liniach (ta sama
\repopath{ChangedLines}, którą pyta bramka); przy \repopath{--baseline} odpadają
zaakceptowane. JSON na stdout nadal zawiera wszystkie findingi, z~\repopath{baselined: true}
przy zaakceptowanych. Gdy zakresu nie da się policzyć, SARIF obejmuje cały skan i~podaje
powód jako notyfikację \repopath{warning} -- filtr, który spadłby do „nic się nie zmieniło”,
opublikowałby pusty log dla PR, którego zmian nikt nie obejrzał.

\paragraph{Przykład uruchomiony.} Repozytorium robocze z~dwoma istniejącymi findingami
zaakceptowanymi w~baseline i~PR, który przeindentowuje jeden z~nich i~dodaje nowy błąd
(2026-09-26, CLI z~gałęzi ticketu):

\begin{lstlisting}
# commit A: Settings.cs reads ConfigurationManager (an error), Orders.cs blocks on .Result (a warning)
$ portcullis scan src --baseline portcullis-baseline.json --write-baseline
portcullis: wrote 3 finding(s) to the baseline 'portcullis-baseline.json'; this run is judged against it.

# commit B, the pull request: re-indents Settings.cs line 7, adds Pricing.cs with a new error
$ portcullis scan src --provenance-range A..B --provenance-repo . \
    --baseline portcullis-baseline.json --sarif pr.sarif
exit 1
summary:  { errorCount: 2, warningCount: 2, infoCount: 0 }
gate:     { blocked: true, scope: "diff", blockingErrorCount: 1, acceptedByBaselineCount: 1 }
baseline: { path: "portcullis-baseline.json", entryCount: 3, acceptedCount: 3 }
\end{lstlisting}

Trzeci wpis baseline to ostrzeżenie pokrycia, które dostaje repozytorium bez konwencji.
Przeindentowana linia jest w~diffie i~jej błąd nadal istnieje, ale baseline go rozpoznaje,
więc blokuje tylko nowy błąd; ten sam zakres bez \repopath{--baseline} daje
\repopath{blockingErrorCount: 2}. SARIF zawiera dokładnie nowy finding
(\repopath{ruleId} \path{PORTCULLIS_MIG_CONFIGURATION_MANAGER}, \repopath{ruleIndex} 1,
\repopath{src/Shop/Pricing.cs}, linia 7, \repopath{baselineState} \repopath{new}),
a~\repopath{properties} to \repopath{scope} \repopath{diff}, \repopath{findingsInScan} 4,
\repopath{excludedOutsideChangedLines} 2, \repopath{excludedByBaseline} 1. Komentarz na PR
mówi: „Blocking this PR — 1 error within this PR's own changed lines. 1 more there is
accepted by the baseline.”

\paragraph{Weryfikacja.} Każdy log produkowany w~testach jest walidowany względem
oficjalnego schematu OASIS SARIF 2.1.0 osadzonego w~zestawie testów (plik skopiowany bez
zmian z~\repopath{oasis-tcs/sarif-spec}, commit \repopath{adbb670}; test przypina jego SHA-256)
biblioteką JsonSchema.Net z~włączonymi asercjami formatu. Schemat deklaruje draft-04, którego
biblioteka nie implementuje, więc jest oceniany jako draft-07; test sprawdza, że używa tylko
słów kluczowych o~tym samym znaczeniu w~obu wersjach i~tylko lokalnych referencji. Kontrole
negatywne pokazują, że walidator odrzuca złą wersję, wynik bez komunikatu, linię 0,
narzędzie bez drivera, niezakodowaną spację w~URI i~nieznany poziom. \repopath{SarifReportTests}
sprawdza treść (pola wyników, oba filtry i~ich liczniki, fallback po degradacji, notę
o~nieaktualnym baseline, nieczytelną konfigurację, kodowanie spacji, \repopath{\#} i~litery
spoza ASCII w~ścieżkach, mapowanie poziomów, deterministyczność), a~testy procesu CLI --
m.in. SARIF z~prawdziwej historii dwóch commitów oraz kod 2 przy braku pliku baseline
i~przy ścieżce SARIF, której nie da się zapisać.

\paragraph{Ograniczenia.} Region ma tylko linię początkową (naruszenie nie zapisuje kolumny
ani końca). Baseline jest związany ze skanowanym katalogiem. Akcja GitHub nie udostępnia
jeszcze \repopath{--sarif} ani \repopath{--baseline} -- instaluje narzędzia z~NuGet, gdzie
nic nie jest opublikowane, więc zmiany nie dało się w~niej przetestować; workflow repozytorium
wywołuje CLI bezpośrednio (sekcja~\ref{sec:portcullis-ci}). To, które alerty GitHub pokaże
na PR i~jak je śledzi między przebiegami, jest poza kontrolą Portcullis: jego część kończy
się na poprawnym, przefiltrowanym pliku.

\zrodla{\path{letsgolegacy.portcullis/docs/SARIF.md},
\path{letsgolegacy.portcullis/src/Portcullis.Engine/Sarif/SarifReport.cs},
\path{letsgolegacy.portcullis/src/Portcullis.Cli/Program.cs},
\path{letsgolegacy.portcullis/tests/Portcullis.Engine.Tests/Sarif/THIRD-PARTY-NOTICE.md},
\path{letsgolegacy.portcullis/tests/Portcullis.Engine.Tests/Sarif/SarifSchemaTests.cs},
\path{letsgolegacy.portcullis/tests/Portcullis.Engine.Tests/Sarif/SarifReportTests.cs},
\path{letsgolegacy.portcullis/tests/Portcullis.Cli.Tests/SarifAndBaselineEndToEndTests.cs}}

\subsection{Komentarz na pull request}\label{sec:portcullis-komentarz}

\repopath{portcullis-ci-comment} (pakiet \repopath{Portcullis.CiComment}) zamienia wynik
skanu w~komentarz Markdown, aktualizuje w~miejscu \textbf{jeden} komentarz sticky zamiast
dodawać nowy przy każdym pushu i~kończy się kodem niezerowym, gdy bramka blokuje. To osobne
narzędzie, bo rozmawia z~API GitHub; dzięki temu \repopath{portcullis scan} nie ma powodu
wymagać \repopath{GITHUB\_TOKEN}.

\begin{lstlisting}
portcullis-ci-comment --current <path|-> [--previous <path|->] [--output <path>] [--pr <number>]
\end{lstlisting}

\repopath{-} oznacza stdin. Nieznany argument, brak \repopath{--current} albo
nieliczbowe \repopath{--pr} to błąd użycia (kod 2). Wynik skanu musi mieć
\repopath{violations} i~\repopath{summary}; nieczytelny albo niekompletny
\repopath{--current} kończy się kodem 2, a~nieczytelny \repopath{--previous} -- ostrzeżeniem
i~komentarzem bez diffu (na pierwszym pushu PR pamięci poprzedniego skanu nie ma, więc to
zwykła ścieżka).

\paragraph{Treść komentarza.}
\begin{itemize}[leftmargin=*]
  \item Ukryty znacznik \repopath{<!-- portcullis:pr-comment -->} i~nagłówek
        \repopath{\#\# Portcullis}.
  \item Linia główna: liczby błędów, ostrzeżeń i~info oraz „N rules evaluated across M files
        in Xms”; osobne warianty dla zera reguł i~dla braku naruszeń.
  \item Adnotacja bramki: przy degradacji -- że nie dało się zawęzić bramki do zmian PR,
        z~powodem; w~zakresie \repopath{diff} -- czy bramka blokuje („Blocking this PR —
        N errors within this PR's own changed lines.”) albo dlaczego błędy widoczne w~liczbach
        nie blokują (poza zmienionymi liniami, zaakceptowane przez baseline); w~zakresie
        \repopath{all} tylko wtedy, gdy baseline zaakceptował jakiś błąd.
  \item Przy \repopath{--previous}: „N new · M resolved since the last push” i~listy
        nowych oraz rozwiązanych.
  \item Findingi pogrupowane po pliku (nagłówek z~ścieżką, „line N — reguła: komunikat”),
        findingi bez pliku pod nagłówkiem „Repository-wide”; zaakceptowane przez baseline
        z~dopiskiem \emph{(accepted by the baseline)}; stopka z~czasem skanu.
\end{itemize}
Pochodzenie AI nigdy nie pojawia się w~komentarzu: formatter nie widzi tego sygnału, widzi
tylko jego skutek w~poziomie naruszenia.

\paragraph{Diff między pushami.} Tożsamość naruszenia to pięć pól kontraktu:
\repopath{ruleId}, \repopath{filePath}, \repopath{line}, \repopath{message},
\repopath{severity}. Naruszenie, któremu zmienił się komunikat albo poziom, liczy się jako
jedno rozwiązane i~jedno nowe -- tak w~M4 pięć eskalowanych ostrzeżeń P10 pojawiło się jako
„rozwiązane” na poziomie \repopath{warning} i~„nowe” na \repopath{error}.
\repopath{fingerprint} i~\repopath{baselined} celowo nie wchodzą do tożsamości: poprzedni
skan pochodzi z~pamięci podręcznej, być może z~silnika sprzed fingerprintów, a~baseline
może się zmienić między pushami.

\paragraph{Publikacja.} Kontekst pochodzi ze zmiennych środowiskowych GitHub Actions.
Token i~repozytorium to \repopath{GITHUB\_TOKEN} oraz \repopath{GITHUB\_REPOSITORY}. Numer
PR pochodzi z~opcji \repopath{--pr} albo z~pola \repopath{pull\_request.number} w~pliku
zdarzenia wskazanym przez \repopath{GITHUB\_EVENT\_PATH}. Adres API to
\repopath{GITHUB\_API\_URL}, a~bez tej zmiennej \repopath{https://api.github.com/}. Bez pełnego
kontekstu narzędzie tylko renderuje na stdout i~pisze „skipping comment publish” -- to tryb
lokalnego suchego przebiegu. Publikacja używa REST API komentarzy issue: lista komentarzy
PR (\repopath{per\_page=100}), wyszukanie znacznika, \repopath{PATCH} istniejącego albo
\repopath{POST} nowego. Porażka publikacji (np. token tylko do odczytu w~PR z~forka, 403)
daje ostrzeżenie i~nie zmienia werdyktu; propaguje się tylko anulowanie. Kod wyjścia:
1, gdy bramka blokuje, 0 w~przeciwnym razie, 2 przy błędzie użycia, nieczytelnym
\repopath{--current} albo błędzie renderowania.

\zrodla{\path{letsgolegacy.portcullis/src/Portcullis.CiComment/PACKAGE-README.md},
\path{letsgolegacy.portcullis/src/Portcullis.CiComment/Program.cs},
\path{letsgolegacy.portcullis/src/Portcullis.CiComment/CliOptions.cs},
\path{letsgolegacy.portcullis/src/Portcullis.CiComment/CommentFormatter.cs},
\path{letsgolegacy.portcullis/src/Portcullis.CiComment/ViolationDiff.cs},
\path{letsgolegacy.portcullis/src/Portcullis.CiComment/StickyComment.cs},
\path{letsgolegacy.portcullis/src/Portcullis.CiComment/GitHubPublishContext.cs},
\path{letsgolegacy.portcullis/src/Portcullis.CiComment/GitHubPullRequestCommentClient.cs},
\path{letsgolegacy.portcullis/src/Portcullis.CiComment/Portcullis.CiComment.csproj},
\path{letsgolegacy.portcullis/docs/M4-INTEGRATION.md}}

\subsection{Dystrybucja i~wersjonowanie}\label{sec:portcullis-dystrybucja}

\begin{xltabular}{\linewidth}{@{}>{\raggedright\arraybackslash}p{2.9cm} >{\raggedright\arraybackslash}p{5.6cm} >{\raggedright\arraybackslash}X@{}}
\toprule
\textbf{Kanał} & \textbf{Artefakt} & \textbf{Dla kogo} \\
\midrule
\endfirsthead
\toprule
\textbf{Kanał} & \textbf{Artefakt} & \textbf{Dla kogo} \\
\midrule
\endhead
NuGet (analyzer) & \texttt{Portcullis.Analyzers} & Domyślny: reguły jako diagnostyki builda w~każdym IDE i~CI. \\
NuGet (narzędzia) & \texttt{Portcullis.Cli} \textrightarrow{} \texttt{portcullis}, \texttt{Portcullis.CiComment} \textrightarrow{} \texttt{portcullis-ci-comment} & Wynik maszynowy, bramka zawężona do diffu, komentarze PR. \\
GitHub Action & \path{konradcinkusz/letsgolegacy.portcullis@v0} (composite) & Bramka PR podłączona od końca do końca. \\
Kontener GHCR & \path{ghcr.io/konradcinkusz/letsgolegacy.portcullis} & CI inne niż GitHub Actions, bez instalacji SDK .NET. \\
\bottomrule
\end{xltabular}

\paragraph{Pakiet analyzerów jako kanał główny.} To jedyny kanał bez żadnej instalacji
w~CI: jedno \repopath{PackageReference} i~naruszenia stają się zwykłymi ostrzeżeniami
i~błędami kompilatora -- w~Visual Studio, Riderze, \repopath{dotnet build}, dowolnym CI i~na
maszynie programisty przed pushem. Dlatego reguły mają osobny projekt
\repopath{netstandard2.0}: analyzer jest ładowany przez proces kompilatora, więc musi się
załadować zarówno w~MSBuild .NET Framework (Visual Studio), jak i~w~\repopath{dotnet build}.
Analyzer zbudowany na \repopath{net8.0} czy \repopath{net10.0} nie zgłasza błędu -- po cichu
nie działa w~Visual Studio.

\begingroup\small
\begin{xltabular}{\linewidth}{@{}>{\raggedright\arraybackslash}p{4.0cm} >{\raggedright\arraybackslash}X@{}}
\toprule
\textbf{Pakiet} & \textbf{Budowa} \\
\midrule
\endfirsthead
\toprule
\textbf{Pakiet} & \textbf{Budowa} \\
\midrule
\endhead
\texttt{Portcullis.Analyzers} (\path{src/Portcullis.Rules}) & \texttt{netstandard2.0}. \texttt{IncludeBuildOutput=false}, a~DLL pakowany do \texttt{analyzers/dotnet/cs} (nie do \texttt{lib/} -- tam analyzer instaluje się bez błędu i~nic nie robi). \texttt{DevelopmentDependency=true}, \texttt{SuppressDependenciesWhenPacking=true} (kompilator sam dostarcza Microsoft.CodeAnalysis). Bez pakietu symboli: \texttt{IncludeSymbols=false}, \texttt{DebugType=embedded}, \texttt{EmbedUntrackedSources=true}. README pakietu z~\texttt{PACKAGE-README.md}. \\
\texttt{Portcullis.Cli} & \texttt{net10.0}, \texttt{PackAsTool}, polecenie \texttt{portcullis}; pakiet symboli \texttt{snupkg}. \\
\texttt{Portcullis.CiComment} & \texttt{net10.0}, \texttt{PackAsTool}, polecenie \texttt{portcullis-ci-comment}; pakiet symboli \texttt{snupkg}. \\
\bottomrule
\end{xltabular}
\endgroup

\paragraph{Akcja composite, nie kontenerowa.} Akcje kontenerowe Docker działają tylko na
hostowanych runnerach Linux -- nie na Windows, nie na macOS i~nie na odchudzonych obrazach
Ubuntu bez Dockera. Akcja composite, która instaluje narzędzie .NET, działa tak samo na
wszystkich. Źródła wskazują, że \repopath{trivy-action} odszedł od Dockera do composite,
\repopath{golangci-lint-action} nigdy go nie używał, a~oba bezpośrednie odpowiedniki .NET są
akcjami JavaScript pobierającymi binarkę. Obraz kontenera jest publikowany mimo to: dla
GitLab CI, Jenkinsa, Woodpeckera czy \repopath{docker run} w~Makefile obraz z~git
i~narzędziem to najmniejsze tarcie. Native AOT odrzucono jako dziś niewykonalne (obsługa
trimmingu w~Roslyn, \repopath{dotnet/roslyn\#68154}, ładowanie analyzerów pod AOT,
\repopath{dotnet/runtime\#83355}).

\paragraph{Dwa problemy znalezione przez konsumpcję pakietów.} Oba wyszły dopiero przy
instalacji spakowanych artefaktów w~jednorazowym projekcie. (1)~Analyzer zbudowany na
\repopath{Microsoft.CodeAnalysis.CSharp} 4.11.0 dawał w~projekcie na SDK .NET 8
\repopath{warning CS9057} (zestaw referencjonuje kompilator 4.11.0.0, nowszy niż działający
4.8.0.0) -- stąd polityka „analyzer floor” (sekcja~\ref{sec:portcullis-roslyn}).
(2)~\repopath{IncludeSymbols=true} z~\repopath{IncludeBuildOutput=false} daje pusty
\repopath{.snupkg}, na którym \repopath{dotnet pack} kończy się \repopath{NU5017};
rozwiązaniem jest osadzenie symboli w~zestawie.

\paragraph{Zweryfikowane i~niezweryfikowane.} 2026-08-18 na spakowanych \repopath{.nupkg}:
analyzer z~samego \repopath{PackageReference} zatrzymał build projektu, który nie zna
repozytorium (\path{PORTCULLIS_P9_CONTROLLER_NO_DBCONTEXT} jako błąd, „Build
FAILED”); oba narzędzia zainstalowane przez \repopath{dotnet tool install -g} przeszły cały
łańcuch na jednorazowym repozytorium (błąd poza diffem raportowany, bramka nie blokuje,
kod 0). Obraz GHCR i~workflow wydania nie zostały wykonane -- wymagają prawdziwego tagu
i~poświadczeń rejestru.

\paragraph{Obraz kontenera.} Wieloetapowy \repopath{Dockerfile}: build na
\path{mcr.microsoft.com/dotnet/sdk:10.0} (restore tylko z~plików projektów, żeby warstwa
była cache'owana; \repopath{ARG VERSION=0.1.0}; \repopath{dotnet publish} CLI i~CiComment
do \repopath{/app}), obraz końcowy \path{mcr.microsoft.com/dotnet/runtime:10.0}. Git jest
instalowany, bo to realna zależność runtime'u (\repopath{--provenance-range} wywołuje
\repopath{git diff} i~\repopath{git blame}); \repopath{safe.directory '*'} pozwala
pracować na zamontowanym repozytorium. Użytkownik bez uprawnień root
(\repopath{portcullis}, UID 1001), katalog roboczy \repopath{/workspace}:
\begin{lstlisting}
ENTRYPOINT ["dotnet", "/app/Portcullis.Cli.dll"]
CMD ["scan", "/workspace"]
\end{lstlisting}
GitLab nadpisuje \repopath{ENTRYPOINT}, więc tam DLL
wywołuje się jawnie. Pakiet GHCR powstaje przy pierwszym pushu jako \textbf{prywatny},
niezależnie od widoczności repozytorium; trzeba go raz upublicznić w~ustawieniach pakietu.

\subsubsection{GitHub Action (\texttt{action.yml})}

\begingroup\small
\begin{xltabular}{\linewidth}{@{}>{\raggedright\arraybackslash}p{2.6cm} >{\raggedright\arraybackslash}p{2.4cm} >{\raggedright\arraybackslash}X@{}}
\toprule
\textbf{Wejście} & \textbf{Domyślnie} & \textbf{Działanie} \\
\midrule
\endfirsthead
\toprule
\textbf{Wejście} & \textbf{Domyślnie} & \textbf{Działanie} \\
\midrule
\endhead
\texttt{path} & \texttt{src} & Katalog skanu względem katalogu głównego repozytorium. \\
\texttt{version} & \texttt{0.1.0} & Wersja instalowanych narzędzi; ma równać się wersji akcji. \\
\texttt{diff-scoped} & \texttt{true} & Na zdarzeniu \texttt{pull\_request} blokować mogą tylko naruszenia na zmienionych liniach, a~linie AI mają ostrzejszy próg. \texttt{false}: każdy błąd w~drzewie blokuje (właściwe dla audytu, zwykle nie dla bramki PR). \\
\texttt{comment} & \texttt{true} & Publikuj lub aktualizuj komentarz sticky. \\
\texttt{fail-on-block} & \texttt{true} & \texttt{false}: raportuj bez egzekwowania (okres wdrażania). \\
\texttt{github-token} & \texttt{github.token} & Token do komentarza; wymaga \texttt{pull-requests: write}. \\
\texttt{dotnet-version} & \texttt{10.0.x} & SDK do zainstalowania; pusty ciąg -- użyj SDK runnera. \\
\bottomrule
\end{xltabular}
\endgroup

Wyjścia: \repopath{blocked}, \repopath{error-count} (wszystkie błędy w~skanie),
\repopath{blocking-error-count} (błędy liczone przez bramkę w~aktywnym zakresie),
\repopath{result-path} (ścieżka surowego JSON). Kroki akcji:

\begin{enumerate}[leftmargin=*]
  \item \repopath{actions/setup-dotnet@v6}, jeśli \repopath{dotnet-version} nie jest puste.
  \item Instalacja narzędzi: \repopath{Portcullis.Cli} przez
        \repopath{dotnet tool install --global} w~wersji z~wejścia \repopath{version}, a~pakiet
        \repopath{Portcullis.CiComment} tylko wtedy, gdy wejście \repopath{comment} ma wartość
        \repopath{true}.
  \item Skan \repopath{\$GITHUB\_WORKSPACE/\$PC\_PATH}. Przy \repopath{diff-scoped: true}
        i~znanym commicie bazowym PR akcja sprawdza (\repopath{git cat-file -e}), czy commit
        jest w~checkoucie; jeśli tak, dodaje \repopath{--provenance-range base..sha}
        i~\repopath{--provenance-repo}, jeśli nie -- wypisuje \repopath{::warning::}
        z~zaleceniem \repopath{fetch-depth: 0} i~skanuje całe drzewo. Kod wyjścia większy
        niż 1 to błąd użycia i~kończy krok; kod 1 jest werdyktem. Wyjścia liczy krótki
        skrypt Python z~JSON (z~fallbackiem na \repopath{summary.errorCount}, gdy brak
        \repopath{gate}).
  \item Publikacja komentarza (\repopath{comment: true} i~zdarzenie PR) -- błąd publikacji
        jest ignorowany (\repopath{|| true}), bo werdykt egzekwuje następny krok, więc
        porażka komentarza nie może udawać zielonej bramki. Akcja nie przekazuje
        \repopath{--previous}, więc jej komentarz nie ma sekcji zmian od poprzedniego pushu.
  \item Egzekwowanie (\repopath{fail-on-block: true}): \repopath{::error::} i~kod 1, gdy
        \repopath{blocked}.
\end{enumerate}

\subsubsection{Wersjonowanie i~wydanie}

Wersje stosują Semantic Versioning, a~CHANGELOG format Keep a~Changelog. Wersja jest jednym
pokrętłem: \repopath{<Version>0.1.0</Version>} w~\repopath{Directory.Build.props}, a~workflow
wydania przekazuje \repopath{-p:Version=<tag>}, więc pakiety NuGet, narzędzia i~obraz niosą
numer tagu, który je zbudował. Druga kopia wersji to domyślne wejście \repopath{version}
w~\repopath{action.yml}: ruchomy tag \repopath{v0} wskazuje otagowany commit, więc
\repopath{uses: ...@v0} uruchamia \repopath{action.yml} z~tego commita, a~jeśli jego domyślna
wersja narzędzi nie istnieje w~NuGet, psuje się workflow każdego konsumenta. Wydanie
według \repopath{CONTRIBUTING.md}:

\begin{enumerate}[leftmargin=*]
  \item Podbić \repopath{<Version>} w~\repopath{Directory.Build.props}.
  \item Podbić domyślne wejście \repopath{version} w~\repopath{action.yml} na tę samą wartość.
  \item Zaktualizować przypięte wersje w~dokumentacji: w~\path{docs/TUTORIAL.md},
        w~\path{docs/index.html} i~w~plikach \repopath{PACKAGE-README.md}.
  \item Commit, potem tag:
\begin{lstlisting}
git tag v0.2.0 && git push origin v0.2.0
\end{lstlisting}
\end{enumerate}

Push tagu \repopath{v*} buduje, testuje, pakuje, publikuje trzy pakiety NuGet, wypycha obraz
do GHCR, tworzy wydanie GitHub i~przesuwa tag główny (\repopath{v0}) na ten commit
(szczegóły jobów -- sekcja~\ref{sec:portcullis-ci}). Jednorazowo przed pierwszym wydaniem:
sekret repozytorium \repopath{NUGET\_API\_KEY} (klucz z~nuget.org ograniczony wzorcem
\repopath{Portcullis.*}, żeby działał dla pakietów, które jeszcze nie istnieją) i~po
pierwszym pushu obrazu -- upublicznienie pakietu GHCR.

\zrodla{\path{letsgolegacy.portcullis/docs/DISTRIBUTION.md}, \path{letsgolegacy.portcullis/action.yml},
\path{letsgolegacy.portcullis/Dockerfile}, \path{letsgolegacy.portcullis/Directory.Build.props},
\path{letsgolegacy.portcullis/CONTRIBUTING.md}, \path{letsgolegacy.portcullis/CHANGELOG.md},
\path{letsgolegacy.portcullis/docs/TUTORIAL.md},
\path{letsgolegacy.portcullis/src/Portcullis.Rules/Portcullis.Rules.csproj},
\path{letsgolegacy.portcullis/src/Portcullis.Cli/Portcullis.Cli.csproj},
\path{letsgolegacy.portcullis/src/Portcullis.CiComment/Portcullis.CiComment.csproj},
\path{letsgolegacy.portcullis/.github/workflows/release.yml}}

\subsection{Polityka wersji Roslyn i~analyzer floor}\label{sec:portcullis-roslyn}

Repozytorium celowo używa dwóch różnych wersji \repopath{Microsoft.CodeAnalysis.CSharp}:

\begin{itemize}[leftmargin=*]
  \item \textbf{\texttt{Portcullis.Rules}: 4.8.0 -- najstarsza, a~nie najnowsza.} Analyzer
        zbudowany na nowszym Roslyn niż kompilator konsumenta daje \repopath{CS9057}
        w~każdym buildzie każdego repozytorium konsumenta. 4.8.0 to kompilator pasma 1xx
        SDK .NET 8 (późniejsze pasma 8.0 dostarczają 4.9--4.11), więc pakiet ładuje się na
        .NET 8 i~każdym nowszym SDK. Referencja ma \repopath{PrivateAssets} równe
        \repopath{all}, obok niej
        \repopath{Microsoft.CodeAnalysis.Analyzers} 3.11.0 i~\repopath{EnforceExtendedAnalyzerRules}.
        Podniesienie tej podłogi odcina każdego konsumenta na starszym SDK -- to decyzja
        o~wspieranych SDK, nie aktualizacja.
  \item \textbf{\texttt{Portcullis.Engine}: 5.9.0 -- najnowsza.} To parser skanera,
        a~skanowany kod to kandydat migracji: projekt .NET 10 domyślnie kompiluje C\#~14,
        który starszy Roslyn czyta jako błędy składni, a~to, co leży wewnątrz, analizuje
        z~drzewa po naprawie błędów albo wcale. \repopath{ScannerLanguageVersionTests}
        sprawdza parsowanie bloku \repopath{extension} i~słowa \repopath{field} bez błędów
        oraz wykrycie \repopath{.Result} wewnątrz bloku \repopath{extension}.
\end{itemize}

Cel \repopath{netstandard2.0} i~stary Roslyn mają konsekwencje w~kodzie reguł: brak
przeciążenia \repopath{Split} z~parametrami \repopath{params char[]}
i~\repopath{StringSplitOptions} (stąd pole \repopath{PathSeparators}), brak \repopath{[NotNullWhen]} na \repopath{string.IsNullOrEmpty}
(stąd \repopath{x is null || x.Length == 0}), brak \repopath{System.Index}/\repopath{Range}
oraz \repopath{CS9210} przy wyrażeniach kolekcji na \repopath{ImmutableArray} z~tej wersji
\repopath{System.Collections.Immutable}.

\paragraph{Dependabot.} Dependabot aktualizuje pakiety NuGet co tydzień. Biblioteki testowe
tworzą jedną grupę \repopath{test-stack} (wzorce \repopath{xunit*} i~\repopath{coverlet*} oraz
pakiet \path{Microsoft.NET.Test.Sdk}). Rodzina \repopath{Microsoft.CodeAnalysis*} jest
wpisana do \repopath{ignore}: podnosi się ją ręcznie, łącznie z~parserem silnika
\repopath{Portcullis.Engine}. Osobne wpisy co tydzień aktualizują ekosystem
\repopath{github-actions} (workflowy oraz \repopath{action.yml} w~katalogu głównym)
i~ekosystem \repopath{docker}.

\paragraph{Job \texttt{analyzer-floor}.} Podłoga jest egzekwowana przez CI, a~nie pamiętana.
Job „Analyzer package on the .NET 8 SDK” w~\repopath{ci.yml}:
\begin{enumerate}[leftmargin=*]
  \item instaluje SDK \repopath{8.0.1xx} i~\repopath{10.0.x};
  \item pakuje \repopath{src/Portcullis.Rules} w~wersji
        \repopath{0.1.0-ci.<run\_id>.<run\_attempt>} do
        \path{fixtures/analyzer-consumer/feed};
  \item buduje projekt \path{fixtures/analyzer-consumer}. Plik \repopath{global.json}
        przypina SDK w~wersji \repopath{8.0.100} z~polityką \repopath{latestPatch}, a~job
        przerywa się, jeśli wersja SDK nie należy do pasma \repopath{8.0.1xx}. Projekt
        docelowo \repopath{net8.0} traktuje ostrzeżenia \repopath{CS9057} i~\repopath{CS8032}
        jako błędy, więc analyzer zbudowany na nowszym Roslyn przerywa build. Plik
        \repopath{nuget.config} kieruje pakiety \repopath{Portcullis.*} wyłącznie do
        lokalnego feedu;
  \item sprawdza, że log builda zawiera poniższą linię -- plik celowo blokuje na tasku,
        więc to dowód, że kompilator analyzery uruchomił, a~nie tylko ich nie odrzucił:
\begin{lstlisting}
PendingOrders.cs(11,55): warning PORTCULLIS_MIG_SYNC_OVER_ASYNC
\end{lstlisting}
\end{enumerate}

Po stronie konsumenta tabela problemów tutoriala mówi: \repopath{warning CS9057} oznacza SDK
starszy niż .NET 8.

\zrodla{\path{letsgolegacy.portcullis/src/Portcullis.Rules/Portcullis.Rules.csproj},
\path{letsgolegacy.portcullis/src/Portcullis.Engine/Portcullis.Engine.csproj},
\path{letsgolegacy.portcullis/CONTRIBUTING.md}, \path{letsgolegacy.portcullis/docs/DISTRIBUTION.md},
\path{letsgolegacy.portcullis/.github/dependabot.yml}, \path{letsgolegacy.portcullis/.github/workflows/ci.yml},
\path{letsgolegacy.portcullis/fixtures/analyzer-consumer/AnalyzerConsumer.csproj},
\path{letsgolegacy.portcullis/fixtures/analyzer-consumer/global.json},
\path{letsgolegacy.portcullis/fixtures/analyzer-consumer/nuget.config},
\path{letsgolegacy.portcullis/fixtures/analyzer-consumer/PendingOrders.cs},
\path{letsgolegacy.portcullis/src/Portcullis.Rules/PathConventions.cs},
\path{letsgolegacy.portcullis/src/Portcullis.Rules/PortcullisConventions.cs},
\path{letsgolegacy.portcullis/tests/Portcullis.Engine.Tests/ScannerLanguageVersionTests.cs},
\path{letsgolegacy.portcullis/docs/TUTORIAL.md}}

\subsection{Kamienie milowe i~zapisy weryfikacji}\label{sec:portcullis-bootstrap}

Portcullis powstał w~prywatnym repozytorium-poprzedniku i~został zaimportowany bez jego
historii git (ticket R0). Dokumenty \repopath{BOOTSTRAP}, \repopath{M4-INTEGRATION},
\repopath{FINDINGS} i~\repopath{MUTATIONS} to datowane zapisy z~tamtego okresu; identyfikatory
w~nich zaktualizowano do nazwy Portcullis, a~same ustalenia zostawiono bez zmian. Wszystkie
opisują to, co naprawdę uruchomiono, z~poleceniami i~wynikami. Repozytorium, na którym
testowano reguły, to prywatna aplikacja .NET Aspire, cytowana jako \repopath{<consumer>}
(dalej: repozytorium referencyjne).

Etapy: M0 -- SPEC; M1D -- fixture \path{fixtures/consumer-violations.json} (lista „musi
złapać”, 11~wpisów); M2 -- szkielet skanera bez reguł; trzy równoległe tory: A~(\repopath{track-a-engine}: reguły P2, P9, P10 i~pierwszy przebieg mutacyjny, M2A/M3A),
B (\repopath{track-b-ci}: akcja i~komentarz PR na makiecie zgodnej z~SPEC §3) i~C
(\repopath{track-c-provenance}: pochodzenie z~git); M4 -- integracja; M5 -- atak na cały
łańcuch; M5A -- bramka zawężona do diffu; M6 -- dogfooding na repozytorium referencyjnym,
zaplanowany po M5A. Potem tickety R0--R3 publicznego repozytorium
(sekcja~\ref{sec:portcullis-status}).

\subsubsection{Bootstrap (M0, M1D, M2)}

Zapis z~2026-08-17, z~końca sesji bootstrapu. Uruchomiono: \repopath{dotnet build} całego
rozwiązania (0 ostrzeżeń, 0 błędów), \repopath{dotnet test} (3 z~3 testów: pusty wynik przy
zerze reguł, pomijanie \repopath{bin} i~\repopath{obj}, pusty wynik zamiast wyjątku dla
nieistniejącej ścieżki), skan repozytorium referencyjnego (156 plików, 0 reguł, 0 naruszeń,
kod 0 -- przykład z~SPEC §3) i~skan własnego \repopath{src/} (6 plików). JSON porównano
z~tabelą SPEC pole po polu. Wszystkie 11 wpisów fixture'a zweryfikowano niezależnie na
źródłach repozytorium referencyjnego.

Ustalenie sesji: z~trzech kategorii naruszeń, które założenia startowe wskazywały „jako
minimum”, dwie już się nie odtwarzały -- naprawił je commit sprzed audytu, z~którego
pochodziły założenia, a~README aplikacji tego nie odnotował. Fixture nie podmienił
przykładów po cichu: \repopath{CONSUMER-001} to teraz sama nieaktualność README (P14),
a~\repopath{CONSUMER-002} to rzeczywisty kształt drugiej luki -- przycisk klienta web, który
w~ogóle nie wywołuje istniejącego, poprawnie zaimplementowanego endpointu (kod TypeScript). Oba mają
\repopath{roslynCheckable: false}. Jedyna kategoria sprawdzalna Roslynem (encje EF Core bez
kontrolera) urosła do dziewięciu wpisów (\repopath{CONSUMER-003}--\repopath{CONSUMER-011}).
Niezbudowane były wtedy: reguły, przebieg mutacyjny, akcja CI, implementacja pochodzenia
i~analiza między projektami (skaner referencjonował tylko \repopath{System.Private.CoreLib}).

\subsubsection{Integracja M4}\label{sec:portcullis-m4}

Zapis z~2026-08-17, po scaleniu torów A, B i~C. \textbf{Pierwsze ustalenie: tory nigdy nie
były razem.} Implementacja \repopath{GitProvenanceProvider} toru C trafiła tylko do gałęzi
\repopath{track-c-provenance}, a~\repopath{main} miał wyłącznie kontrakt
\repopath{ProvenanceModels.cs}. Scalenie (\repopath{git merge --no-ff}) przeszło bez
konfliktów. Build: 0 ostrzeżeń; testy: 59 z~59 (23~CiComment i~36~Engine, w~tym 14 testów
dostawcy pochodzenia), po dodaniu okablowania 67 z~67 (8 nowych
\repopath{ScannerProvenanceTests}).

\textbf{Decyzja projektowa: eskalacja raz, w~skanerze.} SPEC mówi „reguła może czytać
raport pochodzenia”, co sugeruje obsługę w~każdym analyzerze -- przez
\repopath{AnalyzerOptions} lub \repopath{AdditionalFiles}, z~tą samą kontrolą nakładania
linii powielaną w~każdej regule. M4 wybrał jeden przebieg po zebranych naruszeniach:
jeden testowany punkt polityki zamiast trzech i~więcej, testy reguł niezależne od pochodzenia,
identyczny efekt. Odstępstwo od literalnego brzmienia SPEC zapisano jawnie.
\repopath{Scanner.ScanAsync} dostał opcjonalne parametry raportu i~katalogu repozytorium --
raport dostarcza wywołujący (silnik nie zna git), a~katalog repozytorium jest potrzebny, bo
ścieżki git są względne wobec katalogu głównego repozytorium, a~workflow CI skanuje
podkatalog \repopath{src}. CLI dostało \repopath{--provenance-range}
i~\repopath{--provenance-repo}. Testy sprawdzają: eskalację ostrzeżenia w~zakresie AI,
brak zmiany błędu, dopasowanie po linii (nie po pliku), brak eskalacji dla zakresów
\repopath{human} i~\repopath{mixed}, zachowanie bez raportu oraz oba warianty katalogu
repozytorium.

\textbf{Przebieg na repozytorium referencyjnym.} Bez pochodzenia: 156 plików, 3~reguły,
15~naruszeń (4~błędy, 11~ostrzeżeń); trzy z~nich to dokładnie wpisy
\repopath{CONSUMER-003}, \repopath{CONSUMER-005} i~\repopath{CONSUMER-006}
(\path{PORTCULLIS_P9_ORPHAN_ENTITY}), 4~błędy to \path{PORTCULLIS_P9_CONTROLLER_NO_DBCONTEXT}.
Z~pochodzeniem od pierwszego commita do \repopath{HEAD}: 9~błędów, 6~ostrzeżeń -- wszystkie
5~ostrzeżeń \path{PORTCULLIS_P10_CUSTOM_BASE_CLASS} (klasy klienta web) eskalowały.
Sprawdzono ręcznie przez \repopath{git blame} i~\repopath{git log}: jeden eskalowany plik
pochodzi z~commita, którego autorem jest tożsamość agenta AI z~odpowiednim trailerem
\repopath{Co-Authored-By}, a~nieeskalowany \repopath{QrCode.cs:11} -- z~commita człowieka bez
trailera AI. Łańcuch z~\repopath{Portcullis.CiComment} wyrenderował komentarz i~zakończył się
kodem 1; bez \repopath{--pr} i~\repopath{GITHUB\_TOKEN} narzędzie pominęło publikację zamiast
się wywrócić.

\textbf{Luka zostawiona świadomie:} bramka była bezwzględna, więc repozytorium referencyjne
blokowałoby każdy PR. Dodatkowo 9~z~11~wpisów fixture'a to kształt
\repopath{ORPHAN\_ENTITY} o~poziomie \repopath{warning}: ich ponowne wprowadzenie przez
człowieka nie zatrzymałoby bramki bez eskalacji AI. Decyzję o~bramce zawężonej do diffu
zostawiono na M5/M6; podjął ją M5A.

\subsubsection{Atak end-to-end M5 (FINDINGS)}

Zapis z~2026-08-17. M3A dowiódł, że każda reguła osobno zależy od swojego mechanizmu; M5
sprawdza, czy cały łańcuch -- reguły, eskalacja pochodzenia i~bramka CI -- daje poprawną
decyzję o~merge'u dla realistycznych kształtów PR. Dwa źródła dowodu: migawka repozytorium
referencyjnego (156 plików, 6~reguł, 4~błędy, 14~ostrzeżeń: 4~\repopath{CONTROLLER\_NO\_DBCONTEXT},
6~\repopath{ORPHAN\_ENTITY}, 5~\repopath{CUSTOM\_BASE\_CLASS}, 3~\repopath{VENDOR\_SDK\_OUTSIDE\_ADAPTER},
zero P4 i~P15) oraz jednorazowe repozytorium robocze z~pięcioma commitami, każdy izolujący
jedno zachowanie.

\begingroup\small
\begin{xltabular}{\linewidth}{@{}l >{\raggedright\arraybackslash}X >{\raggedright\arraybackslash}p{3.6cm} l l l@{}}
\toprule
\textbf{\#} & \textbf{Stan (commit)} & \textbf{Polecenie} & \textbf{Błędy} & \textbf{Ostrz.} & \textbf{Kod} \\
\midrule
\endfirsthead
\toprule
\textbf{\#} & \textbf{Stan (commit)} & \textbf{Polecenie} & \textbf{Błędy} & \textbf{Ostrz.} & \textbf{Kod} \\
\midrule
\endhead
1 & B: nowa encja-sierota, człowiek & \texttt{scan} & 0 & 1 & \textbf{0} \\
2a & C: druga encja-sierota, commit z~trailerem agenta AI & \texttt{scan} & 0 & 2 & \textbf{0} \\
2b & C & \texttt{scan -{}-provenance-range B..C} & 1 & 1 & \textbf{1} \\
3 & D: kontroler z~\texttt{DbContext}, człowiek & \texttt{scan} & 1 & 2 & \textbf{1} \\
4 & E: niezwiązany, czysty helper & \texttt{scan} & 1 & 2 & \textbf{1} \\
\bottomrule
\end{xltabular}
\endgroup

Każdy kod wyjścia to dosłowny kod procesu (\repopath{echo \$?}), nie wniosek z~JSON.
Scenariusz 2: dwa strukturalnie identyczne naruszenia -- ta sama reguła, ten sam poziom
domyślny -- rozchodzą się w~wyniku zależnie od autora; naruszenie z~commita agenta AI
eskaluje do błędu i~blokuje. Scenariusz 4 to atak na projekt bramki: diff commita E to jeden
nowy plik (6 linii), a~bramka i~tak blokowała, bo liczyła błędy w~całym drzewie. Wnioski:
eskalacja uogólnia się poza jeden historyczny przykład; komentarz, liczby i~kod wyjścia
zawsze się zgadzały; luka „bezwzględna czy diff” odtwarza się na żądanie; mechanizmy
eskalacji i~bramki działają tylko na \repopath{filePath} i~\repopath{line}, więc nie zależą
od reguły. Zapis nie twierdzi, że reguły są wolne od fałszywych alarmów, nie sprawdza
ponownie publikacji na GitHub (potwierdzonej w~M4 na prawdziwym przebiegu Actions
w~repozytorium-poprzedniku) i~nie dotyka M6 -- naprawa repozytorium referencyjnego
zniszczyłaby fixture, z~którego korzystają wszystkie etapy.

\zrodla{\path{letsgolegacy.portcullis/docs/BOOTSTRAP.md}, \path{letsgolegacy.portcullis/docs/M4-INTEGRATION.md},
\path{letsgolegacy.portcullis/docs/FINDINGS.md}, \path{letsgolegacy.portcullis/docs/SPEC.md},
\path{letsgolegacy.portcullis/docs/DIFF-GATE.md}, \path{letsgolegacy.portcullis/CHANGELOG.md},
\path{letsgolegacy.portcullis/fixtures/consumer-violations.json},
\path{letsgolegacy.portcullis/tests/Portcullis.Engine.Tests/ScannerTests.cs},
\path{letsgolegacy.portcullis/tests/Portcullis.Engine.Tests/ScannerProvenanceTests.cs}}

\subsection{Tutorial: procedura wdrożenia}\label{sec:portcullis-tutorial}

Procedura przenosi kroki \repopath{docs/TUTORIAL.md}. Zastrzeżenie tutoriala: pakiety
i~obraz są zbudowane i~zweryfikowane lokalnie, ale nieopublikowane, więc do pierwszego
tagu \repopath{v*} każde \repopath{dotnet add package} i~\repopath{dotnet tool install}
kończy się „package not found”; działa ścieżka ze źródeł (krok~7). Poza przykładami
z~kontenerem każde polecenie tutoriala uruchomiono przed zapisaniem.

\paragraph{Krok 1. Wybór ścieżki.} Ścieżki się składają: większość zespołów używa pakietu
analyzerów do szybkiej informacji lokalnej i~akcji jako bramki.

\begin{xltabular}{\linewidth}{@{}>{\raggedright\arraybackslash}X >{\raggedright\arraybackslash}p{5.2cm} l@{}}
\toprule
\textbf{Cel} & \textbf{Narzędzie} & \textbf{Nakład} \\
\midrule
\endfirsthead
\toprule
\textbf{Cel} & \textbf{Narzędzie} & \textbf{Nakład} \\
\midrule
\endhead
Naruszenia psują build wszędzie, także lokalnie & pakiet \texttt{Portcullis.Analyzers} & jedna linia \\
Bramka PR z~komentarzem na GitHub & akcja & jeden plik workflow \\
To samo na GitLab, Azure DevOps, Jenkins & narzędzie \texttt{portcullis} lub kontener & kilka linii \\
Rozpoznanie przed decyzją & lokalne \texttt{portcullis scan} & jedno polecenie \\
\bottomrule
\end{xltabular}

\paragraph{Krok 2. Pakiet analyzerów.} \repopath{dotnet add package Portcullis.Analyzers}
wystarcza: przy następnym \repopath{dotnet build} naruszenia są zwykłymi diagnostykami
kompilatora, np.:

\begin{lstlisting}
Controllers/WidgetsController.cs(5,22): error PORTCULLIS_P9_CONTROLLER_NO_DBCONTEXT:
    Controller 'WidgetsController' references 'WidgetDbContext' directly. Controllers are
    transport only — bind, authorize, delegate — route through an orchestrator or
    repository instead (architecture-standards P9).
Build FAILED.
\end{lstlisting}

Dla wszystkich projektów naraz -- w~\repopath{Directory.Build.props}; \repopath{PrivateAssets}
chroni przed wejściem Portcullis do zależności publikowanego pakietu zespołu, a~wersja ma być
przypięta (gwiazdka w~\repopath{Version} pozwala regule z~nowego wydania zepsuć build bez
zmiany kodu):

\begin{lstlisting}
<Project>
  <ItemGroup>
    <PackageReference Include="Portcullis.Analyzers" Version="0.1.0" PrivateAssets="all" />
  </ItemGroup>
</Project>
\end{lstlisting}

Ta ścieżka nie ma zawężenia do diffu: każde naruszenie w~każdym pliku przy każdym buildzie
(zob. krok~4).

\paragraph{Krok 3. Bramka na pull request.} \emph{GitHub Actions} -- kompletna bramka
(wejścia i~wyjścia akcji: sekcja~\ref{sec:portcullis-dystrybucja}):

\begin{lstlisting}
# .github/workflows/architecture.yml
name: Architecture
on: pull_request

permissions:
  contents: read
  pull-requests: write        # required for the sticky comment

jobs:
  portcullis:
    runs-on: ubuntu-latest
    steps:
      - uses: actions/checkout@v4
        with:
          fetch-depth: 0      # required — see the warning below
      - uses: konradcinkusz/letsgolegacy.portcullis@v0
        with:
          path: src
\end{lstlisting}

\repopath{fetch-depth: 0} nie jest opcjonalne: płytki klon nie zawiera commita bazowego PR,
więc zawężenie do diffu nie może się odbyć; akcja wypisuje wtedy \repopath{::warning::}
i~przechodzi na całe drzewo. SARIF do code scanning wymaga dziś własnego workflow z~CLI
i~\repopath{github/codeql-action/upload-sarif} (sekcja~\ref{sec:portcullis-ci}). Wyjścia
akcji pozwalają reagować w~kolejnych krokach:

\begin{lstlisting}
      - uses: konradcinkusz/letsgolegacy.portcullis@v0
        id: portcullis
        with: { path: src, fail-on-block: false }
      - if: steps.portcullis.outputs.blocked == 'true'
        run: echo "would have blocked on ${{ steps.portcullis.outputs.blocking-error-count }} violation(s)"
\end{lstlisting}

\emph{GitLab CI} -- obraz nie wymaga SDK na runnerze; GitLab nadpisuje
\repopath{ENTRYPOINT}, więc DLL wywołuje się jawnie. Job kończy się błędem, gdy CLI zwraca 1:

\begin{lstlisting}
portcullis:
  image: ghcr.io/konradcinkusz/letsgolegacy.portcullis:0.1.0
  rules:
    - if: $CI_PIPELINE_SOURCE == "merge_request_event"
  variables:
    GIT_DEPTH: 0              # same reason as fetch-depth: 0 above
  script:
    - dotnet /app/Portcullis.Cli.dll scan "$CI_PROJECT_DIR/src"
        --provenance-range "$CI_MERGE_REQUEST_DIFF_BASE_SHA..$CI_COMMIT_SHA"
        --provenance-repo "$CI_PROJECT_DIR"
        | tee portcullis.json
  artifacts:
    when: always
    paths: [portcullis.json]
\end{lstlisting}

\emph{Azure DevOps}:

\begin{lstlisting}
- task: UseDotNet@2
  inputs: { version: '10.0.x' }

- script: |
    dotnet tool install -g Portcullis.Cli --version 0.1.0
    export PATH="$PATH:$HOME/.dotnet/tools"
    BASE=$(git merge-base "origin/$(System.PullRequest.TargetBranch)" HEAD)
    portcullis scan "$(Build.SourcesDirectory)/src" \
      --provenance-range "$BASE..HEAD" \
      --provenance-repo "$(Build.SourcesDirectory)"
  displayName: Portcullis
\end{lstlisting}

\emph{Dowolne inne CI} (Jenkins, Woodpecker, Makefile) -- \repopath{git merge-base} działa
wszędzie i~nie wymaga zmiennych CI; wariant z~kontenerem nie wymaga .NET:

\begin{lstlisting}
BASE=$(git merge-base origin/main HEAD)
portcullis scan ./src --provenance-range "$BASE..HEAD" --provenance-repo .

docker run --rm -v "$PWD:/workspace" ghcr.io/konradcinkusz/letsgolegacy.portcullis:0.1.0 \
  scan /workspace/src --provenance-range "$BASE..HEAD" --provenance-repo /workspace
\end{lstlisting}

\paragraph{Krok 4. Wdrożenie na kodzie z~istniejącymi naruszeniami.} Według tutoriala tu
umiera większość wdrożeń bramek architektury: bramka blokująca na całym długu blokuje każdy
PR i~w~ciągu tygodnia ktoś ją wyłącza. Trzy niezależne odpowiedzi, które się składają:
\begin{enumerate}[leftmargin=*]
  \item \textbf{Zawężenie do diffu} (domyślnie w~akcji) -- wszystko jest raportowane,
        blokuje tylko to, co w~zmienionych liniach. Dług jest widoczny w~każdym PR, ale nie
        jest niczyim alarmem, i~maleje w~miarę dotykania plików.
  \item \textbf{Stopniowe zaostrzanie w~\texttt{.editorconfig}} -- per reguła i~per
        ścieżka (kanał pakietu analyzerów). Sprawdzono: podniesienie
        \repopath{ORPHAN\_ENTITY} do \repopath{error} zepsuło przechodzący build, a~ustawienie
        \repopath{CONTROLLER\_NO\_DBCONTEXT} na \repopath{none} usunęło błąd. Rozsądny plan:
        pierwszego dnia wszystko na \repopath{suggestion}, potem po jednej regule --
        naprawa, \repopath{warning}, \repopath{error}; każdy krok to jedna linia.
  \item \textbf{Baseline} -- raz \repopath{--write-baseline}, commit pliku, potem
        \repopath{--baseline} w~każdym skanie (sekcja~\ref{sec:portcullis-baseline}).
\end{enumerate}

\begin{lstlisting}
# .editorconfig
root = true

[*.cs]
# start noisy rules as suggestions while you pay down existing violations
dotnet_diagnostic.PORTCULLIS_P10_CUSTOM_BASE_CLASS.severity = suggestion

# turn one all the way off for now
dotnet_diagnostic.PORTCULLIS_P4_SEED_DATA_IN_MODEL.severity = none

# and promote one you are already clean on, so it can never regress
dotnet_diagnostic.PORTCULLIS_P9_ORPHAN_ENTITY.severity = error

# legacy code is exempt; new code is not
[src/Legacy/**.cs]
dotnet_diagnostic.PORTCULLIS_P9_CONTROLLER_NO_DBCONTEXT.severity = none
\end{lstlisting}

\begin{lstlisting}
portcullis scan ./src --baseline portcullis-baseline.json --write-baseline
\end{lstlisting}

\paragraph{Krok 5. Odczyt wyniku.} \repopath{summary} to to, co \textbf{znaleziono};
\repopath{gate} to to, co \textbf{blokuje}; różnica to istniejący dług. Pełny schemat --
sekcja~\ref{sec:portcullis-spec}. Kody wyjścia: 0 -- nie blokuje, 1 -- blokuje, 2 -- błędne
użycie; tylko 2 oznacza awarię Portcullis, 1 to werdykt.

\begin{lstlisting}
{
  "filesScanned": 156,
  "rulesEvaluated": 7,
  "violations": [
    { "ruleId": "PORTCULLIS_P9_ORPHAN_ENTITY", "filePath": "Domain/QrCode.cs",
      "line": 11, "message": "...", "severity": "warning" }
  ],
  "summary": { "errorCount": 4, "warningCount": 14, "infoCount": 0 },
  "gate": { "blocked": true, "scope": "diff", "blockingErrorCount": 5 }
}
\end{lstlisting}

\paragraph{Krok 6. Rozwiązywanie problemów.}

\begingroup\small
\begin{xltabular}{\linewidth}{@{}>{\raggedright\arraybackslash}p{4.6cm} >{\raggedright\arraybackslash}X@{}}
\toprule
\textbf{Objaw} & \textbf{Przyczyna i~poprawka} \\
\midrule
\endfirsthead
\toprule
\textbf{Objaw} & \textbf{Przyczyna i~poprawka} \\
\midrule
\endhead
\texttt{gate.scope} to \texttt{all}, a~oczekiwano \texttt{diff} & Commita bazowego nie ma w~checkoucie. Dodaj \texttt{fetch-depth: 0} (GitHub) lub \texttt{GIT\_DEPTH: 0} (GitLab); akcja ostrzega o~tym jawnie. \\
Bramka przechodzi, choć powinna blokować & Sprawdź \texttt{gate.scope}. Przy \texttt{diff} naruszenie leży poza zmienionymi liniami PR -- bramka działa zgodnie z~projektem. Przy \texttt{all} z~\texttt{degradedReason} git nie rozwiązał zakresu i~skan przeszedł na całe drzewo; powód mówi, co zawiodło. Awaria git już nie przechodzi po cichu. \\
Pakiet się instaluje, ale diagnostyk nie ma & Sprawdź, czy DLL leży w~\texttt{analyzers/dotnet/cs/} wewnątrz \texttt{.nupkg}; analyzer spakowany do \texttt{lib/} instaluje się bez błędu i~nic nie robi. \\
\texttt{warning CS9057: ... references version 4.x of the compiler, which is newer} & Roslyn SDK jest starszy niż ten, na którym zbudowano Portcullis. \texttt{Portcullis.Analyzers} celuje w~Roslyn 4.8 (wersja SDK .NET 8) właśnie po to, żeby tego uniknąć; objaw oznacza SDK starszy niż .NET 8. \\
\texttt{docker pull} obrazu GHCR zwraca „unauthorized” & Pakiet GHCR powstaje przy pierwszym pushu jako prywatny, niezależnie od widoczności repozytorium; trzeba go raz upublicznić w~ustawieniach pakietu. \\
Reguły nie strzelają na kodzie, który na pewno jest zły & Portcullis parsuje źródła bezpośrednio, nie przez MSBuild, więc typ z~pakietu NuGet lub innego projektu się nie rozwiązuje; część reguł na tym polega, inne są przez to stępione (README, „Honest limits”). \\
Reguła strzela na poprawnym kodzie & To błąd do zgłoszenia, nie do obchodzenia -- hałaśliwa bramka jest gorsza niż brakująca reguła. Wycisz ją w~\texttt{.editorconfig} (krok~4) i~otwórz issue z~fragmentem kodu. \\
\bottomrule
\end{xltabular}
\endgroup

\paragraph{Krok 7. Przed pierwszym wydaniem: uruchomienie ze źródeł.} Wymaga tylko SDK
.NET 10:

\begin{lstlisting}
git clone https://github.com/konradcinkusz/letsgolegacy.portcullis
cd letsgolegacy.portcullis
dotnet build portcullis.sln

# scan any repository on your machine
dotnet run --project src/Portcullis.Cli -- scan /path/to/your-repo/src

# with the diff-scoped gate
cd /path/to/your-repo
BASE=$(git merge-base origin/main HEAD)
dotnet run --project /path/to/letsgolegacy.portcullis/src/Portcullis.Cli -- \
  scan ./src --provenance-range "$BASE..HEAD" --provenance-repo .

# the analyzer package from a local feed (add ./artifacts as a source in nuget.config)
dotnet pack src/Portcullis.Rules/Portcullis.Rules.csproj -c Release -o ./artifacts
\end{lstlisting}

\paragraph{Krok 8. Czego jeszcze nie da się zrobić.} Kształt reguły jest stały, choć
konwencje są konfigurowalne; wdrożono sześć z~piętnastu zasad; obraz kontenera i~workflow
wydania nie były wykonane; awaria git degraduje bramkę do całego drzewa zamiast ją
przepuszczać. Jeśli coś w~tutorialu nie działa tak, jak napisano, to błąd dokumentacji wart
zgłoszenia.

\zrodla{\path{letsgolegacy.portcullis/docs/TUTORIAL.md}, \path{letsgolegacy.portcullis/action.yml},
\path{letsgolegacy.portcullis/docs/SARIF.md}}

\subsection{Testy mutacyjne reguł}\label{sec:portcullis-mutacje}

Zasada repozytorium: zestaw testów, który nigdy nie zawiódł, to zestaw, którego nikt nie
przetestował. Reguła nie jest gotowa, gdy strzela, tylko gdy jest dowód, że strzela
\emph{z~właściwego powodu}. Test mutacyjny to dowód, że testy reguły zauważyłyby, gdyby
mechanizm, od którego reguła zależy, przestał działać.

\paragraph{Metoda.} Mutanty żyją w~kodzie testów (\path{tests/Portcullis.Engine.Tests/Rules/Mutants/}),
nigdy w~produkcyjnym \repopath{RuleRegistry.All}. Każdy mutant jest podstawiany w~tym samym
miejscu rejestracji -- \repopath{Compilation.WithAnalyzers} z~jednym analyzerem -- i~przechodzi
tę samą ścieżkę \repopath{GetAnalyzerDiagnosticsAsync}, co skaner. Mutant używa prawdziwego
\repopath{DiagnosticDescriptor} (ten sam identyfikator i~kategoria), więc jego milczenie to
czysta luka wykrywania. Dla każdego wariantu test najpierw sprawdza, że prawdziwa reguła
łapie naruszenie w~fixture (bez tego „mutant nie złapał” nic by nie znaczyło), potem, że
mutant go nie łapie. Mutanty reguł migracyjnych wołają wewnętrzne pomocniki
\repopath{MigrationSymbols} (\repopath{InternalsVisibleTo}), więc różnią się od reguły jednym
mechanizmem, a~nie reimplementacją. Dwa warianty są odwrócone: dla reguły, której sensem jest
semantyka albo cisza, liczy się błąd „strzela, gdy nie powinna”.

\begingroup\footnotesize
\begin{xltabular}{\linewidth}{@{}>{\raggedright\arraybackslash}p{6.9cm} >{\raggedright\arraybackslash}p{2.6cm} >{\raggedright\arraybackslash}X@{}}
\toprule
\textbf{Mutant} & \textbf{Reguła} & \textbf{Co psuje} \\
\midrule
\endfirsthead
\toprule
\textbf{Mutant} & \textbf{Reguła} & \textbf{Co psuje} \\
\midrule
\endhead
\path{KernelLocCeilingMutant} & P2 sufit & Folder jądra w~liczbie pojedynczej (\texttt{ServiceDefault}) -- żaden projekt \texttt{*.ServiceDefaults} nie jest jądrem. \\
\path{KernelEntityReferenceMutant} & P2 encje & Folder encji \texttt{Domains} zamiast \texttt{Domain}. \\
\path{ControllerNoDbContextMutant} & P9 kontroler & Nazwa typu musi być równa \texttt{DbContext}, a~nie się na nią kończyć -- gubi każdy konkretny kontekst. \\
\path{OrphanEntityMutant} & P9 sierota & Nie wyklucza deklaracji \texttt{DbSet<T>}, więc każda encja „używa się” sama. \\
\path{CustomBaseClassMutant} & P10 & \texttt{TypeKind.Struct} zamiast \texttt{TypeKind.Class} -- warunek nigdy nie jest prawdziwy. \\
\path{EnsureCreatedAsyncVariantDroppedMutant} & P4 & Tylko \texttt{EnsureCreated}, bez wariantu \texttt{EnsureCreatedAsync} używanego w~realnym kodzie. \\
\path{OnModelCreatingTypoMutant} & P4 & Nazwa metody \texttt{OnModelCreated} zamiast \texttt{OnModelCreating}. \\
\path{AdapterSegmentOverbroadMutant} & P11 & Lista adapterów z~\texttt{Application} zamiast \texttt{Integrations} -- zwalnia warstwę, w~której leży realne naruszenie. \\
\path{HostBuilderFactoryNarrowedMutant} & P15 & Tylko \texttt{CreateBuilder}; punkt wejścia workera (\texttt{CreateApplicationBuilder}) nie jest rozpoznawany. \\
\path{ConventionCoverageAggregateMutant} & \texttt{Meta}, odwrócony & Raport per nietrafiona konwencja zamiast „nic nie trafiło”: prawdziwa reguła milczy, mutant strzela. \\
\path{SystemWebNestedNamespacesMutant} & \path{MIG_SYSTEM_WEB} & „Jest System.Web” zamiast „jest System.Web lub zagnieżdżona”; na \texttt{System.Web.SessionState} reguła daje 2~zgłoszenia, mutant 0. \\
\path{SystemWebShortNameMutant} & \path{MIG_SYSTEM_WEB}, odwrócony & Dopasowanie po nazwie zamiast po symbolu; na \texttt{HttpContext} z~ASP.NET Core reguła milczy, mutant strzela. \\
\path{HttpContextCurrentOwnerMutant} & \path{MIG_HTTPCONTEXT_CURRENT} & Szuka \texttt{Current} na \texttt{HttpContextBase}, który go nie deklaruje. \\
\path{SyncOverAsyncGenericAwaitersDroppedMutant} & \path{MIG_SYNC_OVER_ASYNC} & Lista awaiterów bez czterech generycznych -- \texttt{GetAwaiter().GetResult()} na \texttt{Task<T>} przechodzi niezgłoszone. \\
\path{ConfigurationManagerPropertiesOnlyMutant} & \path{MIG_CONFIGURATION_MANAGER} & Tylko referencje właściwości; gubi metodę \texttt{GetSection}. \\
\bottomrule
\end{xltabular}
\endgroup

Wszystkie 15 wariantów jest łapanych (\repopath{MutationTests}, 15 testów). W~przebiegach
ręcznych opisanych dla reguł zasad żaden mutant nie przeżył przy pierwszym uruchomieniu
i~żadna asercja nie wymagała osłabienia po fakcie. Przebieg
na repozytorium referencyjnym dał też ustalenie projektowe: pierwszy skan z~P11 zgłosił
czwarte miejsce ponad przewidziane trzy -- \repopath{Program.cs} z~idiomatyczną rejestracją
klienta SDK w~DI. Zamiast zostawić to jako nieopisane ograniczenie, reguła uznaje plik punktu
wejścia za drugie miejsce adaptera; po poprawce skan dał dokładnie trzy przewidziane trafienia.

\paragraph{Stryker.NET (reguły migracyjne).} Stryker mutuje źródło mechanicznie, więc znajduje
luki, dla których nikt nie napisał mutanta. Jest przypięty jako narzędzie lokalne
(\path{.config/dotnet-tools.json}: \repopath{dotnet-stryker} 5.0.0), a~\path{tests/Portcullis.Engine.Tests/stryker-config.json}
ogranicza go do pięciu plików: \repopath{MigrationSymbols.cs} i~czterech analyzerów
migracyjnych. Poziom mutacji \repopath{Standard}, zbiór testów -- cały
\repopath{Portcullis.Engine.Tests}, progi \repopath{high} 100, \repopath{low} 95,
\repopath{break} 95 (poniżej job CI zawodzi), raporty \repopath{progress},
\repopath{cleartext}, \repopath{html}, \repopath{json}, jedyne wykluczenie to
\repopath{ignore-methods} dla \repopath{EnableConcurrentExecution} (usunięcia tego wywołania
żaden test nie zaobserwuje, więc takie mutanty są z~definicji równoważne).

\begin{lstlisting}
dotnet tool restore
cd tests/Portcullis.Engine.Tests
dotnet stryker            # HTML and JSON reports under StrykerOutput/<timestamp>/reports/
\end{lstlisting}

\begingroup\small
\begin{xltabular}{\linewidth}{@{}>{\raggedright\arraybackslash}X *{6}{>{\raggedleft\arraybackslash}p{1.3cm}}@{}}
\toprule
\textbf{Plik} & \textbf{Zabite} & \textbf{Time\-out} & \textbf{Prze\-żyły} & \textbf{Bez po\-krycia} & \textbf{Błąd kom\-pila\-cji} & \textbf{Pomi\-nięte} \\
\midrule
\endfirsthead
\toprule
\textbf{Plik} & \textbf{Zabite} & \textbf{Time\-out} & \textbf{Prze\-żyły} & \textbf{Bez po\-krycia} & \textbf{Błąd kom\-pila\-cji} & \textbf{Pomi\-nięte} \\
\midrule
\endhead
\texttt{MigrationSymbols.cs} & 12 & 0 & 0 & 0 & 3 & 5 \\
\texttt{SystemWebUsageAnalyzer.cs} & 73 & 2 & 0 & 0 & 1 & 24 \\
\texttt{HttpContextCurrentAnalyzer.cs} & 16 & 0 & 0 & 0 & 0 & 5 \\
\texttt{SyncOverAsyncAnalyzer.cs} & 95 & 0 & 0 & 0 & 5 & 23 \\
\texttt{ConfigurationManagerAnalyzer.cs} & 16 & 0 & 0 & 0 & 0 & 5 \\
\midrule
\textbf{Razem} & \textbf{212} & \textbf{2} & \textbf{0} & \textbf{0} & 9 & 62 \\
\bottomrule
\end{xltabular}
\endgroup

Wynik końcowy: \textbf{100,00\%} -- 214 z~214 testowanych mutantów wykrytych (212 zabitych,
2~timeouty, które Stryker liczy jako wykryte: oba zamieniają pętlę w~nieskończoną). Poza
wynikiem: 9~błędów kompilacji (mutacje, które nie są poprawnym C\#), 4~pominięte przez
\repopath{ignore-methods} i~58 odrzuconych przez filtr Strykera „blok już pokryty”.

Pierwszy przebieg dał 82,80\%, drugi 97,70\% (pięć mutantów, które przeżyły), a~liczby były
mniejszym problemem:
\begin{itemize}[leftmargin=*]
  \item tryb bezpieczny Strykera porzucał wszystkie mutacje w~dwóch metodach reguły
        sync-over-async: zaprzeczenie warunku ze zmienną wzorca (\repopath{x is T t \&\& ...})
        zostawia \repopath{t} nieprzypisane i~psuje kompilację mutanta, a~wynik tego nie
        pokazywał. Reguły używają teraz zwykłych rzutowań i~sprawdzeń \repopath{null};
  \item identyfikatory, tytuły i~komunikaty diagnostyk przeżywały jako mutanty łańcuchów:
        deskryptory są polami statycznymi, inicjalizowanymi raz na proces testów. Są teraz
        stałymi (Stryker nie mutuje stałych), każda reguła ma test przypinający pełny
        komunikat, a~\repopath{RuleRegistryTests} przypina identyfikatory;
  \item część ocalałych mutantów to był martwy kod, którego żaden realny przypadek nie
        rozróżnia -- usunięto go lub przebudowano zamiast podtrzymywać sztucznymi testami;
  \item reszta to prawdziwe luki testów, zamknięte testami w~obu kierunkach (m.in. kod
        generowany, przestrzeń tylko z~tym samym początkiem nazwy -- \repopath{System.WebHooks},
        dyrektywa \repopath{using}, która niczego nie wiąże, aliasy i~\repopath{using static}
        typów spoza System.Web, składowa instancyjna o~nazwie \repopath{Current},
        \repopath{ConfigureAwait} na szybkiej ścieżce i~w~kontynuacji). Helper testów oblewa
        teraz każdy test, w~którym analyzer rzucił wyjątek -- Roslyn zgłasza to jako
        \repopath{AD0001} zamiast oblać test, więc reguła, która się wywróciła, wyglądała
        jak reguła, która milczy.
\end{itemize}

\paragraph{Rola i~granice.} Wynik dowodzi, że każda gałąź logiki decyzyjnej czterech reguł
jest przypięta testem, który zawodzi po jej zmianie, i~że ręczne mutanty centralnych
mechanizmów są łapane. Nie dowodzi braku fałszywych alarmów na prawdziwym zmigrowanym kodzie
-- wszystkie fixture'y są syntetyczne, a~przebieg na kandydacie to ticket P8 benchu
(sekcja~\ref{sec:bench}). Strona silnika w~kanale CLI (\path{src/Portcullis.Engine/Semantics/})
jest ćwiczona end-to-end przez \repopath{ScannerMigrationTests}, ale nie jest w~zakresie
Strykera. Podobnie wcześniej: \path{PORTCULLIS_P9_ORPHAN_ENTITY} łapie trzy z~dziewięciu
wpisów fixture'a sprawdzalnych Roslynem; pozostałe omija przez grubszą definicję
„sieroty” -- to granica zakresu, zapisana jako taka, nie błąd.

\zrodla{\path{letsgolegacy.portcullis/docs/MUTATIONS.md},
\path{letsgolegacy.portcullis/tests/Portcullis.Engine.Tests/Rules/MutationTests.cs},
\path{letsgolegacy.portcullis/tests/Portcullis.Engine.Tests/Rules/Mutants/ConventionCoverageAggregateMutant.cs},
\path{letsgolegacy.portcullis/tests/Portcullis.Engine.Tests/stryker-config.json},
\path{letsgolegacy.portcullis/.config/dotnet-tools.json},
\path{letsgolegacy.portcullis/.github/workflows/mutation.yml},
\path{letsgolegacy.portcullis/src/Portcullis.Rules/MigrationSymbols.cs},
\path{letsgolegacy.portcullis/CONTRIBUTING.md}}

\subsection{Testy i~workflowy CI}\label{sec:portcullis-ci}

\paragraph{Projekty testowe} (xUnit, \repopath{net10.0}):
\begin{itemize}[leftmargin=*]
  \item \repopath{Portcullis.Engine.Tests} -- testy każdej reguły w~obu kierunkach i~jej
        mutanty; test rejestru (11 analyzerów, przypięte identyfikatory reguł migracyjnych,
        format identyfikatorów i~kategorii); skaner (wykluczenia,
        nieistniejąca ścieżka, bramka, pochodzenie, baseline, reguły migracyjne przez ścieżkę
        CLI, C\#~14), pochodzenie (parsery wyjścia git, \repopath{ChangedLines}, dostawca na
        jednorazowych repozytoriach git), fingerprinty i~plik baseline, SARIF z~walidacją
        schematu, \repopath{portcullis.json};
  \item \repopath{Portcullis.CiComment.Tests} -- formatowanie komentarza, diff naruszeń,
        znacznik sticky, publikacja (create/update), opcje; dane z~\path{fixtures/mock-engine-output}
        (ręczne wyniki zgodne z~SPEC §3, z~pierwotnymi identyfikatorami z~łącznikami);
  \item \repopath{Portcullis.Cli.Tests} -- ścisłe parsowanie opcji (flaga bez wartości
        i~literówka w~nazwie flagi to błąd, a~nie po cichu pominięty zakres bramki) oraz
        testy procesu CLI dla SARIF i~baseline.
\end{itemize}
Wymóg PR: \repopath{dotnet build} bez ostrzeżeń („zero, nie kilka”) i~zielone
\repopath{dotnet test}.

\paragraph{Workflowy.} Workflowy na \repopath{pull\_request} nie mają filtra gałęzi: PR są
tu układane w~stos (gałąź każdego ticketu wychodzi z~poprzedniej), a~filtr \repopath{main}
pominąłby CI na wszystkich PR poza pierwszym.

\begingroup\small
\begin{xltabular}{\linewidth}{@{}>{\raggedright\arraybackslash}p{3.1cm} >{\raggedright\arraybackslash}p{3.3cm} >{\raggedright\arraybackslash}X@{}}
\toprule
\textbf{Workflow} & \textbf{Wyzwalacz, uprawnienia} & \textbf{Joby i~co sprawdzają} \\
\midrule
\endfirsthead
\toprule
\textbf{Workflow} & \textbf{Wyzwalacz, uprawnienia} & \textbf{Joby i~co sprawdzają} \\
\midrule
\endhead
\path{ci.yml} (CI) & push do \texttt{main}, każdy PR; \texttt{contents: read}; anulowanie poprzednich przebiegów per PR & \emph{Former name absent} -- \texttt{git grep -n -i} dawnej nazwy projektu (wzorzec zapisany wyrażeniem nawiasowym, żeby workflow nie dopasował sam siebie); porażka, gdy nazwa wróci (ticket R1). \emph{Build and test} -- .NET \texttt{10.0.x}, restore, build i~test w~Release (płytki klon wystarcza, testy pochodzenia tworzą własne repozytoria). \emph{Analyzer package on the .NET 8 SDK} -- sekcja~\ref{sec:portcullis-roslyn}. \\
\path{codeql.yml} (CodeQL) & PR, push do \texttt{main}, harmonogram \texttt{23 3 * * 1}; \texttt{security-events: write} & Analiza CodeQL kodu C\# samego Portcullis (\texttt{build-mode: manual}, build Release) -- „to sprawdza bramkę”. \\
\path{mutation.yml} (Mutation) & PR i~push do \texttt{main}, tylko przy zmianach w~regułach, testach silnika, manifeście narzędzi lub samym workflow & \emph{Stryker (migration rules)}, limit 30 min: \texttt{dotnet tool restore}, \texttt{dotnet stryker}; raport jako artefakt \texttt{mutation-report}; próg \texttt{break} 95. Osobny workflow, bo przebieg trwa ok. sześciu minut na czterech rdzeniach. \\
\path{portcullis-pr-check.yml} (Portcullis PR Check) & PR \texttt{opened}, \texttt{synchronize}, \texttt{reopened}; \texttt{pull-requests: write} & Portcullis bramkuje własne PR: checkout z~\texttt{fetch-depth: 0}, build, przywrócenie poprzedniego wyniku z~cache (klucz per PR i~commit, fallback na ostatni wpis PR), skan \texttt{src} z~zakresem \texttt{base..sha} i~\texttt{--baseline}. Krok skanu zawodzi tylko przy kodzie > 1 lub pustym pliku wyniku (\texttt{dotnet run} zwraca 1 także przy własnej awarii, więc sprawdza się artefakt, nie sam kod). Krok komentarza (zawsze) uruchamia CiComment, z~\texttt{--previous}, jeśli jest poprzedni wynik, i~to on jest bramką (kod 1 przy blokadzie). Wynik trafia do cache na następny push. \\
\path{portcullis-sarif.yml} (Portcullis SARIF) & PR \texttt{opened}, \texttt{synchronize}, \texttt{reopened}; \texttt{security-events: write} & \emph{Filtered SARIF}: skan z~zakresem PR, \texttt{--baseline} i~\texttt{--sarif}; raportuje, nigdy nie blokuje (kod 1 zapisany, awaria to kod 2 lub brak SARIF/JSON). Podsumowanie joba: findingi w~skanie, odrzucone przez każdy filtr, wyniki w~SARIF, zakres, werdykt, notyfikacje. Artefakt \texttt{portcullis-sarif}; \texttt{upload-sarif@v4} z~kategorią \texttt{portcullis} tylko dla PR z~tego samego repozytorium (PR z~forka ma token tylko do odczytu). \\
\path{release.yml} (Release) & push tagu \texttt{v*}; \texttt{workflow\_dispatch} z~wejściem \texttt{version} (buduje, nie publikuje); \texttt{contents: write}, \texttt{packages: write} & \emph{build}: wersja z~tagu, zgodność domyślnego \texttt{version} w~\texttt{action.yml} z~wydaniem, build, test, pack trzech pakietów, konsument testowy, w~którym musi pojawić się \texttt{PORTCULLIS\_P9\_ORPHAN\_ENTITY} (inaczej pakiet jest „martwy”). \emph{publish-nuget} (wymaga \texttt{NUGET\_API\_KEY}, \texttt{--skip-duplicate}), \emph{publish-container} (GHCR, tagi wersji i~\texttt{latest}, notyfikacja o~prywatnym pakiecie), \emph{release-notes} (wydanie GitHub z~artefaktami), \emph{moving-tag} (\texttt{vN} przesuwany dopiero po sukcesie NuGet i~GHCR). \\
\path{secret-scan.yml} (Secret scan) & PR, push do \texttt{main}; \texttt{contents: read} & \emph{gitleaks} 8.28.0 na całej historii (\texttt{fetch-depth: 0}): \texttt{gitleaks detect --redact} z~\texttt{.gitleaks.toml}. \\
\bottomrule
\end{xltabular}
\endgroup

Bramka (\repopath{portcullis-pr-check.yml}) i~SARIF (\repopath{portcullis-sarif.yml}) są
rozdzielone, żeby tylko job SARIF miał \repopath{security-events: write}, a~job z~tokenem
komentarza -- nie. Oba używają tego samego zatwierdzonego \repopath{portcullis-baseline.json}.
Pierwszym PR, na którym przebiegł workflow SARIF, był PR, który go dodał.

\zrodla{\path{letsgolegacy.portcullis/.github/workflows/ci.yml}, \path{letsgolegacy.portcullis/.github/workflows/codeql.yml},
\path{letsgolegacy.portcullis/.github/workflows/mutation.yml},
\path{letsgolegacy.portcullis/.github/workflows/portcullis-pr-check.yml},
\path{letsgolegacy.portcullis/.github/workflows/portcullis-sarif.yml},
\path{letsgolegacy.portcullis/.github/workflows/release.yml},
\path{letsgolegacy.portcullis/.github/workflows/secret-scan.yml},
\path{letsgolegacy.portcullis/docs/SARIF.md}, \path{letsgolegacy.portcullis/CONTRIBUTING.md},
\path{letsgolegacy.portcullis/tests/Portcullis.Cli.Tests/ScanCommandOptionsTests.cs},
\path{letsgolegacy.portcullis/src/Portcullis.Cli/ScanCommandOptions.cs}}

\subsection{Bezpieczeństwo}\label{sec:portcullis-bezpieczenstwo}

\paragraph{Polityka (\texttt{SECURITY.md}).} Wspierana wersja to \repopath{main}; żadne
wydanie nie istnieje, więc nie ma do czego backportować poprawek. Podatność zgłasza się
prywatnie: preferowane jest prywatne zgłaszanie podatności GitHub (Security → Advisories →
Report a~vulnerability), a~gdy jest niedostępne -- publiczne issue z~prośbą o~prywatny kanał,
bez szczegółów. Najbardziej pomagają: wejście, które wyzwala problem (fragment C\#, układ
repozytorium, plik workflow), to, co zyskuje atakujący, i~wersja lub commit. Potwierdzenie
w~ciągu tygodnia; poprawka trafia najpierw na \repopath{main}, a~gdy będą wydania -- wydanie
z~poprawką i~advisory z~podziękowaniem.

Zakres: pakiet analyzerów (kod ładowany przez kompilator w~każdym buildzie i~IDE
konsumenta), CLI i~kontener (uruchamiane na checkoucie, wywołują \repopath{git}), akcja
i~\repopath{Portcullis.CiComment} (token z~\repopath{pull-requests: write}, komentarz budowany
z~wyniku skanu) oraz pipeline wydań. Warte zgłoszenia: spreparowane repozytorium, ścieżka,
nazwa gałęzi lub komunikat commita prowadzące do wstrzyknięcia poleceń, przejścia ścieżek,
ujawnienia tokenu w~logach lub komentarzu albo wykonania kodu w~buildzie konsumenta.
\textbf{Nie} są podatnościami: reguła, która nie łapie naruszenia (to błąd, a~Portcullis jest
bramką zgodności, nie skanerem bezpieczeństwa), fałszywy alarm oraz obejście bramki przez
kogoś, kto może edytować workflow.

\paragraph{Mechanizmy w~kodzie.} Git jest uruchamiany z~listą argumentów, bez powłoki;
powód degradacji jest spłaszczany do jednej linii przed wstawieniem do komentarza Markdown;
porażka publikacji komentarza nie zmienia werdyktu; uprawnienia są rozdzielone między
workflow bramki i~SARIF; upload SARIF jest pomijany dla PR z~forka; kontener działa jako
użytkownik bez uprawnień root. \repopath{CODEOWNERS} obejmuje wszystkie ścieżki i~wymienia
jawnie te istotne dla bezpieczeństwa: \path{/.github/}, \path{/.gitleaks.toml},
\path{/action.yml}, \path{/portcullis-baseline.json}, \path{/src/Portcullis.Rules/}.
\repopath{.gitignore} wyklucza lokalne sekrety (\repopath{.env*}, \repopath{secrets.env})
i~wyniki przebiegów Second Key (\repopath{*.skcap}, \repopath{*.skrun} -- nagrania mogą
zawierać dane klientów) poza próbkami.

\paragraph{Skan sekretów.} Dwie połowy: hook pre-commit łapie błąd, zanim stanie się
historią, a~job CI łapie kontrybutorów bez hooka.
\begin{itemize}[leftmargin=*]
  \item \repopath{.gitleaks.toml}: domyślne reguły gitleaks (\repopath{useDefault = true})
        i~reguła \path{sqlserver-connection-string-with-password} dla connection stringów
        SQL Server z~hasłem w~treści. \repopath{secretGroup = 1} wskazuje samo hasło; bez tego
        gitleaks brałby pierwszą grupę („Server”) i~placeholder nigdy nie trafiałby na
        allowlistę. Allowlista przepuszcza wartości zaczynające się od \repopath{<}, np.
        \repopath{Password=<from SK\_DB\_PASSWORD>} w~fixture'ach i~dokumentacji.
  \item \path{scripts/hooks/pre-commit}: \repopath{gitleaks protect --staged --redact}
        z~konfiguracją repozytorium; bez zainstalowanego gitleaks commit jest odrzucany
        (celowo), a~przy znalezisku hook każe odczytać sekret ze środowiska albo -- dla
        fixture'a, który musi wyglądać jak sekret -- dopisać wyjątek do
        \repopath{.gitleaks.toml} w~tym samym commicie, żeby był widoczny w~review.
  \item \path{scripts/install-hooks.sh} kopiuje hook do \path{.git/hooks/}; bez gitleaks kończy
        się kodem 127, bo instalator zgłaszający sukces, gdy nic nie będzie skanowane, jest
        gorszy niż żaden.
  \item \repopath{secret-scan.yml} skanuje całą historię przy każdym PR i~pushu do
        \repopath{main} -- sekret dodany i~usunięty później też został wypchnięty.
\end{itemize}
Szerzej o~bezpieczeństwie frameworku -- sekcja~\ref{sec:bezpieczenstwo}.

\zrodla{\path{letsgolegacy.portcullis/SECURITY.md}, \path{letsgolegacy.portcullis/.gitleaks.toml},
\path{letsgolegacy.portcullis/scripts/hooks/pre-commit}, \path{letsgolegacy.portcullis/scripts/install-hooks.sh},
\path{letsgolegacy.portcullis/.github/workflows/secret-scan.yml}, \path{letsgolegacy.portcullis/.github/CODEOWNERS},
\path{letsgolegacy.portcullis/.gitignore}, \path{letsgolegacy.portcullis/Dockerfile},
\path{letsgolegacy.portcullis/src/Portcullis.Engine/Provenance/GitCommandRunner.cs},
\path{letsgolegacy.portcullis/src/Portcullis.Engine/Provenance/GitProvenanceProvider.cs},
\path{letsgolegacy.portcullis/src/Portcullis.CiComment/Program.cs}}

\subsection{Zasady wkładu}\label{sec:portcullis-wklad}

Zasada domu: \textbf{twierdzenia to rzeczy, które uruchomiono}. Każde udokumentowane
zachowanie ma dać się prześledzić do polecenia, które ktoś wykonał, i~wyniku, który dało;
to, co niezbudowane, niezweryfikowane albo działające częściowo, dokumentacja nazywa wprost.
Nieprawdziwe twierdzenie w~dokumentacji jest błędem do zgłoszenia. Środowisko: tylko SDK
.NET 10 (polecenia jak w~kroku~7 tutoriala) oraz \path{scripts/install-hooks.sh} dla hooka
skanu sekretów, który wymaga gitleaks.

\paragraph{Dodanie reguły.}
\begin{enumerate}[leftmargin=*]
  \item Zaimplementować w~\path{src/Portcullis.Rules/} w~kształcie istniejącego analyzera;
        identyfikator z~podkreśleniami (Roslyn odrzuca łączniki).
  \item Zarejestrować w~\repopath{RuleRegistry.All} i~zaktualizować \repopath{RuleRegistryTests}.
  \item Testy jednostkowe w~obu kierunkach; test „nie strzela” jest ważniejszy, bo fałszywe
        alarmy wyłączają bramki.
  \item Mutant w~\path{tests/.../Rules/Mutants/} i~przypadek w~\repopath{MutationTests}:
        prawdziwa reguła łapie, mutant nie. Psuć konkretny mechanizm, nie coś
        przypadkowego. Dla reguł migracyjnych także Stryker.NET i~dopisanie pliku do listy
        \repopath{mutate} w~\repopath{stryker-config.json}.
  \item Uruchomić na prawdziwym kodzie, nie tylko fixture'ach, i~zapisać wynik -- także to,
        co reguła zgłosiła niesłusznie.
  \item Udokumentować w~\repopath{docs/MUTATIONS.md} i~tabeli README (reguła migracyjna -- także
        w~\repopath{docs/rules/MIGRATION.md}).
\end{enumerate}
Lepiej mała reguła o~wysokiej pewności niż szeroka, która wymaga strojenia.

\paragraph{Pull request.} Portcullis bramkuje własne PR względem zatwierdzonego baseline:
nowy błąd na zmienionej linii blokuje; należy go naprawić, a~nie powiększać baseline. Jeśli
akceptacja jest właściwa, baseline regeneruje się poniższym poleceniem i~uzasadnia w~PR:

\begin{lstlisting}
dotnet run --project src/Portcullis.Cli -- scan src --baseline portcullis-baseline.json --write-baseline
\end{lstlisting}

Checklista szablonu PR:
\begin{checklista}
  \item \repopath{dotnet build portcullis.sln} czysty -- zero ostrzeżeń, nie „kilka”.
  \item \repopath{dotnet test portcullis.sln} zielony.
  \item Nowa reguła ma test jednostkowy w~obu kierunkach \textbf{i}~mutanta
        (\repopath{CONTRIBUTING.md}, „Adding a~rule”).
  \item Dokumentacja zaktualizowana tam, gdzie zmiana dotyka udokumentowanego zachowania.
  \item Żadne twierdzenie w~PR ani w~dokumentach, których dotyka, nie jest mocniejsze niż
        to, co zweryfikowano.
\end{checklista}
Szablon prosi też o~sekcje: co się zmienia, co naprawdę uruchomiono (polecenia i~wynik),
dowód działania (dla reguły: przypadek strzelający, niestrzelający i~mutant; dla bramki lub
CI: werdykt przed i~po na prawdziwym diffie) oraz znane ograniczenia. Szablony issue
rozróżniają fałszywy alarm, inny błąd i~propozycję reguły; podatność idzie prywatnie przez
\repopath{SECURITY.md}.

\zrodla{\path{letsgolegacy.portcullis/CONTRIBUTING.md}, \path{letsgolegacy.portcullis/.github/pull_request_template.md},
\path{letsgolegacy.portcullis/.github/ISSUE_TEMPLATE/config.yml},
\path{letsgolegacy.portcullis/.github/ISSUE_TEMPLATE/false_positive.yml},
\path{letsgolegacy.portcullis/.github/ISSUE_TEMPLATE/bug_report.yml},
\path{letsgolegacy.portcullis/.github/ISSUE_TEMPLATE/rule_proposal.yml}}

\subsection{Historia zmian (CHANGELOG)}\label{sec:portcullis-changelog}

CHANGELOG ma format Keep a~Changelog 1.1.0, wersje -- Semantic Versioning 2.0.0.
\textbf{Nic nie zostało wydane}, więc wszystkie wpisy są pod \repopath{[Unreleased]},
a~\repopath{0.1.0} to numer przyszłego pierwszego wydania.

\begin{itemize}[leftmargin=*]
  \item \textbf{Added}: silnik reguł na modelu semantycznym Roslyn; dziewięć diagnostyk
        sześciu zasad (P2, P4, P9, P10, P11, P15) z~testami w~obu kierunkach i~mutantami;
        dwie diagnostyki pokrycia konwencji; konfigurowalne konwencje
        (\repopath{portcullis.json} i~\repopath{.globalconfig}); bramka zawężona do diffu;
        pochodzenie AI jako wejście progu poziomu; cztery kanały dystrybucji (zbudowane
        i~zweryfikowane lokalnie, nieopublikowane); komentarz sticky; pipeline wydań
        sterowany tagami z~bramką zgodności \repopath{action.yml}; dokumentacja dla czytelników
        z~zewnątrz; zestaw reguł migracyjnych (R2, Stryker 100\%, 214 z~214); SARIF 2.1.0
        z~filtrem zmienionych linii i~plikiem baseline (R3).
  \item \textbf{Changed}: import do \repopath{letsgolegacy.portcullis} jako komponent C5 (R0),
        bez historii repozytorium-poprzednika i~bez wewnętrznych notatek planistycznych;
        projekty \repopath{net8.0} przeszły na \repopath{net10.0}, analyzery zostały na
        \repopath{netstandard2.0}; nazwa Portcullis ustalona 2026-09-13, przed jakąkolwiek
        publikacją (poprzednia kolidowała z~istniejącym produktem w~tej samej dziedzinie),
        i~zmiana nazwy dokończona z~kontrolą CI (R1); skan CLI kompiluje się względem
        zestawów frameworku i~deklaracji API .NET Framework; wynik skanu niesie fingerprinty
        i~baseline (addytywnie); jedna klasa sprawdza zakres diffu dla bramki, eskalacji
        i~SARIF.
  \item \textbf{Fixed}: bramka nie przechodzi już przy awarii git (fail-open,
        \repopath{ProvenanceStatus}); CLI respektuje konfigurację (wcześniej
        \repopath{options: null}); ścieżka komentarza PR nie wywraca się na zniekształconym
        wyniku; dwa findingi w~jednej linii mają stałą kolejność między przebiegami; dwa zdania
        pozostałe po zmianie nazwy.
  \item \textbf{Known limitations}: analiza pojedynczej kompilacji w~CLI (pakiet analyzerów
        tego ograniczenia nie ma); dziewięć z~piętnastu zasad niewdrożonych; projekt
        nieopublikowany.
\end{itemize}

\zrodla{\path{letsgolegacy.portcullis/CHANGELOG.md}}

\subsection{Ograniczenia}\label{sec:portcullis-ograniczenia}

\begin{itemize}[leftmargin=*]
  \item \textbf{Analiza jednej kompilacji w~CLI.} Źródła są parsowane bezpośrednio, nie
        przez MSBuild, więc typ z~pakietu NuGet lub innego projektu się nie rozwiązuje. Skan
        referencjonuje zestawy frameworku i~deklaracje API .NET Framework, więc reguły
        migracyjne wiążą te same symbole co w~prawdziwym buildzie; ich skutki tego
        ograniczenia opisano per reguła (sekcja~\ref{sec:portcullis-migracja}).
  \item \textbf{Konwencje są konfigurowalne, kształt reguły nie.} Jądro, domena i~adapter
        są rozpoznawane po nazwach folderów i~plików, nie po niczym zadeklarowanym w~kodzie.
  \item \textbf{Konwencja, która nic nie dopasowała, jest raportowana} -- ostrzeżeniami,
        które same nie blokują.
  \item \textbf{Dziewięć z~piętnastu zasad jest niewdrożonych}: pięć wymaga kontekstu
        międzyplikowego, między repozytoriami albo z~historii git, cztery dotyczą innych
        artefaktów (Dockerfile, \repopath{fly.toml}, YAML CI, proza). P5 jest rozdzielona:
        skan literałów sekretów dałoby się zbudować dziś, „dokładnie jedna usługa trzyma
        klucz podpisu” to niezmiennik między repozytoriami.
  \item \textbf{Poziom da się stroić tylko w~kanale pakietu analyzerów}; akcja nie udostępnia
        \repopath{--sarif} ani \repopath{--baseline}; nic nie jest opublikowane, a~obraz GHCR
        i~workflow wydania nie były wykonane.
\end{itemize}

\zrodla{\path{letsgolegacy.portcullis/README.md}, \path{letsgolegacy.portcullis/CHANGELOG.md},
\path{letsgolegacy.portcullis/docs/CONFIGURATION.md}, \path{letsgolegacy.portcullis/docs/SARIF.md},
\path{letsgolegacy.portcullis/docs/SPEC.md}}

\subsection{Status według WORKPLAN}\label{sec:portcullis-status}

Portcullis kontynuuje prywatne repozytorium-poprzednika; w~tym repozytorium zmiana nazwy
staje się kompletna i~dochodzi zestaw reguł migracyjnych. Identyfikatory ticketów są wspólne
z~backlogiem całego frameworku; jeden ticket to jeden pull request.

\begin{xltabular}{\linewidth}{@{}l >{\raggedright\arraybackslash}X >{\raggedright\arraybackslash}p{4.4cm} l@{}}
\toprule
\textbf{ID} & \textbf{Zakres} & \textbf{Kryterium ukończenia} & \textbf{Status} \\
\midrule
\endfirsthead
\toprule
\textbf{ID} & \textbf{Zakres} & \textbf{Kryterium ukończenia} & \textbf{Status} \\
\midrule
\endhead
R0 & Import kodu (źródła, testy, fixture'y, dokumentacja użytkownika) z~zielonym CI & Build i~pełny zestaw testów przechodzą w~CI tego repozytorium & zrobione (\#1) \\
R1 & Dokończenie zmiany nazwy: repozytorium, pakiety, przestrzenie nazw & Stara nazwa nie występuje w~drzewie, egzekwowane przez CI & zrobione (\#2) \\
R2 & Reguły migracyjne: System.Web, \texttt{HttpContext.Current}, \texttt{.Result}/\texttt{.Wait()}, \texttt{ConfigurationManager} zamiast \texttt{IOptions} & Cztery analyzery z~testami sprawdzonymi mutacyjnie & zrobione (\#3) \\
R3 & SARIF 2.1 z~filtrem zmienionych linii i~plikiem baseline & PR produkuje przefiltrowany SARIF & zrobione (\#4) \\
P8 support & Reguły na PR zmigrowanego kandydata w~benchu nopCommerce & SARIF dołączony do pakietu dowodowego (śledzone w~benchu) & planowane \\
\bottomrule
\end{xltabular}

Dalsze fazy: reguły warstw w~stylu ArchUnitNET i~kolejne reguły (faza 02), pakiety reguł
per estate (faza 03), analyzery Java dla modernize-java (faza 04); sposób dystrybucji
(publiczna albo licencjonowana) pozostaje decyzją otwartą. Status całego frameworku --
sekcja~\ref{sec:status}.

\zrodla{\path{letsgolegacy.portcullis/docs/WORKPLAN.md}, \path{letsgolegacy.portcullis/README.md}}
